% Part: core
\maketitle
\noteprovenance

\begin{abstract}
A degree-seven counterexample due to the erdosproblems.com contributor
\href{https://www.erdosproblems.com/forum/thread/1041\#post-8861}{\texttt{ani}},
formalised in Lean here, refutes the universal total-variation
formulation: every continuous path inside the strict unit lemniscate joining
two distinct roots has total variation greater than $2$.  Its correspondence
with the historical curve-length question has not received independent human
review.
For a monic trinomial with roots in the open unit disc, the root equation
bounds the polynomial along every root-to-origin segment. Any two distinct roots
are therefore joined inside $\{|f|<1\}$ by a path of length less than $2$.
For any squarefree monic polynomial $f$ of degree $n\ge2$, an area-growth
argument gives such a path when its least critical-value modulus $\mu$
is at most $13/25$, without a root-location hypothesis. Scaling gives
length less than $(5/2)\mu^{1/n}$ inside $\{|f|<(25/13)\mu\}$.
We develop this argument, an independent inverse-ray averaging estimate,
and a construction using an isolated simple critical value. We also prove
results for collinear roots and several coefficient families.
A weighted Poisson identity gives the mean bound, with sharp constant,
$\sum_{j=1}^{n-1}|f(c_j)|^{4/(n-1)}\le(n-1)R^{4n/(n-1)}$
for roots in a closed disc of radius $R$, with critical points counted
with multiplicity. Examples distinguish this bound from path-length control.
The final sections examine inverse-sheet topology and refute an earlier
spanning-tree estimate. The positive results retain their stated hypotheses
and do not repair the refuted universal formulation.
\end{abstract}

The \href{erdos-1041-lemniscate-newton-flow.pdf}{short note} centres on
the trinomial identity in Section~\ref{sec:trinomials}. That proof does
not use Newton flow. The area-growth comparison is proved here in
Section~\ref{sec:low-critical-closure}. The other constructions are
independent of that proof. This record provides their derivations and
keeps failed estimates beside the witnesses that refute them.

\begin{center}
\begin{tabular}{@{}p{0.27\linewidth}p{0.65\linewidth}@{}}
\toprule
Short-note passage & Argument in this record\\
\midrule
\raggedright Section 3: small critical value &
Section~\ref{sec:low-critical-closure}: inverse-map lengths,
packing, area growth and the rational certificate.\\[3pt]
\raggedright Supplement after Section 3: area estimate &
Section~\ref{sec:constant-factor}: regular-level selection,
low and high lifts, and the adjacent-pair average.\\[3pt]
\raggedright Section 6: isolated critical value &
Section~\ref{sec:critical-value-separation}: the two-sheeted component,
square-root map, Bergman segment bound and capacity estimate.\\[3pt]
\raggedright Sections 13--14: chords and collinear roots &
Section~\ref{sec:frontier}: the \hyperlink{binomial-chord-calculation}{exact binomial chord calculation};
Section~\ref{subsec:sharp-collinear-chebyshev}: Chebyshev comparison
and \hyperlink{collinear-source-comparison}{related extremal problems}.\\[3pt]
\raggedright Section 15: critical-value means &
Section~\ref{sec:frontier}: the reflected derivative,
weighted Poisson identity, boundary passage and equality cases.\\[3pt]
\raggedright Section 17: limits and examples &
Section~\ref{sec:frontier}, subsections
\hyperlink{closed-class-limits}{Limits of contained paths at fixed degree}
and \hyperlink{blaschke-product-examples}{Blaschke-product examples and limits in varying degree}.\\
\bottomrule
\end{tabular}
\end{center}

\section{The historical question}\label{sec:problem}

\begin{problem}[Erd\H{o}s \#1041]\label{res:problem}
Let $f(z)=\prod_{i=1}^{n}(z-z_i)$ be monic of degree $n\ge2$, with
$z_i\in\dbar$, the open unit disc.  Must two of the roots be joined by a curve of length less than
$2$ lying in the open lemniscate $\{z\in\mathbb{C}:|f(z)|<1\}$?
\end{problem}

The open-disc hypothesis is essential for this formulation.  For
$f(z)=z^2-1$, whose roots lie on the unit circle, every continuous path from
$-1$ to $1$ meets the imaginary axis.  There $|f(iy)|=1+y^2\ge1$, so no
such path lies in the open unit lemniscate.  Repeated roots, counted as distinct
occurrences, give a constant path; the substantive case is therefore
squarefree.

\noindent The problem numbering follows Bloom's Erd\H{o}s problem catalogue
\cite{bloom}.  \phantomsection\label{res:ani-degree-seven-counterexample-long}
A degree-seven counterexample due to the
\href{https://www.erdosproblems.com/forum/thread/1041\#post-8861}{erdosproblems.com contributor \texttt{ani}} makes the unrestricted total-variation assertion false.  For
one explicit monic polynomial of degree seven whose roots are distinct and
lie in the open unit disc, every continuous path in the strict unit lemniscate
joining two distinct roots has total variation greater than~$2$.\footnote{Lean
checks this instance in
\href{https://github.com/wcook04/plectis-erdos/blob/f70679d47178bf1174309cb29f5b31748609cc04/lean/ErdosProblems/Erdos1041/Counterexample/CatalogueAdapter.lean\#L25}{\texttt{erdos1041\_ani\_degree\_seven}}.
The formal result concerns one polynomial; it does not formalise the reported
small-parameter family or a Hausdorff-measure comparison, and its correspondence
with the historical formulation has not received independent human review.}
Thus this polynomial has no continuous root-to-root parametrisation of total variation less than~$2$.
The original source is Problem~5 on printed p.~139 of
Erd\H{o}s--Herzog--Piranian \cite[Problem~5, p.~139]{ehp1958}; the preceding paragraph records the
known input that one component of the lemniscate contains at least two zeros.

The main constructions use different kinds of information. The
trinomial identity controls a prescribed path directly. The area argument
instead uses the number of roots in a sublevel component. Source comparisons for
component counts and lemniscate topology accompany the inverse-sheet
discussion in Section~\ref{sec:gap}; the Newton-flow terminology is
introduced in Section~\ref{sec:newton}.

Two recent manuscripts are relevant.  The 48-page manuscript posted by
\texttt{shtuka} on 24 March 2026
\cite[Theorem~1, p.~1]{march2026} claims the unrestricted
statement.  Its Proposition~12 (p.~16, with proof continuing through p.~30)
supplies the spanning-tree decomposition used in the final proof.  The defect
was located publicly in the problem's discussion thread: on 25 March 2026 Tao
observed that the invocation of Lemma~8 there is unjustified and that the flow
lines need not organise into connected trees, and on 26 March 2026 the
manuscript's author agreed that the statement of Proposition~12 itself, not
only its printed proof, is incorrect, and set the strategy aside.  Section~\ref{sec:gap}
records an independent diagnosis of the same local saddle defect, together
with the Cassini obstruction to the printed global tree-length bound.  That
obstruction refutes the metric estimate used in the final proof, not Erdős \#1041.  Pendyala's independent June 2026
preprint \cite[Thm.~1, p.~1]{june2026} proves the degree-four case through a
finite four-point radial lemma and a short polygonal connector.  That is the
degree-four result directly comparable to the root-pair problem. Together
with the cubic theorem established here, it settles these two degrees; it
does not supply the general-degree conclusion. The all-degree estimates below either impose a bound on the least
critical-value modulus or allow a larger length constant and containment
level. They address different hypotheses from the quartic theorem.

\paragraph{What the formal links establish.}
An ordinary mathematical argument is distinct from a proof checked by Lean.
A linked declaration may prove a finite inequality without constructing the
curve used here. Moreover, a declaration may be present in the versioned
sources but absent from their recorded checked build. We distinguish these
cases where the declarations are used. Neither a successful build nor an
exact finite calculation establishes an unformalised analytic step.

\section{Trinomials}
\label{sec:trinomials}

Write $E_f=\{z\in\mathbb C:|f(z)|<1\}$.  The first theorem gives a prescribed
path between every pair of zeros of a trinomial, in every degree.

\begin{theorem}[trinomial root connections]
\label{res:trinomial-all-degree}
Let $n,m$ be integers with $1\le m<n$, and let $f(z)=z^n+az^m+b$ have every
zero in $\dbar$.  For every zero $\zeta$, the segment $[0,\zeta]$ lies in
$E_f$.  Consequently any two zeros $\zeta_1,\zeta_2$ are joined in $E_f$ by
the broken line $\zeta_1\to0\to\zeta_2$, of length $|\zeta_1|+|\zeta_2|<2$.
\end{theorem}
\leannote{Lean: \lproof{ErdosProblems/Erdos1041/PaperTrinomial.lean}{26}{all_spokes}, \lproof{ErdosProblems/Erdos1041/PaperTrinomial.lean}{38}{complete_trinomial}.}

\begin{proof}
Vieta's formula gives $|b|<1$.  At a zero $\zeta$ the root equation
$\zeta^n+a\zeta^m+b=0$ eliminates the middle coefficient:
\begin{equation}\label{eq:trinomial-cancellation}
 f(t\zeta)=b(1-t^m)+\zeta^n(t^n-t^m).
\end{equation}
For $0\le t<1$ the two scalar weights $1-t^m$ and $t^m-t^n$ are nonnegative
and sum to $1-t^n$, so
\[
 |f(t\zeta)|\le|b|(1-t^m)+|\zeta|^n(t^m-t^n)<1-t^n\le1 .
\]
At $t=1$ the value is zero.  Concatenating two such segments gives the length
assertion.
\end{proof}

The coefficient $a$ carries no hypothesis; the root equation removes it before
absolute values are taken.  The conclusion concerns segments to zeros and
makes no assertion that $E_f$ is star-shaped.

\paragraph{The cancellation mechanism.}
\label{subsec:abel-mechanism}
For $f(z)=\sum_{k=0}^nc_kz^k$ and a zero $\zeta$, put
$S_j=\sum_{k=0}^jc_k\zeta^k$.  Finite summation by parts gives
\begin{equation}\label{eq:abel-control-polygon}
 f(t\zeta)=\sum_{j=0}^{n-1}(t^j-t^{j+1})S_j\qquad(0\le t\le1),
\end{equation}
since the coefficient of $c_k\zeta^k$ on the right is $t^k-t^n$ and
$\sum_{k<n}c_k\zeta^k=-c_n\zeta^n$.  The weights
in~\eqref{eq:abel-control-polygon} are nonnegative and sum to $1-t^n$, so the
whole radial segment lies in the closed unit sublevel set whenever every partial sum lies in the closed
unit disc. For a trinomial, $S_j=b$ for $j<m$ and $S_j=-\zeta^n$
for $m\le j<n$; these two values may coincide. This is exactly the
estimate above. With two intermediate coefficients, the root-disc
hypothesis need not put every partial sum in the unit disc. The next
example shows that a prescribed radial segment can then escape.

\begin{example}[an escaping root spoke]\label{ex:sextic-spoke}
Let $0<r<1$ satisfy $r^6>320/327$, and set
\[
 f_r(z)=z^6+\tfrac15r^2z^4-\tfrac15r^4z^2-r^6
       =(z^2-r^2)\bigl(z^4+\tfrac65r^2z^2+r^4\bigr).
\]
Every zero has modulus $r$: for the quartic factor put $z=rw$, so that the
squared roots solve $v^2+(6/5)v+1=0$, whose two roots are complex conjugates
of product one.  Nevertheless
\[
 f_r(r/2)=-\frac{327}{320}\,r^6,
\]
so the segment $[0,r]$ leaves $E_{f_r}$.  This rules out a universal assertion
about every origin-to-zero segment.  It exhibits no counterexample to the
existence of some short connection.
\end{example}

\paragraph{Sources and scope.}
\label{bdry:trinomial}
The identity~\eqref{eq:abel-control-polygon} is checked as
\href{https://github.com/wcook04/plectis-erdos/blob/0ba585f632fbbbb43af2ee9532d4e83752af3f67/ErdosProblems/Erdos1041/AbelControlPolygon.lean\#L123}{the Abel summation identity}, the
constant-term bound as
\href{https://github.com/wcook04/plectis-erdos/blob/0ba585f632fbbbb43af2ee9532d4e83752af3f67/ErdosProblems/Erdos1041/AbelControlPolygon.lean\#L258}{the constant-term estimate from the root-disc hypothesis},
the radial estimate as
\href{https://github.com/wcook04/plectis-erdos/blob/0ba585f632fbbbb43af2ee9532d4e83752af3f67/ErdosProblems/Erdos1041/AbelControlPolygon.lean\#L219}{containment along a trinomial root segment},
and their combination with the length bound as
\href{https://github.com/wcook04/plectis-erdos/blob/0ba585f632fbbbb43af2ee9532d4e83752af3f67/ErdosProblems/Erdos1041/AbelControlPolygon.lean\#L330}{the two segment inequalities and their length bound}.
Example~\ref{ex:sextic-spoke} is checked as
\href{https://github.com/wcook04/plectis-erdos/blob/0ba585f632fbbbb43af2ee9532d4e83752af3f67/ErdosProblems/Erdos1041/AbelControlPolygon.lean\#L555}{the sextic evaluation}.  What the
kernel proves is the pair of radial inequalities and the bound
$|\zeta_1|+|\zeta_2|<2$ on the sum of the two radii.  Assembling those into a
single rectifiable path object, and the passage to the one-dimensional
Hausdorff measure used by the original formulation of the problem, are ordinary
steps taken here.  No novelty claim is made for this elementary trinomial argument.

\section{A small least critical value}
\label{sec:low-critical-closure}

The next criterion restricts the smallest critical-value modulus, but
places no condition on the coefficients or the root locations. Throughout
this section
\[
 \mu=\min_{f'(c)=0}|f(c)|
\]
is the least critical-value modulus.

\begin{theorem}[a small least critical value forces a short connector]
\label{res:low-critical-thirteen-twentyfifths}
Let $f$ be squarefree and monic of degree $n\ge2$ with $\mu\le13/25$.  Then
two distinct roots of $f$ are joined inside $\{|f|<1\}$ by a rectifiable
curve of length strictly below $2$.  No hypothesis is placed on the locations
of the roots, on the number of roots in any component, or on the capacity of
any component.
\end{theorem}
\leannote{Not formalised: the degree-two case and the exact rational arithmetic of the stopping-time comparison are checked, and the analytic argument in higher degrees is not. See the \hyperref[sec:coverage]{coverage section}.}

For $f(z)=z^n-b$ with $0<|b|<1$, the only critical point is $0$ and
$\mu=|b|$. Thus this criterion includes $|b|\le13/25$ but excludes
$13/25<|b|<1$, although the trinomial theorem gives the required path
in both ranges. Its value is that it also applies to polynomials with
arbitrary coefficient patterns, including roots outside the unit disc. For example,
$(z-3)^n-1/2$, $n\ge2$, has $\mu=1/2$ and all its roots satisfy
$|z|>2$. The required containment is $|f(z)|<1$, not $|z|<1$.

\begin{corollary}[scale-free form]
\label{res:low-critical-scale-free}
Every squarefree monic $f$ of degree $n\ge2$ has two distinct roots joined
inside $\{|f|<(25/13)\mu\}$ by a curve of length below
$2\bigl((25/13)\mu\bigr)^{1/n}$.
\end{corollary}
\leannote{Lean: \lproof{ErdosProblems/Erdos1041/PaperCompleteR21/LowCriticalScaleTransport.lean}{460}{scaledLowCritical_of_lowCritical}. Conditional on Theorem~\ref{res:low-critical-thirteen-twentyfifths} as stated; see the \hyperref[sec:coverage]{coverage section}.}

\begin{proof}
Apply Theorem~\ref{res:low-critical-thirteen-twentyfifths} to
$s^{-n}f(sz)$ with $s=\bigl((25/13)\mu\bigr)^{1/n}$, whose least
critical-value modulus is $13/25$, and scale back. Squarefreeness gives
$\mu>0$, so the scaling factor is positive.
\end{proof}

The proof has two steps. Failure of the length bound first forces the
roots to be far apart in the hyperbolic metric of a sublevel component.
Packing those roots gives a lower bound for their number. Averaging paths
between adjacent boundary points then converts that root count into an
area-growth inequality. The rational certificate shows that this growth
would exceed P\'olya's area bound before the value level reaches $1$.

\paragraph{Consequences of having no short path.}
Assume no two distinct roots are joined inside $\{|f|<1\}$ by a curve of
length below $2$.  Choose a critical point at level $\mu$ and two descending
inverse arcs into distinct one-root components of $\{|f|<\mu\}$; below $\mu$
each component maps conformally onto the value disc, and distinct local
inverse arcs at the first critical point enter distinct components, since two
inverse images of a nearby regular value in one component would contradict its
degree one.  Their union is a compact connected set containing two roots.
Fix a level $t\in(\mu,1)$ at which no critical value has modulus $t$,
so that the level is regular. Let $C_t$ be the component of
$\{|f|<t\}$ containing that set, let $k\ge2$ be its root count, and put
$x=\log(t/\mu)$ and
$a=\operatorname{Area}(C_t)/\pi$; P\'olya's inequality \cite[Theorem~1]{crane}
gives $a\le t^{2/n}<1$.
The component is a Jordan domain. List its roots as
$\zeta_1,\ldots,\zeta_k$, choose a Riemann map
$\varphi:\mathbb D\to C_t$, and write $b_j=\varphi^{-1}(\zeta_j)$.
The component-wise form of \cite[(2.2)]{eks2010} is
\[
 \frac{f(\varphi(z))}{t}
 =e^{i\theta}\prod_{j=1}^k\frac{z-b_j}{1-\overline{b_j}z}.
\]
Indeed, regularity extends $\varphi$ across the Jordan boundary. The
quotient of the two sides without the phase has no zero or pole on the
closed disc and has boundary modulus one. The maximum principle,
applied to that quotient and its reciprocal, makes it a unimodular
constant. The product has degree $k$, the root count of this component,
not necessarily the full degree $n$. Distances between the $b_j$ below are
hyperbolic distances in $\mathbb D$, normalised by
$d(0,s)=2\operatorname{artanh}s$ for $0\le s<1$. Intrinsic distances
in $C_t$ are infima of Euclidean lengths of curves in that component.

The length estimate converts area into a bound for the image of a compact real
interval $I\subset(-1,1)$. Here is the coefficient calculation underlying
the Bergman-kernel argument. Write
$\varphi'(z)=\sum_{\ell\ge0}\alpha_\ell z^\ell$. Injectivity and the
area formula give
$\sum_{\ell\ge0}|\alpha_\ell|^2/(\ell+1)=a$.
To express length as a linear integral, choose
$\omega(x)=\overline{\varphi'(x)}/|\varphi'(x)|$ on $I$; the denominator
is nonzero because $\varphi$ is conformal. Cauchy--Schwarz gives
\[
 \operatorname{length}(\varphi(I))^2
 =\left|\sum_{\ell\ge0}\alpha_\ell\int_I\omega(x)x^\ell\,dx\right|^2
 \le a\sum_{\ell\ge0}(\ell+1)
       \left|\int_I\omega(x)x^\ell\,dx\right|^2.
\]
Compactness of $I$ permits termwise integration of the series.
Since $\sum_{\ell\ge0}(\ell+1)(xy)^\ell=(1-xy)^{-2}>0$ for real
$x,y\in I$, taking absolute values in the resulting double integral yields
\[
 \operatorname{length}(\varphi(I))^2
 \le a\int_I\!\int_I\frac{dx\,dy}{(1-xy)^2}.
\]
This is precisely the positivity of the disc Bergman kernel
$K(z,w)=\pi^{-1}(1-z\overline w)^{-2}$ on the real diameter.
For $I=[u,v]\subset(-1,1)$, direct integration gives
\[
 \int_u^v\!\int_u^v\frac{dx\,dy}{(1-xy)^2}
 =\log\frac{(1-uv)^2}{(1-u^2)(1-v^2)}.
\]
Thus $I=[0,r]$ gives
$\operatorname{length}(\varphi([0,r]))^2\le a\log(1/(1-r^2))$.
Moving two preimages to $\pm s$ and taking $I=[-s,s]$ instead gives
\[
 \operatorname{length}(\varphi([-s,s]))^2\le4a\operatorname{artanh}(s^2).
\]
A pair of roots at hyperbolic distance $d$ may be placed at $\pm\tanh(d/4)$,
so a connector shorter than $2$ exists as soon as
$\operatorname{artanh}(\tanh^2(d/4))<1/a$.  Failure therefore gives
\begin{equation}\label{eq:lc-separation}
 d(b_i,b_j)\ \ge\ D:=4\operatorname{artanh}\sqrt{\tanh(1/a)},
 \qquad \cosh(D/2)=e^{2/a}.
\end{equation}

There is also a point $h$ in the chosen compact connected set whose
intrinsic distance in $C_t$ from every root is at least $1$.
Otherwise the open intrinsic balls of radius $1$ about the roots would
cover that set. Two balls with distinct centres cannot meet: paths from
their centres to a common point would concatenate to length below $2$.
But a connected set containing two distinct roots cannot be covered by
these disjoint open sets. This contradiction proves the assertion.
Now choose the Riemann map with $\varphi(0)=h$, and continue to write
$b_j=\varphi^{-1}(\zeta_j)$. Put $d_j=d(0,b_j)$ and
$\lambda(d)=-\log\tanh(d/2)$.  The one-root estimate at $h$
gives $\tanh(d_j/2)\ge\sqrt{1-e^{-1/a}}$. The centre $h$ lies on the
chosen inverse arcs, so $0<|f(h)|\le\mu$; it is not a root because its
intrinsic distance from every root is at least $1$. The Blaschke identity
at $0$ gives $|f(h)|/t=\prod_j|b_j|$, hence
$\sum_j\lambda(d_j)=\log(t/|f(h)|)\ge\log(t/\mu)=x$.
Together these estimates give
\begin{equation}\label{eq:lc-radius-and-budget}
 \lambda(d_j)\le\frac{\delta(a)}2,\quad
 \delta(a)=-\log\bigl(1-e^{-1/a}\bigr),
 \qquad
 \sum_{j=1}^{k}\lambda(d_j)\ \ge\ x .
\end{equation}

\paragraph{A lower bound for the number of roots.}
Testing only consecutive pairs in angular order loses essential information.
For $k=2m$, put $m$ points at hyperbolic radius $d_{\mathrm{low}}$ and
$m$ at radius $d_{\mathrm{low}}+D$, alternating in cyclic order. The
reverse triangle inequality separates every consecutive pair by at least
$D$, regardless of the angular gaps, whereas
\[
 \sum_j\lambda(d_j)
 =m\bigl(\lambda(d_{\mathrm{low}})+\lambda(d_{\mathrm{low}}+D)\bigr)
 \longrightarrow\infty\qquad(m\to\infty).
\]
The even number of points is needed for cyclic alternation. These point
configurations test only the consecutive-pair constraints; they need not
satisfy separation for nonconsecutive pairs or arise from a polynomial
component. To use the missing pairwise information, we impose the packing
bound on every circle centred at the origin.

\begin{lemma}[circle-slice packing]\label{res:circle-slice-packing}
Under~\eqref{eq:lc-separation}, for every $r>0$,
\[
 \sum_{j=1}^{k}w(d_j,r)\le\pi,\qquad
 w(d,r)=\arccos\Bigl(\operatorname{clamp}
   \frac{\cosh d\cosh r-\cosh(D/2)}{\sinh d\sinh r}\Bigr),
\]
where $\operatorname{clamp}$ truncates its argument to $[-1,1]$.
\end{lemma}
\leannote{Lean: \lproof{ErdosProblems/Erdos1041/PaperCompleteR21/HyperbolicLawOfCosines.lean}{563}{circle_slice_packing}, \lproof{ErdosProblems/Erdos1041/PaperCompleteR21/HyperbolicLawOfCosines.lean}{141}{cosh_dist_polar}, \lproof{ErdosProblems/Erdos1041/PaperCompleteR21/HyperbolicLawOfCosines.lean}{166}{dist_polar_I}, \lproof{ErdosProblems/Erdos1041/PaperCompleteR21/HyperbolicLawOfCosines.lean}{177}{exists_polar}, and 2 further declarations in the \hyperref[sec:coverage]{coverage section}.}

\begin{proof}
The open balls $B_j=B_{\mathrm{hyp}}(b_j,D/2)$ are pairwise disjoint, since a
common point would force $d(b_i,b_j)<D$.  Fix $r>0$.  By the hyperbolic law of
cosines the point of the hyperbolic circle of radius $r$ at angle $\theta$
lies in $B_j$ exactly when
$\cosh d_j\cosh r-\sinh d_j\sinh r\cos(\theta-\theta_j)<\cosh(D/2)$.
The set of angles satisfying this inequality has measure $2w(d_j,r)$:
clamping gives measure zero for an empty intersection and $2\pi$ for a
full circle, with tangent endpoint sets of measure zero. The intersections
of the balls with this circle are disjoint, so their angular measures
sum to at most $2\pi$.
\end{proof}

\begin{theorem}[a lower bound for the number of roots]\label{res:dual-arity-floor}
Fix radii $r_1,\ldots,r_p>0$ and weights $\sigma_1,\ldots,\sigma_p\ge0$, put
$\Sigma=\sum_i\sigma_i$ and
\[
 U=\sup_{d\ge d_{\mathrm{low}}(a)}
   \Bigl[\lambda(d)-\sum_i\sigma_i\,w(d,r_i)\Bigr],
 \qquad
 \lambda\bigl(d_{\mathrm{low}}(a)\bigr)=\frac{\delta(a)}2 .
\]
If $U>0$, then failure forces $k\ge(x-\pi\Sigma)/U$.
\end{theorem}
\leannote{Lean: \lproof{ErdosProblems/Erdos1041/PaperCompleteR21/HyperbolicLawOfCosines.lean}{594}{dual_arity_floor_sup}, \lproof{ErdosProblems/Erdos1041/PaperCompleteR21/HyperbolicLawOfCosines.lean}{579}{dual_arity_floor}, \lproof{ErdosProblems/Erdos1041/PaperCompleteR21/HyperbolicLawOfCosines.lean}{506}{le_of_lam_le}, \lproof{ErdosProblems/Erdos1041/PaperCompleteR21/HyperbolicLawOfCosines.lean}{141}{cosh_dist_polar}, and 1 further declaration. Conditional on the radius and budget bound that this proof derives from the standing failure hypothesis; see the \hyperref[sec:coverage]{coverage section}.}

\begin{proof}
By~\eqref{eq:lc-radius-and-budget}, Lemma~\ref{res:circle-slice-packing} and
the radius bound,
\[
 x\le\sum_j\lambda(d_j)
  =\sum_j\Bigl[\lambda(d_j)-\sum_i\sigma_i w(d_j,r_i)\Bigr]
    +\sum_i\sigma_i\sum_j w(d_j,r_i)
  \le kU+\pi\Sigma. \qedhere
\]
\end{proof}

The weighted sum of circle-slice inequalities is a dual certificate: it
bounds the contribution of every permitted root configuration without
having to optimise over all configurations directly. The weights give a
usable bound only when
$\Sigma$ and the supremum $U$ have rigorous bounds. In particular, checking
the expression on a finite grid does not bound the supremum between grid
points. Zero weights recover only the individual bound
$k\ge2x/\delta(a)$; nonzero weights incorporate the circle-packing
information and can improve the result.
A certificate at area $a'\ge a$ also applies at $a$. Indeed,
$d_{\mathrm{low}}(a)\ge d_{\mathrm{low}}(a')$, so the admissible range of
$d$ only shrinks. Also $D(a)\ge D(a')$, which enlarges every slice angle
$w(d,r_i)$ at fixed $d,r_i$. Since the weights are nonnegative, the
expression defining $U$ can only decrease. Thus the same certified pair
$(\Sigma,U)$ remains valid.
The numerical comparison also uses two bounds that do not involve chosen
weights. They are useful where the circle-slice certificate is weaker.
We derive them here so that every root-count input to the comparison is
visible.

\paragraph{Two further root-count bounds.}
Order the distances as $d_1\le\cdots\le d_k$. The triangle inequality gives
$d_j\ge D-d_1$ for $j\ge2$. The function $\lambda$ is decreasing and convex,
since $\lambda'(d)=-1/\sinh d$ and
$\lambda''(d)=\cosh d/\sinh^2d>0$. On
$[d_{\mathrm{low}},D/2]$, the convex function
$\lambda(d)+(k-1)\lambda(D-d)$ takes its maximum at an endpoint.
If $d_1\ge D/2$, each summand is at most $\lambda(D/2)$ instead.
(The definitions give $d_{\mathrm{low}}<D/2$.) Thus
\[
 x\le\sum_j\lambda(d_j)
 \le\max\left\{
  \frac{\delta(a)}2+(k-1)\lambda(D-d_{\mathrm{low}}),\;
  k\lambda(D/2)\right\},
\]
and consequently
\[
 k\ge\min\left\{
  1+\frac{x-\delta(a)/2}{\lambda(D-d_{\mathrm{low}})},\;
  \frac{x}{\lambda(D/2)}\right\}.
\]
The minimum is essential: the two expressions arise from alternative
positions of the nearest root, not from simultaneous restrictions.

For the second bound, integrate over the disjoint hyperbolic balls of
radius $D/2$. Put $E=e^{2/a}-1$, so each ball has hyperbolic area $2\pi E$.
At most one contains the origin; its contribution to
$\sum_j\lambda(d_j)$ is at most $\delta(a)/2$.
On every other ball, $-\log|z|$ is harmonic. To average it, first move
the ball's centre to $0$ by a disc automorphism. Hyperbolic area is
invariant under this change of coordinates and has radial density
$4/(1-|z|^2)^2$ there. The ordinary mean-value property on each centred
Euclidean circle therefore gives hyperbolic area mean
$\lambda(d_j)$. If the original ball has $0$ on its boundary, use
increasing smaller concentric balls and monotone convergence.

There are $m=k-1$ nonexceptional balls if one contains the origin, and
$m=k$ otherwise. Their union has area $2\pi mE$. Each superlevel set of
$-\log|z|$ is a centred ball, so its intersection with any set of this
area has area at most the smaller of the two areas. The centred ball
of area $2\pi mE$ attains that upper bound at every superlevel.
Integrating these intersection areas over the superlevel parameter
therefore bounds the integral over the union by the integral over that
centred ball. Its radius $R$ satisfies $\cosh R=1+mE$.
With $v=\cosh r-1$, the identity
$-\log\tanh(r/2)=\tfrac12\log(1+2/v)$ reduces the radial integral to
$(2E)^{-1}\int_0^{mE}\log(1+2/v)\,dv$. Evaluating it gives
\[
 \frac1{2\pi E}\int_{B_{\mathrm{hyp}}(0,R)}-\log|z|\,dA_{\mathrm{hyp}}
 =\frac{\frac{mE}{2}\log(1+2/(mE))+\log(1+mE/2)}{E}
 \le\frac{1+\log(1+mE/2)}{E},
\]
where the improper integral is finite because $v\log v\to0$ as
$v\downarrow0$, and the last inequality uses $\log(1+u)\le u$ for $u>0$.
Adding the possible exceptional contribution and inverting each case yields
\[
 k\ge\min\left\{
  1+\frac2E\bigl(e^{E(x-\delta(a)/2)-1}-1\bigr),\;
  \frac2E\bigl(e^{Ex-1}-1\bigr)\right\}.
\]
Again the minimum accounts for the two possible cases. This estimate uses
area rearrangement, whereas the preceding one uses only ordered distances;
neither replaces the circle-slice inequalities.

The comparison takes the maximum of $2$, $2x/\delta(a)$, these two bounds,
and the certified circle-slice bounds. Each may first be rounded up because
$k$ is an integer. Any of these estimates derived with an upper bound for
$a$ remains valid at the actual smaller area: the root configuration still
satisfies the weaker separation and radius conditions at that upper bound.

\paragraph{From the root count to an area contradiction.}
At a direction avoiding the finitely many critical-value arguments, lift the
value radius from $0$ to $te^{i\theta}$ from each of the $k$ roots in $C_t$
and split each lift at level $\mu$. Put
$A_0=\operatorname{Area}(C_t\cap\{|f|<\mu\})$. This is the sum of the
areas of the $k$ one-root components inside $C_t$, not of all the root
components of $f$; in particular, $0\le A_0\le\pi a$.
On each of these components $f/\mu$ is conformal, since it is a proper
map with no critical point onto a simply connected disc. Take
$\psi(z)=\sum_{\ell\ge0}\gamma_\ell z^\ell$ to be its inverse, so
$f(\psi(z))=\mu z$ and $\psi(0)$ is the selected root. The radial curves in
this coordinate therefore trace exactly the required value-ray lifts. Cauchy--Schwarz
in the radial variable and Parseval give
\[
 \operatorname{mean}_\theta
   \Bigl(\int_0^1|\psi'(re^{i\theta})|\,dr\Bigr)^2
 \le\sum_{\ell\ge1}\frac{\ell^2|\gamma_\ell|^2}{2\ell-1}
 \le\sum_{\ell\ge1}\ell|\gamma_\ell|^2
 =\frac{\operatorname{Area}(\text{component})}\pi.
\]
First stop the radial integrals short of the unit circle, then pass to
the endpoint by monotone convergence. Thus no smoothness
of the boundary at level $\mu$ is assumed. Summing the $k$ square-root
area bounds gives a mean total low lift length at most $\sqrt{kA_0/\pi}$.
At a regular intermediate level $u$, let $A_C(u)$ be the area of
$C_t\cap\{|f|<u\}$ and $P_C(u)$ its total boundary length.
The argument principle and the coarea formula give
\begin{equation}\label{eq:lc-perimeter}
 P_C(u)^2\le2\pi k\,u\,A_C'(u),
\end{equation}
because $|dz|=u\,d(\arg f)/|f'|$ on the level curve, its total argument
variation is $2\pi k$, and Cauchy--Schwarz applies.  Integrating
$P_C(u)/(2\pi u)$ from $\mu$ to $t$ bounds the mean high lift length by
$\sqrt{kx(\pi a-A_0)/(2\pi)}$. Cauchy--Schwarz combines the low and high
bounds as
\[
 \sqrt{\frac{kA_0}{\pi}}+
 \sqrt{\frac{kx(\pi a-A_0)}{2\pi}}
 \le\sqrt{ka(1+x/2)}.
\]
Choose a direction whose total lift length is at most
$\sqrt{ka(x+2)/2}$. Order its $k$ boundary endpoints cyclically and
join each adjacent pair using the two lifts and the intervening boundary
arc. Each path lies in $\overline{C_t}\subset\{|f|<1\}$ and, under
failure, has length at least $2$. Each lift occurs twice in the sum of
these lengths, while the boundary occurs once. Therefore
\[
 2k\le\sqrt{2ka(x+2)}+P_C(t).
\]
Now let $t=\mu e^x$ vary, and write
$a(x)=\operatorname{Area}(C_{\mu e^x})/\pi$ for the component containing
the chosen first pair; its root count $k$ may increase at later mergers.
On a regular interval, $a'(x)=tA_C'(t)/\pi$. Applying
\eqref{eq:lc-perimeter} at the outer level and rearranging therefore gives
\begin{equation}\label{eq:lc-area-growth}
 a'(x)\ \ge\ \frac1{2\pi^2}
 \Bigl[\,2\sqrt k-\sqrt{2a(x)(x+2)}\,\Bigr]_{+}^{2}
\end{equation}
on each regular interval for the component containing the selected pair,
using the largest of the root-count bounds just derived. At merger levels
this component only gains
area, which strengthens the integrated comparison.  P\'olya's inequality caps
$a$ at $1$ until the level reaches $1$, so a trajectory forced past that cap
before $t=1$ contradicts the assumption.

No positive initial area has to be assumed. While $a\le1$, for
$0<x\le3/10^5$ the bound $k\ge2$ alone gives
\[
 a'(x)\ge\frac1{2\pi^2}
 \left(2\sqrt2-\sqrt{2(2+3/10^5)}\right)^2>\frac1{30}
\]
at regular levels. For example, the last inequality follows using
$\pi<22/7$, $\sqrt2>140/99$ and
$\sqrt{2(2+3/10^5)}<20001/10000$. Integration, including the
nonnegative merger jumps, gives $a(3/10^5)>10^{-6}$. This is the
universal starting bound used by the certificate.

The comparison is integrated cell by cell without assuming that the
independently certified lower bounds increase with the table index.
On a cell $[x_\ell,x_r]$, suppose $a(x_\ell)\ge a_\iota$ and choose
a trial upper bound $m$ for the area. Let $g$ be a certified lower bound
for the right side of~\eqref{eq:lc-area-growth} throughout that cell
under the trial assumption $a\le m$. If
$m<a_\iota+(x_r-x_\ell)g$, then $a(x_r)>m$. Indeed, the contrary
assumption $a(x_r)\le m$ would imply $a\le m$ throughout the cell by
monotonicity. Integrating the differential inequality, with the
nonnegative merger jumps, would then give $a(x_r)>m$, a contradiction.

\paragraph{Certificate.}
The comparison gives an upper bound $X$ on the logarithmic time needed
for the forced area to exceed $1$. It therefore gives a contradiction
whenever $\mu e^X<1$. The stated constant $13/25$ is a convenient rational
choice with room to spare, not an optimality claim. The companion program
\href{https://github.com/wcook04/plectis-erdos/blob/f36a98bf3d3e6f65f1074e3b800e3293d5b8a51a/ErdosProblems/Erdos1041/scripts/check_erdos1041_angular_budget_closure.py}{\texttt{check\_erdos1041\_angular\_budget\_closure.py}} evaluates
every accepted inequality in exact rational arithmetic, with directed rounding
on the transcendental evaluations. Its fixed rational brackets for $\pi$
and $\log2$ have different elementary checks: Machin's identity
$\pi=16\arctan(1/5)-4\arctan(1/239)$ for the former, and
$\log2=2\sum_{j\ge0}3^{-(2j+1)}/(2j+1)$ for the latter.
Alternating-series bounds and a geometric tail bound, respectively,
suffice to verify them. The full replay uses $18$ area-table levels,
$14$ radii in each weighted bound,
$126$ attempted choices of weights, maximum step $1/400$, and the
initial lower area $10^{-6}$ at $x=3/10^5$. Only accepted rational
certificates enter the comparison; its early steps may be shorter than
$1/400$. The recorded full replay
certifies
\[
 X=\frac{635762889599}{1000000000000}<0.6357629,
 \qquad \frac{13}{25}e^{X}<1,
\]
which gives Theorem~\ref{res:low-critical-thirteen-twentyfifths}.
The recorded quick-mode replay gives $X=664373027131/1000000000000$
and the weaker threshold $51/100$.
The floating optimiser proposes weights but supplies no proof. The
argument uses only the subsequently certified rational bounds on their
sum and on the supremum $U$.  Each accepted
pair supplies a weight sum $\Sigma$ and a certified upper bound for the
supremum $U$ in Theorem~\ref{res:dual-arity-floor}. For orientation, the
recorded quick-mode pairs at $a=1$, rounded to six decimal places, are
$(0.085674,0.045865)$, $(0.126361,0.021829)$, $(0.163157,0.010513)$ and
$(0.208928,0.004318)$, giving $kU+\pi\Sigma$ approximately
$0.4526$, $0.5716$, $0.6808$ and $0.7945$ at $k=4,8,16,32$, respectively.
These rounded pairs do not themselves certify an inequality; the comparison
uses the unrounded rational values and their directed bounds.
To bound the supremum over $d$, the checker subdivides intervals and bounds
each term in the direction needed for an upper bound. The function
$\lambda$ is decreasing. For fixed $r$, the slice angle $w(d,r)$ has no
strict interior minimum. Indeed, before clipping the angle to $0$ or $\pi$,
its cosine is
\[
 H(d)=\frac{\cosh d\cosh r-\cosh(D/2)}{\sinh d\sinh r},\qquad
 H'(d)=\frac{\cosh(D/2)\cosh d-\cosh r}{\sinh^2d\sinh r}.
\]
The numerator of $H'$ is increasing. Thus $H$ has at most one interior
minimum, and $w=\arccos H$ has at most one interior maximum; clipping
preserves the endpoint-minimum property. On an interval $[d_-,d_+]$,
therefore,
\[
 \lambda(d)-\sum_i\sigma_iw(d,r_i)
 \le\lambda(d_-)-\sum_i\sigma_i
       \min\{w(d_-,r_i),w(d_+,r_i)\}.
\]
Beyond $d=\max_i r_i+D/2$, all the slice angles vanish, so only the
decreasing function $\lambda$ remains. These are interval bounds, not
values sampled on a grid.
The closing inequality is independently checkable: with
$X=635762889599/10^{12}$ and
$e^X\le\sum_{j\le14}X^j/j!+(X^{15}/15!)/(1-X/16)$, rational arithmetic gives
$(13/25)e^X<0.982000386<1$. The same majorant also gives
\[
 \frac{529}{1000}e^X<0.998996547<1.
\]
Thus the recorded full certificate in fact gives the path conclusion for
$\mu\le529/1000$, without changing the geometric argument or its hypotheses.
This is an arithmetic consequence of the recorded stopping time, not a new
certificate replay. We retain $13/25$ in the theorem and its scaling
corollaries as the simpler constant.

\paragraph{Sources and scope.}
\label{bdry:low-critical-closure}
The analytic chain above is ordinary mathematics.  Its general inputs are the
Riemann mapping theorem, the Bergman kernel, the argument principle, the
coarea formula, and P\'olya's area inequality
$\operatorname{Area}\{|f|\le t\}\le\pi t^{2/n}$
\cite[printed pp.~280--282]{polya1928}, \cite[Theorem~1]{crane}.  The exact
rational certificate is checked by a separate exact-arithmetic program,
not by the Lean kernel. The two additional root-count bounds above are
the ordered-distance estimate in
\href{https://github.com/wcook04/plectis-erdos/blob/f36a98bf3d3e6f65f1074e3b800e3293d5b8a51a/ErdosProblems/Erdos1041/ClusterSeparationLowCriticalClosure.md}{the earlier distance comparison}
and the area-rearrangement estimate in
\href{https://github.com/wcook04/plectis-erdos/blob/f36a98bf3d3e6f65f1074e3b800e3293d5b8a51a/ErdosProblems/Erdos1041/HyperbolicPackingArityFloor.md}{the hyperbolic packing argument}.
Their derivations above are part of the analytic proof. No complete
Lean proof or independent review of this theorem is claimed. The finite summation implication identified in
the companion paper assumes the geometric inequalities; it does not
formalise the analytic chain or the certificate replay. Prior art for
the assembled statement is unassessed. No novelty is claimed for the
slice inequality, which follows directly from disjointness of the balls
and the hyperbolic law of cosines.  The public
source of record is
\href{https://github.com/wcook04/plectis-erdos/blob/f36a98bf3d3e6f65f1074e3b800e3293d5b8a51a/ErdosProblems/Erdos1041/AngularBudgetLowCriticalClosure.md}{\texttt{AngularBudgetLowCriticalClosure.md}}, which
carries the full analytic argument, the table of certified weighted bounds
and the instructions for replaying the certificate. The stopping-time test
requires $\mu<e^{-X}$; the stated $13/25$ theorem and the $529/1000$
consequence are two rational cutoffs within that range. Other theorems
apply to some polynomials outside it, so failure of this numerical test
does not identify the class left untreated by all the results.
Fixed-degree arguments remain stronger at $n=4$ and $n=5$, where the
corresponding thresholds are $61/100$ and $139/250$; from $n=6$ on the
all-degree constant $13/25$ is the better statement.  These threshold
arguments do not settle the unrestricted root-connector conclusion.

\begin{corollary}[scaled low-critical connection]\label{res:scaled-low-critical-path}
Every squarefree monic polynomial $f$ of degree $n\ge2$ has two distinct
zeros joined by a rectifiable curve of length less than
$(5/2)\mu^{1/n}$ in $\{|f|<(25/13)\mu\}$, where
$\mu=\min_{f'(c)=0}|f(c)|$.
\end{corollary}
\leannote{Lean: \lproof{ErdosProblems/Erdos1041/PaperCompleteR21/LowCriticalScaleTransport.lean}{491}{scaledLowCriticalFiveHalves_of_lowCritical}. Conditional on Theorem~\ref{res:low-critical-thirteen-twentyfifths} as stated, and only in degrees above two; see the \hyperref[sec:coverage]{coverage section}.}
\begin{proof}
For $n\ge3$, apply Theorem~\ref{res:low-critical-thirteen-twentyfifths}
to $g(z)=s^{-n}f(sz)$ with $s=((25/13)\mu)^{1/n}$.
Its least critical modulus is $13/25$; rescaling gives length less than
$2(25/13)^{1/n}\mu^{1/n}<(5/2)\mu^{1/n}$, since
$25/13<(5/4)^3$.  For $n=2$, write $f(z)=(z-h)^2-d^2$.
The two radial segments through $h$ have total length $2|d|=2\mu^{1/2}$
and lie in the closed level $\mu$, which is inside the stated open level.
\end{proof}
The next section gives a different construction whose estimates retain
the root count and capacity of the chosen component.

\section{A path estimate from area and boundary length}
\label{sec:constant-factor}

Corollary~\ref{res:scaled-low-critical-path} has both a smaller length
constant and a smaller containment level than the next theorem:
$5/2<71/10$ and $25/13<2$. The independent area argument below is useful
for a different reason: it keeps the root count and component capacity
in the length estimate. Its $\mu\le1/2$ corollaries illustrate that
geometric dependence; their ranges are already covered by
Theorem~\ref{res:low-critical-thirteen-twentyfifths}.  For a monic degree-$n$ polynomial put
\[
 K_t=\{z:|f(z)|\le t\},\qquad
 \mu=\min_{f'(c)=0}|f(c)|,\qquad \rho=\mu^{1/n}.
\]

\begin{theorem}[a uniform path bound at level $2\mu$]
\label{res:constant-factor-path}
For every monic polynomial $f$ of degree $n\ge2$, two zero occurrences are
joined by a possibly degenerate path of length at most
\[
 \frac{71}{10}\,\rho
\]
inside $K_{2\mu}$.  If $f$ is squarefree, their locations are distinct.
If $\mu\le1/2$, the construction may be chosen inside $\{|f|<1\}$ with
length at most $5.7$.
\end{theorem}
\leannote{Lean: \lproof{ErdosProblems/Erdos1041/PaperCompleteR21/ConstantFactorAreaCriteria.lean}{677}{cfa_constant_factor_path}, \lproof{ErdosProblems/Erdos1041/PaperCompleteR21/ConstantFactorAreaCriteria.lean}{355}{cfaBracket_two_three_twentieths_lt}, \lproof{ErdosProblems/Erdos1041/PaperCompleteR21/ConstantFactorAreaCriteria.lean}{393}{cfaBracket_five_point_seven}, \lproof{ErdosProblems/Erdos1041/PaperCompleteR21/ConstantFactorAreaCriteria.lean}{313}{cfa_degenerate}, and 3 further declarations. Conditional on the level and direction averaging construction that this proof produces; see the \hyperref[sec:coverage]{coverage section}.}

The constant is the rationally certified specialization of a two-parameter
bound.  If the selected component contains $k\ge2$ roots, then for
every $r\in(0,1)$ and $\lambda>1$ the proof constructs a path in
$K_{\lambda\mu}$ whose length is at most
\[
 \sqrt{\frac2k}\left(
   \frac{\sqrt2\,r}{(1-r)^2}
   +\lambda^{1/n}\left(
      \sqrt{\log(\lambda/r)}+
      \frac{\pi}{\sqrt{\log\lambda}}
    \right)
 \right)\rho.                                      \tag{CF}
\]
For $n\ge3$, the choice $\lambda=2$, $r=3/20$ makes the bracket less
than $71/10$: the largest case is $n=3$, where exact rational bounds give
$66517563/9392500<71/10$.  Degree two has the exact root-segment bound
$2\rho$.

\begin{proof}
A repeated zero gives the constant path, so assume $f$ is squarefree and
$\mu>0$. Degree two is the segment of length $2\rho$ in $K_\mu$; assume
$n\ge3$. Put $T=\lambda\mu$, let $C_T$ be the component of $K_T$
containing a critical point at level $\mu$, and write
$A(\sigma)=\operatorname{Area}(K_\sigma\cap C_T)$.
We will first choose a regular level $t$, then a direction in the value
plane, and finally an adjacent pair of roots. This order ensures that
the three estimates below apply to the same path.

At a regular level $\sigma$, a component $C'$ containing $k'$ roots has
$|dz|=\sigma\,d\varphi/|f'|$ on its boundary and total argument variation
$2\pi k'$. Cauchy--Schwarz, followed by the coarea formula, gives
\[
 \mathcal H^1(\partial C')^2\le2\pi k'\,\sigma A'(\sigma).
                                                        \tag{CF1}
\]
Here $A'$ sums the area derivatives of all the components in $C_T$,
so it bounds the contribution from $C'$. P\'olya's inequality
$A(T)\le\pi T^{2/n}$ \cite[Theorem~1]{crane} and averaging with respect
to $d\sigma/\sigma$ on $(\mu,T)$ give a regular level $t$ such that
$tA'(t)\le\pi T^{2/n}/\log\lambda$. The component $C_t$ containing
the chosen first merger has $k\ge2$ roots and
\[
 \mathcal H^1(\partial C_t)
 \le\pi\sqrt{2k/\log\lambda}\,T^{1/n}.
\]
Averaging over levels avoids estimating the area derivative at the
critical level itself.

Choose a direction avoiding all critical-value arguments. Lift its radial
segment from each of the $k$ roots to $\partial C_t$, splitting each lift
at value modulus $r\mu$. Below $\mu$, the inverse branches of $f/\mu$
are univalent. Their derivatives at $0$ are $\mu/f'(z_i)$, where $z_i$
ranges over the roots. The area formula for their disjoint images gives
\[
 \sum_i\frac{\mu^2}{|f'(z_i)|^2}\le\rho^2.        \tag{CF2}
\]
Indeed, the squared derivative at the centre is at most the image area
divided by $\pi$, and the total image area is at most $\pi\mu^{2/n}$.
Koebe distortion and Cauchy--Schwarz therefore bound the sum of the
lengths below $r\mu$, uniformly in the direction. Write
$\ell_i^{\rm high}(\theta)$ for the length of the $i$th piece above
$r\mu$ in direction $\theta$. At an intermediate regular level, apply
(CF1) to all the lower components contained in $C_t$; their root counts
sum to $k$. Integrating their total perimeter with measure
$d\sigma/(2\pi\sigma)$ gives the mean of $\sum_i\ell_i^{\rm high}$.
Cauchy--Schwarz in $\sigma$, using an area increment at most
$\operatorname{Area}(C_t)\le\pi T^{2/n}$, gives
\[
 \frac1{2\pi}\int_0^{2\pi}\sum_i\ell_i^{\rm high}(\theta)\,d\theta
 \le\sqrt{\frac{k\operatorname{Area}(C_t)}{2\pi}
                   \log\frac{t}{r\mu}}
 \le\sqrt{\frac k2}\,T^{1/n}\sqrt{\log(\lambda/r)}.
\]
Choose a direction avoiding the critical-value arguments at which the
sum is at most this bound. The excluded directions form a finite set
and hence have measure zero.

Order its $k$ endpoints cyclically on $\partial C_t$. Each adjacent pair
of endpoints gives a path consisting of two lifted segments and the
intervening boundary arc. In the sum over these $k$ paths, every lifted
segment occurs twice and the boundary arcs partition $\partial C_t$.
Thus some pair has length at most the sum of the following three bounds:
\[
\begin{aligned}
 \frac2k\sum_i\ell_i^{\rm low}
 &\le\frac{2r\rho}{\sqrt{k}(1-r)^2},\\
 \frac2k\sum_i\ell_i^{\rm high}
 &\le\sqrt{\frac2k}\sqrt{\log(\lambda/r)}\,T^{1/n},\\
 \frac{\mathcal H^1(\partial C_t)}k
 &\le\sqrt{\frac2k}\frac{\pi T^{1/n}}{\sqrt{\log\lambda}}.
\end{aligned}
\]
The first holds in every direction, the second in the direction just
chosen, and the third at the previously selected level. Substituting
$T^{1/n}=\lambda^{1/n}\rho$ proves~(CF).
For the final assertion, the repeated-zero case is already a constant path,
and in degree two the root segment has length
$2\rho=2\sqrt\mu\le\sqrt2$.  For $n\ge3$, take $\lambda=2$ and $r=3/20$.
Since $k\ge2$ and $\mu\le1/2$, (CF) gives
\[
 \begin{aligned}
 \operatorname{length}
 &\le \frac{60\sqrt2}{289}\mu^{1/n}
 +(2\mu)^{1/n}\left(
   \sqrt{\log(40/3)}+\frac{\pi}{\sqrt{\log2}}
 \right)\\
 &\le \frac{60\sqrt2}{289}+\sqrt{\log(40/3)}
      +\frac{\pi}{\sqrt{\log2}}\\
 &=5.676476\ldots<5.7.
 \end{aligned}
\]
Here $2\mu\le1$, and the mean-value choice may be taken at a regular level
strictly below the top of its positive-measure window.  Thus the entire path
lies in $\{|f|<1\}$.
\end{proof}

The bracket~(CF) retains two quantities that the constant $71/10$ discards,
namely the root count $k$ of the selected component through the factor
$\sqrt{2/k}$, and the capacity of that component through the area input.
Keeping either one turns the constant-factor theorem into the target
conclusion on an explicit region.

The number of roots in the first merged component can be much smaller
than the degree. At a first merger caused by one simple critical point it
is $2$, so the following thresholds do not apply. For $z^n-b$, however,
all $n$ roots merge at the same level, giving $k_0=n$. For example,
$n\ge17$ and $0<|b|\le1/2$ satisfy the first pair of hypotheses. These
conditions illustrate what the area argument gains from simultaneous
mergers; the preceding $13/25$ theorem already covers these three ranges
without requiring a large root count.

\begin{corollary}[a criterion using the number of roots at the first merger]
\label{res:constant-factor-arity}
Let $f$ be monic with every root in the open unit disc, let $c_*$ be a
critical point with $|f(c_*)|=\mu$, and let $k_0$ be the number of roots,
counted with multiplicity, in the component of $K_\mu$ containing $c_*$.
Then Erd\H{o}s~\#1041 holds for $f$ in each of the three cases
\[
 \mu\le\tfrac12\ \text{and}\ k_0\ge17,\qquad
 \mu\le\tfrac14\ \text{and}\ k_0\ge12,\qquad
 \mu\le\tfrac18\ \text{and}\ k_0\ge10 .
\]
\end{corollary}
\leannote{Lean: \lproof{ErdosProblems/Erdos1041/PaperCompleteR21/ConstantFactorAreaCriteria.lean}{765}{cfa_arity_criterion}, \lproof{ErdosProblems/Erdos1041/PaperCompleteR21/ConstantFactorAreaCriteria.lean}{440}{cfaArityBracket_case_one}, \lproof{ErdosProblems/Erdos1041/PaperCompleteR21/ConstantFactorAreaCriteria.lean}{455}{cfaArityBracket_case_two}, \lproof{ErdosProblems/Erdos1041/PaperCompleteR21/ConstantFactorAreaCriteria.lean}{470}{cfaArityBracket_case_three}, and 2 further declarations. Conditional on the level and direction averaging construction in the proof of Theorem~\ref{res:constant-factor-path}; see the \hyperref[sec:coverage]{coverage section}.}

\begin{proof}
If $\mu=0$, a repeated zero gives the constant path. Assume henceforth
$\mu>0$. Every selected component $C_t$ in the proof of
Theorem~\ref{res:constant-factor-path} contains the first-merge component, so
$k\ge k_0$.  For the first case take $\lambda=2$, $r=13/100$; then
$(2\mu)^{1/n}\le1$ and $\rho\le1$, and the exact bounds
$\sqrt2<283/200$, $\sqrt{\log(200/13)}<5/3$ and
$\pi/\sqrt{\log2}<(22/7)/(104/125)=1375/364$ make the bracket in~(CF) at most
$15668813/2755116$, whose square is $34-12570881068535/7590664173456<34$.
Hence $k_0\ge17$ gives squared length below $(2/17)\cdot34=4$.  For the second
case take $\lambda=4$, $r=3/25$: the bracket is below $6075221/1273888$, whose
square is $24-2038665078215/1622790636544<24$, and $2/k_0\le1/6$ gives squared
length below $4$.  For the third take $\lambda=8$, $r=11/100$: the bracket is
below $55629121/12475575$, whose square is
$20-18200328379859/155639971580625<20$, and $2/k_0\le1/5$ again gives squared
length below $4$.  In each case $\lambda\mu\le1$, so the freedom in
the choice of the regular level $t$ keeps $t<1$ and the containment strict.
\end{proof}

The next criterion uses logarithmic capacity, a measure of the size of a
compact set that scales linearly under dilation. It compares the size of
one component with that of the entire polynomial sublevel set. If that
sublevel set is connected, the ratio $\kappa$ below is $1$; a cutoff such
as $\kappa\le1/3$ therefore requires a genuinely smaller component. This
is additional geometric information, not a consequence of a small number
of roots in the component. The ratio requires $\mu>0$. When $\mu=0$,
$f$ has a repeated zero and the constant path already gives the conclusion;
there is no capacity ratio to evaluate.

\begin{corollary}[a criterion using component capacity]
\label{res:constant-factor-capacity}
Keep the hypotheses of Corollary~\ref{res:constant-factor-arity} with
$0<\mu\le1/2$, let $C$ be the component of $\{|f|<2\mu\}$ containing $c_*$, and
put $\kappa=\operatorname{cap}(\overline C)/(2\mu)^{1/n}$.  If
$\kappa\le\tau_{k_0}$, where
\[
 \tau_k=\frac{\sqrt{2k}-A}{B},\qquad
 A=\frac{283}{3610},\qquad B=\frac{52029}{9100},
\]
then Erd\H{o}s~\#1041 holds for $f$. In particular $\kappa\le1/3$ suffices for
every root count $k_0\ge2$, and the rational cutoffs
$2/5,\,12/25,\,1/2,\,7/12,\,16/25,\,2/3,\,7/10$ suffice at
$k_0=3,\ldots,9$, rising to $39/40$ at $k_0=16$.
\end{corollary}
\leannote{Lean: \lproof{ErdosProblems/Erdos1041/PaperCompleteR21/ConstantFactorAreaCriteria.lean}{538}{cfa_capacity_criterion}, \lproof{ErdosProblems/Erdos1041/PaperCompleteR21/ConstantFactorAreaCriteria.lean}{583}{cfaTau_third}, \lproof{ErdosProblems/Erdos1041/PaperCompleteR21/ConstantFactorAreaCriteria.lean}{591}{cfaTau_cutoffs}, \lproof{ErdosProblems/Erdos1041/PaperCompleteR21/ConstantFactorAreaCriteria.lean}{572}{cfaTau_ge_of_sq}, and 2 further declarations. Conditional on the level and direction averaging construction in the proof of Theorem~\ref{res:constant-factor-path}, together with the area-capacity inequality; see the \hyperref[sec:coverage]{coverage section}.}

\begin{proof}
Repeat the averaging proof of Theorem~\ref{res:constant-factor-path} with
$A(\sigma)=\operatorname{Area}(K_\sigma\cap C)$ for $\sigma<2\mu$.
This keeps the chosen paths in the open component $C$, even if distinct
components have touching boundaries at level $2\mu$. Replace the global
area input by the component form
$\operatorname{Area}(C)\le\pi\operatorname{cap}(\overline C)^2
 =\pi\kappa^2(2\mu)^{2/n}$ of the area--capacity inequality
\cite[Theorem~6]{crane}.  Only the high-lift and boundary terms acquire the
factor $\kappa$; the low-lift term~(CF2) is unchanged.  Taking $\lambda=2$,
$r=1/20$, and the exact bounds above together with
$\log40<12641/3402<(97/50)^2$, gives
$\operatorname{length}<\sqrt{2/k_0}\,(A+B\kappa)$.  The definition of
$\tau_{k_0}$ makes the right side at most $2$, and for each displayed rational
$q_k$ integer arithmetic gives $(A+Bq_k)^2<2k$.  Discarding
$\sqrt{2/k_0}\le1$ altogether gives the uniform cutoff $\kappa\le1/3$.
\end{proof}

By the component-capacity formula, $\kappa=e^{-\Sigma/n}$ with $\Sigma$ the sum
of exterior Green function values at the roots excluded from $C$.  Indeed, let
$\Omega$ be the unbounded component of
$\widehat{\mathbb C}\mathbin{\backslash}\overline C$,
let $G$ be its Green function with pole at infinity, extended by zero off
$\Omega$, and let $g(\cdot,a)$ be its Green function with pole at $a\in\Omega$.
The function $\frac1n\log\frac{|f|}{2\mu}-G+\frac1n\sum_a g(\cdot,a)$, summed
over the roots $a\in\Omega$ with multiplicity, extends harmonically to
$\Omega$, is bounded, and vanishes on $\partial\Omega\subseteq\partial C$, so it
vanishes identically; its value at infinity is $\log\kappa+\frac1n\sum_aG(a)$
by the symmetry $g(\infty,a)=G(a)$.  Failure of these component-sensitive
criteria alone, in the range $\mu\le1/2$, forces $k_0\le16$ and exterior
Green sum $\Sigma<n\log(40/39)$ when $k_0=16$.  This is a limitation of
the area-based criteria.  The low-critical theorem already excludes an actual
counterexample throughout that range.

The proof of Theorem~\ref{res:constant-factor-path} averages over
levels in $[\mu,\lambda\mu]$. To keep the path at the first critical
level instead, one would need a perimeter estimate of the following kind.
The regular-level argument~(CF1) gives no finite pointwise bound at the
critical level: even at a simple critical point, the area derivative
diverges logarithmically.

The addendum \cite[Proposition~2]{revision2026} reports a
Runge-approximation construction of monic level-one components with one
simple zero and arbitrarily long analytic boundaries. It imposes neither
a unit root disc nor a global lower bound on critical-value moduli.
That proof was not supplied or available for checking in this revision;
the report is not an input to the path estimates here. The question below
requires a stronger hypothesis than the existence of one unramified
component: the level must lie below every critical-value modulus of the
polynomial. A construction without that global condition cannot by itself
answer the question.

\begin{problem}[perimeter below the smallest critical-value modulus]\label{prob:conjecture-p}
Let $f$ be squarefree and monic of degree $n\ge2$, and put
$\mu=\min_{f'(c)=0}|f(c)|$.  Is there a constant $\beta$, independent of
$n,f$ and $0<\sigma<\mu$, such that every component $C$ of
$\{|f|\le\sigma\}$ satisfies
\[
 \mathcal H^1(\partial C)\le\beta\sigma^{1/n}?
\]
Equivalently, for each univalent inverse branch defined on $\mathbb D$
by $f(\phi_j(w))=\sigma w$, is
\[
 \int_0^{2\pi}|\phi_j'(e^{it})|\,dt\le\beta\sigma^{1/n}
 \qquad(1\le j\le n)?
\]
There is no root-location hypothesis.  If such constants exist, write
$\beta_*$ for their infimum.
\end{problem}
All these components contain one zero. The condition $\sigma<\mu$
controls all inverse branches at once; it is stronger than requiring
only the selected component to contain one zero. The examples below
compare particular constants, not the necessity of this global
hypothesis for every possible perimeter bound.
For $1<R<\mu/\sigma$, the area estimate gives
\[
 \frac{1}{2\pi}\int_0^{2\pi}|\phi_j'(e^{it})|\,dt
 \le \frac{(\sigma R)^{1/n}}{\sqrt{R^2-1}}.
\]
Indeed, circular means of $|\phi_j'|^2$ increase, the image of the annulus
$1<|w|<R$ has area at least $\pi(R^2-1)$ times the mean at $1$, and
P\'olya's preimage-area bound \cite[Theorem~1]{crane} bounds that area by
$\pi(\sigma R)^{2/n}$.  The loss as $R\downarrow1$ belongs to this
argument and does not prove an endpoint divergence.

For $d>0$ and $f(z)=z^2-d^2$, one has $\mu=d^2$ and $\rho=d$.  For $\sigma<d^2$ the component
of $\{|f|\le\sigma\}$ around $d$ contains one root. Its parametrisations
$d\sqrt{1+(\sigma/d^2)e^{i\theta}}$ converge uniformly, as $\sigma\uparrow d^2$,
to one loop of Bernoulli's lemniscate $|z^2-d^2|=d^2$, of length $d\,\Gamma(1/4)^2/(2\sqrt\pi)$ by the substitution
$\varphi(w)=d\sqrt{1+w}$ and $\Gamma(1/4)\Gamma(3/4)=\pi\sqrt2$.  At
$\sigma=d^2$ the two loops meet at the critical point, and the closed sublevel
set is a single component containing both roots. Lower semicontinuity
of length under the stated uniform convergence gives
$\beta_*\ge\Gamma(1/4)^2/(2\sqrt\pi)=3.70814935\ldots$;
convergence of the lengths themselves is not required.
For $z^n-r^n$, polar parametrisation of one limiting loop gives its
length divided by $r$ as
\[
 \frac{2^{1/n}}{n}\int_{-\pi/2}^{\pi/2}
     (\cos\theta)^{1/n-1}\,d\theta
 =\frac{2^{1/n}}{n}\sqrt\pi\,
     \frac{\Gamma(1/(2n))}{\Gamma(1/(2n)+1/2)}.
\]
The evaluation is Euler's beta integral
\href{https://dlmf.nist.gov/5.12.E1}{(5.12.1)}, in its
\href{https://dlmf.nist.gov/5.12.E2}{trigonometric form (5.12.2)}.
To verify the claimed decrease with $n$, put $p=1/n$. The expression is
$2^{p+1}\sqrt\pi\,\Gamma(1+p/2)/\Gamma((1+p)/2)$; its logarithmic
derivative is
\[
 \log2+\tfrac12\left(
    \frac{\Gamma'(1+p/2)}{\Gamma(1+p/2)}-
    \frac{\Gamma'((1+p)/2)}{\Gamma((1+p)/2)}\right)>0.
\]
Here $\Gamma'/\Gamma$ is increasing on the positive real axis, by the
positive trigamma series \href{https://dlmf.nist.gov/5.15.E1}{(5.15.1)}.
The expression therefore decreases with $n$ and tends to $2$.

The comparison with the next example also has an elementary bound,
independent of rounded gamma values. One loop of $z^2-d^2$ has
parametrisation
\[
 z(s)=\sqrt2\,d\,
   \frac{\cos s(1+i\sin s)}{1+\sin^2s},
 \qquad -\pi/2\le s\le\pi/2,
\]
whose speed is $\sqrt2\,d/\sqrt{1+\sin^2s}$. Convexity puts
$(1+x)^{-1/2}$ below its secant on $[0,1]$. Integrating that secant
at $x=\sin^2s$ bounds the loop length divided by $d$ by
$(\pi/2)(1+\sqrt2)$, as used below.

The quadratic loop is exceeded in degree eight, even with the global
subcritical hypothesis. For $p(z)=z^8-(3/2)z$, every critical point $c$
satisfies $c^7=3/16$ and $p(c)=-(21/16)c$. Hence
\[
 \mu=\frac{21}{16}\left(\frac3{16}\right)^{1/7}>1,
 \qquad \mu^7=\frac{3\cdot21^7}{16^8}>1.
\]
Thus $\sigma=1$ is below every critical-value modulus, not merely a
regular level for the selected component. The component $C$ of
$\{|p|\le1\}$ containing the origin lies in
$\{|z|<4/5\}$, because $|p(z)|\ge6/5-(4/5)^8>1$ on $|z|=4/5$, and it contains a
neighbourhood of the closed disc of radius $5/8$, because
$|p(z)|\le(5/8)^8+15/16<1$ on that disc.  The other roots satisfy
$|z|^7=3/2$, so $C$ contains exactly one root, and
$\mathcal H^1(\partial C)>5\pi/4>(\pi/2)(1+\sqrt2)\ge\Gamma(1/4)^2/(2\sqrt\pi)$.  Moreover,
for fixed $a>1$ and all sufficiently large $N$, the one-root component of
$\{|z^N-az|\le1\}$ contains every disc of radius $r<1/a$, and a circle of
radius between $1/a$ and $1$ isolates it by Rouch\'e's theorem; letting
$a\downarrow1$ shows that every admissible $\beta$ is at least $2\pi$.
For the latter family all critical values have modulus
$(1-1/N)a(a/N)^{1/(N-1)}\to a>1$, so it also gives the necessary bound
$\beta_*\ge2\pi$ under the corrected global subcritical hypothesis.
Kuznetsova--Tkachev \cite[Theorems~1--2]{kuznetsova2003} establish
Laplace-transform and log-convexity properties for regular level-length
functions.  Those results do not provide a uniform bound at the first
critical level.  Total lemniscate length and one-branch endpoint length are
different extremal quantities.

The preceding examples do not establish a universal perimeter bound.
The implication below applies to an individual polynomial; a constant
$\beta$ independent of the polynomial and degree would give a uniform
path bound. The construction uses two components meeting at a first
critical point and joins each root to that point at a cost of at most
half its component's perimeter. It does not give the constant $2$ of
the historical question.

\begin{theorem}[a path from a subcritical perimeter bound]
\label{res:conjecture-p-consumer}
Let $f$ be squarefree and monic of degree $n\ge2$, and put
$\mu=\min_{f'(c)=0}|f(c)|>0$. Suppose $\beta>0$ satisfies
\[
 \mathcal H^1(\partial C)\le\beta\sigma^{1/n}
\]
for every $0<\sigma<\mu$ and every component $C$ of $\{|f|\le\sigma\}$.
Then two distinct roots of $f$ are joined inside $K_\mu=\{|f|\le\mu\}$
by a rectifiable path of length at most $\beta\mu^{1/n}$.
\end{theorem}
\leannote{Lean: \lproof{ErdosProblems/Erdos1041/PaperCompleteR21/SubcriticalPerimeterPath.lean}{265}{subcritical_perimeter_path_paper}, \lproof{ErdosProblems/Erdos1041/PaperCompleteR21/SubcriticalPerimeterPath.lean}{245}{subcritical_perimeter_path}, \lproof{ErdosProblems/Erdos1041/PaperCompleteR21/SubcriticalPerimeterPath.lean}{203}{connectedAtMost_half_perimeter}, \lproof{ErdosProblems/Erdos1041/PaperCompleteR21/SubcriticalPerimeterPath.lean}{216}{halfPerimeterJoin_of_jordanArcDatum}, and 4 further declarations. Conditional on the two-component split at the first critical level that this proof constructs; see the \hyperref[sec:coverage]{coverage section}.}

\begin{proof}
\hypertarget{perimeter-to-path}{}
Put $\rho=\mu^{1/n}$. On each component of $\{|f|<\mu\}$,
$f$ is a proper unramified covering of the disc $\{|w|<\mu\}$, hence
univalent. Equivalently, exhaust by regular levels below $\mu$ and apply
the component count \cite[proof of Proposition~2.1]{eks2010}.
At a critical point $c_*$ with
$|f(c_*)|=\mu$, the local degree $d\ge2$ gives $d$ nearby preimages of
an inward radial value. Univalence places them in different one-root
components. Choose two such components $U_a,U_b$, containing roots $a,b$
and with $c_*$ in both closures; simplicity of $c_*$ is unnecessary.
For the boundary passage, take $f(\phi(w))=\mu w$ on $\mathbb D$.
Fix $|w_0|=1$. The finitely many preimages of $\mu w_0$ have disjoint
neighbourhoods. Properness of $f$ and connectedness of a sufficiently
small disc cap force $\phi$ to stay in one of them as $w\to w_0$.
At its preimage $c$, write
\[
 f(z)-f(c)=(z-c)^d h(z),\qquad h(c)\ne0.
\]
A local analytic $d$th root of $h$ makes $(z-c)h(z)^{1/d}$ a conformal
coordinate. In each inverse sector, $\phi$ is therefore an analytic
function of $(w-w_0)^{1/d}$ and extends continuously to $c$.
The finitely many critical boundary values are the only exceptional
points; at every other boundary value the ordinary inverse function
theorem applies. Moreover, near $e^{i\theta_0}=w_0$ the boundary
speed is $O(|\theta-\theta_0|^{1/d-1})$, which is integrable even when
$d>2$. Thus the extension is continuous on the closed disc and its
boundary curve is rectifiable. The identity
$f(\phi(e^{i\theta}))=\mu e^{i\theta}$ makes the boundary values injective,
so the curve is Jordan.

As $r\uparrow1$, the curves $\phi(re^{i\theta})$ converge uniformly to
that boundary and have length at most $\beta(\mu r)^{1/n}$ by the
hypothesis. Lower semicontinuity of length gives boundary length at most
$\beta\rho$ for each component.

A bounded rectifiable Jordan domain $U$ has the following elementary
property: any $a\in U$ can be joined to any $c\in\partial U$ in
$\overline U$ with length at most $\mathcal H^1(\partial U)/2$.
Choose a line through $a$ but not $c$, and let $[u,v]$ be the closure
of the connected interval in its intersection with $U$ containing $a$.
Of the two boundary arcs from $u$ to $v$, write $A$ for the one containing
$c$ and $A'$ for the other. The paths from $a$ through $u$ or $v$ and
then along $A$ to $c$ have total length
\[
 |a-u|+|a-v|+\operatorname{length}(A)
 =|u-v|+\operatorname{length}(A)
 \le\mathcal H^1(\partial U),
\]
because $|u-v|\le\operatorname{length}(A')$. One path is at most half
this sum. The connected interval ensures containment of the straight
parts even when $U$ is not convex.

Apply this property in $U_a,U_b$ with common boundary point $c_*$.
Concatenating the two paths gives length at most
$\tfrac12(\mathcal H^1(\partial U_a)+\mathcal H^1(\partial U_b))
\le\beta\rho$, as required.
Only these two component perimeters are used; the subcritical hypothesis
bounds them by the limiting argument above. Boundary arcs are allowed,
so containment is in $K_\mu$, not the open sublevel.
\end{proof}

\paragraph{Sources and scope.}
\label{bdry:constant-factor}
Theorem~\ref{res:constant-factor-path} is ordinary mathematics.  Its degenerate
case $\mu=0$, its degree-two case and every numerical constant in its statement
are checked in Lean; the analytic construction for degree at least three is
used there as a hypothesis (Section~\ref{sec:coverage}).  The inputs of that
construction are the capacity identity for polynomial
preimages \cite[\S2]{crane}, P\'olya's area inequality \cite[Theorem~1]{crane},
Cauchy--Schwarz, the Koebe distortion theorem for univalent maps, and
averaging over adjacent boundary endpoints.  The public source
of record is
\href{https://github.com/wcook04/plectis-erdos/blob/f214a6b45528dc5eefe20ffadc35f2e981627d4c/research_corpus/Erdos1041/UnconditionalConstantFactorBound.md}{the full area and boundary-length argument}, which
carries every rational verification quoted above. Its Section~10 gives
a weaker conditional bound involving an auxiliary parameter and a
logarithm; Theorem~\ref{res:conjecture-p-consumer} states the stronger
$\beta\rho$ consequence proved here. That section's proposed numerical
perimeter constant is refuted by the degree-eight example above and is
not assumed in the theorem. The unconditional construction retains
level $2\mu$ and constant $71/10$.
The scaled low-critical corollary has both a smaller constant and a smaller
level.  The proof retained here is ordinary, with component-sensitive
consequences on their stated regions; its rational cutoffs are exact integer
inequalities.  The earlier perimeter samples concerned only a selected component.
They do not test the stronger condition that the level lie below every
critical value.
The sharp root-connector conclusion is not established for the unrestricted
class by the arguments in this record.
The adjacent classical literature on
lemniscate length concerns the arclength of the level curve $\{|p|=1\}$, which
is Erd\H{o}s~\#114.  Fryntov and Nazarov recall its history
\cite[Introduction]{fryntovnazarov2008}: Dolzhenko's bound $4\pi n$,
Pommerenke's bound $74n^2$ of 1961, and Borwein's bound $8\pi en$
\cite{borwein1995}.  Eremenko and Hayman proved $9.173n$
\cite[Theorem~1]{eremenkohayman1999}, Fryntov and Nazarov the asymptotically
sharp $2n+o(n)$ \cite{fryntovnazarov2008}, and Tao resolved the problem for
large $n$ \cite{tao2025}.  These results bound the length of a level curve, not the least internal
root-pair path.  This distinction identifies the quantity being estimated;
no claim of novelty follows from the scope or outcome of a literature search.

\section{Degree three}
\label{sec:degree-three}

\begin{theorem}[the cubic case]
\label{res:degree-three}
Let $f(z)=\prod_{j=1}^{3}(z-z_j)$ with $|z_j|<1$, the roots listed with
multiplicity.  Then two listed root occurrences are joined inside
$\{|f|<1\}$ by a polygonal path of length strictly below $2$.  If $f$ is
squarefree the two are distinct.
\end{theorem}
\leannote{Lean: \lproof{ErdosProblems/Erdos1041/PaperCubicCompletion.lean}{297}{cubic_paper_complete}.}

\begin{proof}
Multiple roots give a constant path, so assume $f$ squarefree.
The Erd\H{o}s--Herzog--Piranian component lemma
\cite[remark before Problem~5, p.~139]{ehp1958} gives a component of
$\{|f|<1\}$ containing two roots. Join them by a compact path in that
open component, and choose a regular value modulus $t<1$ larger than
$\max|f|$ on the path. The component of $\{|f|<t\}$ containing it has
$k\ge2$ roots and hence $k-1\ge1$ critical points, by the component-wise
Riemann--Hurwitz count in the proof of \cite[Proposition~2.1]{eks2010}.
One of these critical points has value of modulus below $t<1$. This
choice of $t$ avoids assuming that the level $|f|=1$ is regular.

Suppose first that $f'$ has two distinct zeros.  Choose a critical point $c$
minimising $|f(c)|$, write the other one as $c+\delta$, and put $v=f(c)$, so
$0<|v|<1$.  Monicity gives the exact expansion
$f(c+d)=d^3-\tfrac32\delta d^2+v$.  Choose $\alpha$ with $\alpha^3=v$ and set
$b=\delta/\alpha$; dividing by $v$ gives
$f(c+\alpha w)/v=P_b(w)$ with
\[
 P_b(w)=w^3-\tfrac32bw^2+1 .
\]
At the other critical point $f(c+\delta)=v(1-b^3/2)$, so minimality of $|v|$
is exactly the hypothesis $|1-b^3/2|\ge1$.

Under that hypothesis $P_b$ has at least two zeros, with multiplicity, in the
closed unit disc.  In the strict region a unit-circle zero $w$ would give
$b=\tfrac23(w+w^{-2})$. Write $w^3=e^{2iu}$; substitution gives
\[
 |1-b^3/2|^2
 =1+\frac{64}{729}\cos^4u\,(16\cos^2u-27)\le1.
\]
Since $16\cos^2u-27<0$, equality requires $\cos u=0$, which gives
$b=0$. Thus the strict region has no unit-circle zero.  The radial deformation $b\mapsto sb$, $s\ge1$, stays in
the strict region, since $z=b^3/2$ and $|1-z|^2>1$ give
$|1-s^3z|^2-1=s^3(s^3|z|^2-2\Re z)>0$.  For large $s$, Rouch\'e's theorem on
$|w|=1$ compares $P_{sb}$ with $-\tfrac32sbw^2$, whose modulus
$\tfrac32s|b|$ exceeds $2\ge|w^3+1|$, so $P_{sb}$ has exactly two zeros in
the open unit disc; no zero crosses the circle along the deformation.  For the equality case with $b\ne0$, the same radial deformation lies in
the strict region for every $s>1$, and continuity of the root multiset
allows $s\downarrow1$. If $b=0$, then $P_b(w)=w^3+1$ and all three
roots lie on the unit circle.

If $P_b(w)=0$ and $0\le t\le1$, then
$P_b(tw)=1-t^2-t^2(1-t)w^3$, so $|w|\le1$ gives
$|P_b(tw)|\le(1-t^2)+t^2(1-t)|w|^3\le1-t^3\le1$.  The whole segment from $0$
to $w$ therefore lies in the closed unit sublevel set.  Selecting two normalized roots $w_1,w_2$ with
$|w_i|\le1$, the two segments from $c$ to $c+\alpha w_i$ lie in
$\{|f|\le|v|\}\subset\{|f|<1\}$ and have combined length at most
$2|\alpha|=2|v|^{1/3}<2$.

If instead $f'$ has a double zero $c$, then $f(c+d)=d^3+v$. With $\alpha^3=-v$ and
$\omega=e^{2\pi i/3}$, its roots are
$c+\alpha$, $c+\alpha\omega$, $c+\alpha\omega^2$.  Averaging their squared
moduli gives $|c|^2+|\alpha|^2<1$, so $|\alpha|<1$; any two radial spokes
have total length $2|\alpha|<2$, and along either spoke the value has modulus
$|v|(1-t^3)<1$ away from the root endpoint.
\end{proof}

The unit level need not be regular, even when all roots lie in the open
unit disc. For example,
$Q(z)=(z-1)(z+1)(z+9/10)$ has $Q(1/3)=-148/135<-1$.
Its minimum $v$ on $[-9/10,1]$ occurs at an interior critical point and
satisfies $v<-1$. Set $r=|v|^{-1/3}<1$. The monic cubic
$r^3Q(z/r)$ has the three distinct roots $-r,-9r/10,r$ in the open unit
disc and a critical value equal to $-1$. Its unit level is therefore
singular, although the cubic path theorem applies.

\noindent\href{https://github.com/wcook04/plectis-erdos/blob/bb4e24651b37cd096d3749ef299c3317930b3a3b/ErdosProblems/Erdos1041/PaperCubicCompletion.lean#L297}{A formal-source implementation includes the polygonal path, strict variation bound and repeated-root case; its recorded build status is stated below.}

\paragraph{Sources and scope.}
\label{bdry:degree-three}
The source snapshot contains \href{https://github.com/wcook04/plectis-erdos/blob/3d6d938d696fed0fb71dd55115a18a73738ff223/lean/ErdosProblems/Erdos1041/PaperCubicCompletion.lean#L297}{the complete degree-three path implementation} and its supporting cubic source.
This earlier implementation lies outside the main repository's recorded
checked build. The supplied release snapshot also contains \href{https://github.com/wcook04/plectis-erdos-lean/blob/52f29ad173b04e3bac941b3663f2b9aebe5de0bb/ErdosProblems/Erdos1041/PaperCubicMonic.lean#L32}{the cubic theorem under monicity, degree three and the open-disc hypothesis}
and \href{https://github.com/wcook04/plectis-erdos-lean/blob/52f29ad173b04e3bac941b3663f2b9aebe5de0bb/Solutions/ExternalVerification1041CubicPath.lean#L38}{its expanded solution statement}. The release proof obtains a root enumeration from
monicity and degree; it does not assume an additional analytic path
supplier. It constructs a continuous polygonal curve with bounded
variation, values strictly below one and length strictly below two.
Its endpoints are distinct when the polynomial is squarefree; repeated
root occurrences are handled by the enumerated version.
The supplied index classifies these release declarations as
\texttt{release\_only}, not \texttt{ci\_checked}. The solution source
contains a proof rather than the placeholder in the corresponding
challenge file. This identifies the full formal target without claiming
a fresh compilation or axiom audit. The proof above remains an ordinary
mathematical argument.
Pendyala \cite[Theorem~1, pp.~1--3]{june2026} proves the separate quartic
case; no historical priority for the cubic theorem is asserted here.

\section{An all-degree critical-value separation theorem}
\label{sec:critical-value-separation}

The next theorem applies when one simple critical value is separated
from all the others by a disc in the value plane. After a change of variables,
the chosen critical point is $0$ and its value is $1$. The hypothesis below
then says that the disc of centre $w_0$ and radius $S$ contains no other
critical value. It is a condition to check, not an isolation property
guaranteed for every polynomial.

The simplicity assumption means that the chosen critical point is a simple
zero of the derivative, so two inverse branches meet there. A multiple
chosen critical point is excluded; the other critical points need not be
simple. Also, separation of roots in the $z$-plane
is not enough: different critical points can have nearly equal values.
An exact quartic example below makes this distinction explicit.
For a positive example, take $f(z)=z^3-3a^2z$ with $0<a<1/\sqrt3$.
Its roots $0,\pm\sqrt3a$ lie in the unit disc. At $c=a$ the critical
value is $v=-2a^3$, and the other normalised critical value is $-1$.
Thus the disc centred at $1$ with radius $4/3$ contains $0$ and $1$
but excludes the other critical value. Radius $2$ is allowed as well:
the other critical value then lies on the boundary. The proof must
therefore allow a nonregular outer level.

\begin{theorem}[separation of one simple critical value]
\label{res:critical-value-separation}
Let $P$ be a polynomial of degree $n\ge3$ whose leading coefficient has
modulus one, with
\[
 P(0)=1,\qquad P'(0)=0,\qquad P''(0)\ne0.
\]
Fix $w_0\in[0,1]$ and $S>\max(w_0,1-w_0)$.  Suppose every other critical
point $d\ne0$ satisfies
\[
 |P(d)-w_0|\ge S.                                      \tag{4}
\]
Put $p=w_0(1-w_0)$.  The two local solutions of
$P(Z(\xi))=1-\xi^2$, $Z(0)=0$, continue along the real segment to one
injective root-to-root connector $\Gamma$.  Its endpoints are distinct roots,
$\Gamma\subseteq\{|P|\le1\}$, and
\begin{equation*}\label{eq:disk-family-length}
 \operatorname{length}(\Gamma)^2
 \le 2\Bigl(\frac{S}{n-1}\Bigr)^{2/n}
 \log\!\frac{S^2+S+p}{S^2-S+p}.                       \tag{5}
\end{equation*}
Consequently the connector is shorter than $2$ whenever
\begin{equation*}\label{eq:disk-family-coefficient}
 \Bigl(\frac{S}{n-1}\Bigr)^{2/n}
 \log\!\frac{S^2+S+p}{S^2-S+p}<2.                    \tag{6}
\end{equation*}
\end{theorem}
\leannote{Lean: \lproof{ErdosProblems/Erdos1041/PaperCompleteR21/CriticalValueSeparationTransport.lean}{344}{discSep_separation_long}, \lproof{ErdosProblems/Erdos1041/PaperCompleteR21/CriticalValueSeparationTransport.lean}{305}{discSep_squared_length_le}, \lproof{ErdosProblems/Erdos1041/PaperCompleteR21/CriticalValueSeparationTransport.lean}{334}{discSep_length_lt_two}, \lproof{ErdosProblems/Erdos1041/PaperCompleteR21/CriticalValueSeparationTransport.lean}{179}{discSep_bergman_factor}, and 2 further declarations. Conditional on the connector and area construction that this proof produces; see the \hyperref[sec:coverage]{coverage section}.}

\begin{proof}
Let $Q=D(w_0,S)$ and let $U$ be the component of $P^{-1}(Q)$ containing
$0$.  Both $0$ and $1$ lie in $Q$, and (4) says that $0$ is the only critical
point in $U$.  A component of a polynomial sublevel set is simply connected: on any
Jordan curve in the component, the maximum principle bounds the same
polynomial throughout its interior, which must therefore remain in the
component. Condition (4) permits critical values on $\partial Q$, so
we must not assume that $\partial U$ is regular. For $1-w_0<S'<S$,
let $U'$ be the component of $\{|P-w_0|<S'\}$ containing $0$.
This level is regular, and $0$ contributes its only unit of ramification.
The component count in the proof of \cite[Proposition~2.1]{eks2010}
therefore gives $\deg(P|_{U'})=2$.

These nested components exhaust $U$. Indeed, a point of $U$ can be
joined to $0$ by a compact path in $U$; the maximum of $|P-w_0|$ on
that path is strictly below $S$, so the path eventually lies in $U'$.
For any fixed regular value in $Q$, all its finitely many preimages in
$U$ therefore lie in one such $U'$ once $S'$ is sufficiently large.
The proper map $P:U\to Q$ consequently also has degree two. This
establishes the count without a smoothness assumption at the outer level.

Set $a=1-w_0$.  On $U$, the function $1-P$ has only its double zero at $0$,
so it has a single-valued square root
\[
 \xi(z)^2=1-P(z).
\]
Taking this square root removes the double branching at $0$.
The map $\xi$ is proper into
\[
 \widetilde Q=\{\xi:|\xi^2-a|<S\}.
\]
Away from $0$, the identity $2\xi\xi'=-P'$ gives $\xi'\ne0$;
at $0$, $(\xi'(0))^2=-P''(0)/2\ne0$. The target is star-shaped about
$0$: because $S>a$, the convex disc $D(a,S)$ contains $0$, so
$t^2\xi^2\in D(a,S)$ whenever $\xi\in\widetilde Q$ and $0\le t\le1$.
It is therefore connected and simply connected. A proper local
biholomorphism has open and closed image, hence is a covering of this
target. Since $U$ is connected, it is a conformal bijection. Its inverse
$Z$ continues both local inverse branches through the critical point.  Since $[-1,1]\subset\widetilde Q$, the
curve $Z([-1,1])$ joins the two distinct points over $P=0$, and
$P(Z(\xi))=1-\xi^2\in[0,1]$ gives its containment.

We next map the parameter domain to the unit disc so that the curve
becomes a straight interval. This allows its length to be estimated by
the area of $U$, without replacing the domain by a larger disc.
Squaring sends $\widetilde Q$ onto $D(a,S)$. To map this disc to
$\mathbb D$ while fixing $0$, compose $w\mapsto(w-a)/S$ with the disc
automorphism taking $-a/S$ to $0$. The result is the M\"obius map
\[
 w\longmapsto\frac{Sw}{S^2+aw-a^2}.
\]
Taking its square root after substituting $w=\xi^2$ gives the conformal
bijection
\[
 \zeta(\xi)=\xi\sqrt{\frac{S}{S^2+a\xi^2-a^2}}
     :\widetilde Q\longrightarrow\mathbb D,
\]
where the square-root factor is chosen positive at $0$. It sends
$[-1,1]$ to $[-q,q]$, with
\[
 q^2=\frac{S}{S^2+p},\qquad
 S^2-S+p=(S-w_0)(S-(1-w_0))>0.
\]
The last inequality is exactly where $S>\max(w_0,1-w_0)$ ensures
$0<q<1$, as required by the length estimate below.
For $\Phi=Z\circ\zeta^{-1}:\mathbb D\to U$, the Bergman segment inequality,
proved by the kernel estimate of Section~\ref{sec:low-critical-closure} with
the double integral taken over $[-q,q]^2$, gives
\begin{align*}
 \operatorname{length}(\Gamma)^2
 &\le \frac2\pi\log\frac{1+q^2}{1-q^2}
       \operatorname{Area}(U)\\
 &=\frac2\pi\log\frac{S^2+S+p}{S^2-S+p}
       \operatorname{Area}(U).                       \tag{7}
\end{align*}

It remains to bound the area of this component. The factor $n-1$ comes
from the roots outside it, rather than from P\'olya's inequality alone.
The auxiliary capacity estimate is as follows: for a monic degree-$n$
polynomial $F$, a regular value $t>0$ of $|F|$, and a component $V$ of
$\{|F|<t\}$ containing $k<n$ zeros counted with multiplicity,
\[
 \frac{\operatorname{cap}(\overline V)^n}{t}<\frac{k}{2n-k}.
\]
Multiplying $F$ by a unimodular constant leaves the estimate unchanged.
This is a separate component estimate, not part of the
Riemann--Hurwitz statement cited above. We give the harmonic-measure
argument behind it, following
\href{https://github.com/wcook04/plectis-erdos/blob/8efbccc235df64a38d83f5dc7b1949e2ad18270d/research_corpus/Erdos1041/problem/ExteriorBlaschkeFibreCapacityGap.md}{Theorem~2 and Corollary~3 of the component-capacity proof}.

Let $m=n-k$ and let $\psi$ map $|\zeta|>1$ conformally onto the
exterior of $\overline V$, with positive leading coefficient
$\operatorname{cap}(\overline V)$. The $m$ excluded roots have preimages
$\xi_j$ with $|\xi_j|>1$. Put $b_j=1/\overline{\xi_j}$ and form
\[
 B(z)=\prod_{j=1}^m\frac{z-b_j}{1-\overline{b_j}z}.
\]
Reflection in the unit circle shows that $F(\psi(\zeta))/t$ equals
$\zeta^n\overline{B(1/\overline\zeta)}$ up to a constant of modulus one:
their zeros and pole at infinity agree, and both have modulus one on
the boundary. Comparing leading coefficients gives
\[
 |B(0)|=\frac{\operatorname{cap}(\overline V)^n}{t}.
\]
Regularity of the level permits analytic continuation of $\psi$ across
the boundary. On the unit circle,
\[
 \frac{d}{d\theta}\arg F(\psi(e^{i\theta}))
 =n-|B'(e^{i\theta})|>0.
\]
Positivity follows because a regular polynomial level is mapped locally
in the positive boundary direction. Thus $|B'|<n$ there.

The boundary-fibre identity used next is Lemma~1 of the same
component-capacity proof. Here is its harmonic-measure argument.
For a continuous real function $h$ on the unit circle, let $H$ be its
harmonic extension to the disc. Since $H\circ B$ is harmonic, the mean
value property and the Poisson formula give
\[
 \frac1{2\pi}\int_0^{2\pi}h(B(e^{i\theta}))\,d\theta
 =H(B(0))
 =\frac1{2\pi}\int_0^{2\pi}h(e^{i\phi})
   \frac{1-|B(0)|^2}{|e^{i\phi}-B(0)|^2}\,d\phi.
\]
Differentiating the factors gives the positive angular derivative
\[
 \frac{d}{d\theta}\arg B(e^{i\theta})
 =|B'(e^{i\theta})|
 =\sum_{j=1}^m\frac{1-|b_j|^2}{|e^{i\theta}-b_j|^2}>0.
\]
Changing variables through the $m$ boundary inverse branches in the
first integral and comparing the continuous densities therefore yields
\[
 \sum_{B(\zeta)=w}\frac1{|B'(\zeta)|}
 =\frac{1-|B(0)|^2}{|w-B(0)|^2}\qquad(|w|=1).
\]
In particular each boundary fibre consists of $m$ distinct points. Choosing
$w=-B(0)/|B(0)|$ yields
\[
 \frac mn<\frac{1-|B(0)|}{1+|B(0)|},
 \qquad |B(0)|<\frac{n-m}{n+m}=\frac{k}{2n-k}.
\]
Here $B(0)\ne0$ because every excluded-root preimage is finite. This proves
the component estimate and explains how the excluded roots enter it.

For the regular inner components $U'$ above, the degree-two count gives
exactly two zeros of $P-w_0$, counted with multiplicity. Since $2<n$,
the preceding estimate applies with $k=2$ and gives
\[
 \frac{\operatorname{cap}(\overline{U'})^n}{S'}
 <\frac{k}{2n-k}=\frac1{n-1}.
\]
P\'olya's area--capacity inequality
$\operatorname{Area}(K)\le\pi\operatorname{cap}(K)^2$ \cite[Theorem~6]{crane}
then yields
$\operatorname{Area}(U')<\pi(S'/(n-1))^{2/n}$. The exhaustion proved
at the start and continuity of area from below give, without assuming
regularity of the limiting boundary,
\[
 \operatorname{Area}(U)\le\pi\Bigl(\frac{S}{n-1}\Bigr)^{2/n}.
\]
Substitution into (7) proves (5), and (6) makes its right side strictly less
than $4$.
\end{proof}

\begin{corollary}[uniform radius $4/3$]
\label{res:critical-value-thresholds}
Inequality~\eqref{eq:disk-family-coefficient} holds for every $n\ge3$,
every $w_0\in[0,1]$, and every $4/3\le S\le2$.  Thus, if $f$ is monic with
roots in the open unit disc, $c$ is a simple critical point with
$v=f(c)\ne0$ and $|v|<1$, and
\[
 \left|\frac{f(d)}v-w_0\right|\ge\frac43
\]
for every other critical point $d$, then two roots of $f$ are joined inside
$\{|f|<1\}$ by a curve of length strictly below $2$.
\end{corollary}
\leannote{Lean: \lproof{ErdosProblems/Erdos1041/PaperCompleteR21/CriticalValueSeparationTransport.lean}{645}{discSepCoefficient_lt_two}, \lproof{ErdosProblems/Erdos1041/PaperCompleteR21/CriticalValueSeparationTransport.lean}{702}{discSep_uniform_radius}, \lproof{ErdosProblems/Erdos1041/PaperCompleteR21/CriticalValueSeparationTransport.lean}{629}{discSep_ratio_le_branch}, \lproof{ErdosProblems/Erdos1041/PaperCompleteR21/CriticalValueSeparationTransport.lean}{637}{discSep_branch_le_seven}. Conditional on the connector and area construction in the proof of Theorem~\ref{res:critical-value-separation}; see the \hyperref[sec:coverage]{coverage section}.}

\begin{proof}
For $S\ge4/3$ and $p\ge0$,
\[
 \frac{S^2+S+p}{S^2-S+p}\le\frac{S+1}{S-1}\le7.
\]
Also $S/(n-1)\le1$ when $S\le2$ and $n\ge3$, while $\log7<2$.
This proves~\eqref{eq:disk-family-coefficient}.  For the unnormalised
polynomial, set $z=c+|v|^{1/n}w$. Then $P(w)=f(z)/v$ has leading
coefficient $|v|/v$, of modulus one. Scaling the connector multiplies its
length by $|v|^{1/n}<1$ and takes $\{|P|\le1\}$ to
$\{|f|\le|v|\}\subset\{|f|<1\}$. The selected value $v$ need not attain
the minimum defining $\mu$.
\end{proof}

Taking $w_0=1$ and $S=2$ gives the earlier criterion
$|1-f(d)/v|\ge2$ in every degree $n\ge3$. If $c$ attains the minimum
modulus among \emph{all} critical values and $v\ne0$, the
Fekete--resultant bound gives $|v|\le R^n<1$ when all roots lie in a disc
of radius $R<1$. Deleting zero critical values before taking the minimum
would invalidate that inference. For example, $(z-R)^2(z+R)$ with $0<R<1$ has critical
values $0$ and $32R^3/27$, so its least nonzero critical-value modulus
exceeds $R^3$. The repeated root already gives a constant path, but it
does not give this nonzero value the required bound. For a critical point
chosen in another way, $|v|<1$ must remain an explicit hypothesis.
In degree three the choice $w_0=1$ already works at
$S=6/5$, since
\[
 \Bigl(\frac35\Bigr)^{2/3}\log 11<2.
\]
This sharper cubic constant is part of the numerical Lean kernel cited below.

The separation hypothesis of Theorem~\ref{res:critical-value-separation}
is a condition on values, not on the positions of roots or critical points.
To distinguish these conditions, consider the family
\[
 g_\varepsilon(z)=z^4-\varepsilon z^3-2z^2+3\varepsilon z+\tfrac12,
 \qquad 0<\varepsilon\le\tfrac1{32}.
\]
Its derivative and two critical values are
\[
 g_\varepsilon'(z)=(z^2-1)(4z-3\varepsilon),\qquad
 g_\varepsilon(\pm1)=-\tfrac12\pm2\varepsilon.
\]
At the endpoints of each of the intervals
$(-3/2,-5/4)$, $(-3/4,-1/2)$, $(1/2,3/4)$ and $(5/4,3/2)$,
the values of $g_\varepsilon$ have opposite signs throughout this parameter
range. Each interval therefore contains exactly one root. Together with
the critical points $-1,3\varepsilon/4,1$, these give seven points at
pairwise distances greater than $1/4$, uniformly in $\varepsilon$.
Nevertheless the two displayed critical values differ by $4\varepsilon\to0$.
The earlier numerical choice is the exact member $\varepsilon=1/262144$:
\[
 g(z)=z^4-\frac{z^3}{262144}-2z^2+\frac{3z}{262144}+\frac12,
 \qquad
 g'(z)=4(z+1)(z-3/1048576)(z-1),
\]
for which that difference is $1/65536$.
The monic polynomial $2^{-4}g_\varepsilon(2z)$ puts every root inside
$|z|<3/4$ and retains a spatial separation greater than $1/8$.
For the critical points $\pm1/2$ of the scaled polynomial, the ratio
of the two values is unchanged and obeys
\[
 \frac{g_\varepsilon(-1)}{g_\varepsilon(1)}
 =\frac{1+4\varepsilon}{1-4\varepsilon}\le\frac97<\frac43.
\]
The reciprocal ratio lies between $0$ and $1$. Hence, for any centre
$w_0\in[0,1]$, neither selected value satisfies the theorem's
radius-$4/3$ isolation condition: the other normalised value is less
than $4/3$ away. This does not exclude a suitable isolated value at a
\emph{different} critical point.

\paragraph{Sources and scope.}
\label{bdry:critical-value-separation}
The two-sheeted component, its square-root uniformisation, the Bergman segment
inequality, the component-capacity estimate, P\'olya's area--capacity
inequality and the exhaustion are ordinary arguments. They are collected in
\href{https://github.com/wcook04/plectis-erdos/blob/8efbccc235df64a38d83f5dc7b1949e2ad18270d/research_corpus/Erdos1041/problem/DiskFamilyCriticalValueSeparation.md}{the proof using an isolated critical value}, whose capacity step is cited
separately above.
The separate
\href{https://github.com/wcook04/plectis-erdos/blob/f36e325c8b05fe5c3ce8a2ca32a699f6dc6336b9/ErdosProblems/Erdos1041/DiskFamilyCriticalValueSeparation.lean}{Lean numerical kernel}
checks the coefficient bound for $n\ge3$, $4/3\le S\le2$ and $p\ge0$, and
the implication from the squared length bound to length below $2$.  It does
not formalise the analytic hypotheses producing~\eqref{eq:disk-family-length}.
The adjacent axiom-audit source names those numerical declarations; it is not
an axiom audit of the ordinary analytic theorem.

The proof requires both a simple critical point and a disc centred on
$[0,1]$ that isolates its normalised value. For example, the trinomial
$z^n-b$ with $n>2$ and $0<|b|<1$ has a short path by
Theorem~\ref{res:trinomial-all-degree}, but its only critical point is not
simple, so this method does not apply. The restriction is on the chosen
critical point and on the distance of the other critical values from the
chosen disc. The other critical points may be multiple, and their values
need not be separated from one another. No sharpness is claimed for
$S=4/3$, and no assertion is made that every polynomial has a suitable
isolated simple value. In particular, this is not a proof of the
unrestricted assertion of Erd\H{o}s~\#1041.

A related area comparison is Dubinin's Theorem~1. The
source is \cite{dubinin}, Theorem~1 on printed page~85 of the POMI original.
It concerns a holomorphic function that gives a full $n$-fold covering of an
annulus $t_1<|w|<t_2$, and, in the notation of that theorem, with $E$ its
explicitly defined complementary set, it states
\[
 \Bigl(\frac{t_2}{t_1}\Bigr)^{2/n}\le\frac{m(E\cup D)}{m(E)}.
\]
The full covering and the explicit complementary set are hypotheses of the
theorem.  Tao cited it on the Erd\H{o}s Problem~\#1041 discussion page on
25 March 2026 for the relative scaling factor $s^{2/n}$ between the areas of
two nested sublevel sets, which is the reading of the displayed inequality in
which $E$ and $E\cup D$ are those two sublevel sets and the covering
hypothesis holds.  No step of this note uses that relative inequality.  The
disk-family theorem instead uses P\'olya's absolute area--capacity inequality
on a component containing fewer than $n$ roots, after the
component-capacity estimate supplies its strict
capacity gap.  The absolute sublevel inequality
$\operatorname{Area}\{|P|<T\}\le\pi T^{2/n}$ remains the area input to
Theorem~\ref{res:low-critical-thirteen-twentyfifths} and to
Theorem~\ref{res:constant-factor-path}.  Dubinin's theorem is a neighbouring
result under a full covering hypothesis and supplies no root connector.

\section{Collinear roots and two sparse polynomial families}
\label{sec:solved-polynomial-families}

The following results impose concrete restrictions: roots on one line,
a quintic with two missing coefficients, or a polynomial of the form
$P((z-h)^q)$ with $P$ cubic. None covers a general polynomial with roots
in the unit disc. We first state the finite inequalities used in their
proofs, then explain how the inequalities produce contained paths. The
formal links distinguish these two steps.

\subsection{Collinear roots and Chebyshev comparison}
\label{subsec:sharp-collinear-chebyshev}

The extreme zeros of $T_n$ are $\pm\cos(\pi/(2n))$. Scaling them to
$\pm1$ makes the comparison polynomial have the same endpoint zeros as
the polynomial under study; dividing by its leading coefficient makes
both polynomials monic. Their difference then has smaller degree.
Accordingly, put
\[
 r_n=\cos\frac{\pi}{2n},
 \qquad C_n=\frac{1}{2^{n-1}r_n^n}.
\]

\begin{theorem}[Chebyshev comparison]
\label{prop:sharp-collinear-chebyshev-comparator}
Let $m\ge0$, let $p\in\mathbb R[X]$ be monic of degree $m+2$, and let
\[
 -1<c_0<\cdots<c_m<1,\qquad |c_i|\le1.
\]
Suppose $p(-1)=p(1)=0$ and
$p(c_i)p(c_{i+1})<0$ for $0\le i<m$.  Then
\[
 \min_{0\le i\le m}|p(c_i)|\le C_{m+2}.
\]
\end{theorem}
\leannote{Lean: \lproof{ErdosProblems/Erdos1041/SharpCollinearChebyshev.lean}{133}{exists_peak_le_comparisonBound}.}

This is exactly the statement selected as
\href{\repobase/ExternalVerification1041SolvedFamilies/Solution.lean\#L74}{the formal Chebyshev endpoint}.
It proves the finite alternation inequality; it does not yet construct
the segment between two roots.

All real-rooted polynomials, and their images under rotations and
translations, satisfy the collinearity hypothesis. Three noncollinear
roots already fall outside it. The gain is an explicit sharp bound for
the modulus along one adjacent-root segment, not a bound obtained by
assuming that segment is contained.

\begin{theorem}[a sharp bound for collinear roots]
\label{thm:sharp-collinear-diameter}
Let $f$ be a monic polynomial of degree $n\ge2$ whose zero occurrences are
collinear, and let $D$ be their diameter.  Some two adjacent zero occurrences
are joined by a segment of length at most $D$ on which
\[
 |f(z)|\le
 \frac{(D/2)^n}{2^{n-1}\cos^n(\pi/(2n))}.          \tag{9}
\]
The constant in \emph{(9)} is best possible in every degree.  Equality is
attained by affine images of the zeros of $T_n$ whose extreme zeros have
distance $D$.
\end{theorem}
\leannote{Lean: \lproof{ErdosProblems/Erdos1041/PaperCompleteR21/CollinearDiameterWhole.lean}{367}{sharp_collinear_root_diameter}, \lproof{ErdosProblems/Erdos1041/PaperCompleteR21/CollinearDiameterWhole.lean}{967}{sharp_collinear_root_diameter_monic}, \lproof{ErdosProblems/Erdos1041/PaperCompleteR21/CollinearDiameterWhole.lean}{928}{exists_collinear_factorisation}, \lproof{ErdosProblems/Erdos1041/PaperCompleteR21/CollinearDiameterWhole.lean}{212}{exists_gap_le_comparisonBound}, and 5 further declarations in the \hyperref[sec:coverage]{coverage section}.}

\begin{proof}
A repeated zero gives the constant path, so assume the zero values are
distinct. Let $m$ be the midpoint of the two extreme zeros, and choose
$\theta\in\mathbb R$ so that their line is $m+e^{i\theta}\mathbb R$.
Put $R=D/2$. Translation by $m$, rotation by $e^{i\theta}$ and scaling
by $R$ give the polynomial
\[
 q(w)=R^{-n}e^{-in\theta}f(m+Re^{i\theta}w),
\]
which is monic with real zeros
$-1=y_1<\cdots<y_n=1$, and
$|f(m+Re^{i\theta}w)|=R^n|q(w)|$.

Compare $q$ with the monic endpoint-normalized Chebyshev polynomial
\[
 q_*(x)=\frac{T_n(r_nx)}{2^{n-1}r_n^n}.
\]
Both $q$ and $q_*$ vanish at $\pm1$, and $|q_*|\le C_n$ on
$[-1,1]$.  For each gap $[y_i,y_{i+1}]$, choose $c_i$ at which
$|q|$ is maximal. The maximum lies in the interior: $q$ vanishes at the
endpoints and nowhere inside the gap. Since every root is simple, the
signs at the $n-1$ chosen points alternate. If every gap maximum were
larger than $C_n$, then $q-q_*$ would have the same nonzero signs there.
The intermediate value theorem would give $n-2$ zeros between consecutive
$c_i$, in addition to the two zeros at $\pm1$. These $n$ zeros are distinct,
but $q-q_*$ is nonzero and has degree at most $n-1$ because its leading
terms cancel.
This contradiction selects a gap on which $|q|\le C_n$.  Scaling back proves
\emph{(9)}.

For sharpness, take
\[
 y_k=\frac{\cos((2k-1)\pi/(2n))}{r_n},
 \qquad 1\le k\le n.
\]
These are the zeros of $q_*$, have extremes $\pm1$, and every adjacent gap
contains a scaled Chebyshev extremum where $|q_*|=C_n$.  No smaller universal
constant can work.
\end{proof}

\begin{corollary}[collinear Erd\H{o}s case]
\label{cor:collinear-erdos-1041}
If the zero occurrences of a monic polynomial of degree $n\ge2$ lie on one line
in the open unit disc, two of them are joined by a curve of length strictly
below $2$ inside $\{|f|<1\}$.
\end{corollary}
\leannote{Lean: \lproof{ErdosProblems/Erdos1041/PaperCompleteR21/CollinearDiameterWhole.lean}{515}{collinear_erdos_1041}, \lproof{ErdosProblems/Erdos1041/PaperCompleteR21/CollinearDiameterWhole.lean}{984}{collinear_erdos_1041_monic}, \lproof{ErdosProblems/Erdos1041/PaperCompleteR21/CollinearDiameterWhole.lean}{928}{exists_collinear_factorisation}, \lproof{ErdosProblems/Erdos1041/PaperCompleteR21/CollinearDiameterWhole.lean}{367}{sharp_collinear_root_diameter}.}

\begin{proof}
A repeated zero gives the constant path, so assume the zeros are distinct.
Their diameter satisfies $0<D<2$. Since
$\cos(\pi/(2n))\ge1/\sqrt2$, the defining formula gives
$C_n\le2^{1-n/2}\le1$, with equality possible only at $n=2$.
Thus \emph{(9)} is strictly below one, and the selected segment has
length at most $D<2$.
\end{proof}

Corollary~\ref{cor:collinear-erdos-1041} already follows from a theorem of
Erd\H{o}s, Herzog and Piranian \cite[Theorem~1, p.~126]{ehp1958}: if the zeros
of a monic polynomial of degree $n$ are real, lie in $[-1,1]$ and have
centroid in $[0,1]$, then $\{|f|<1\}\cap\mathbb R$ contains an interval
holding at least $n/2$ of the zeros.  Apply it to the normalized polynomial
$q$, replacing it by the monic polynomial $(-1)^nq(-w)$ if its centroid
is negative. This reverses the root line without changing modulus bounds. For
$n\ge3$ two consecutive zeros in that interval are joined by a segment of
length at most $D$ inside $\{|f|<(D/2)^n\}$, and for $n=2$ the segment between
the zeros lies in $\{|f|\le(D/2)^2\}$.
Theorem~\ref{thm:sharp-collinear-diameter} adds the sharp level and its
equality configurations.

\paragraph{Critical sequences and the extremal configuration.}
\hypertarget{collinear-source-comparison}{}
For distinct real zeros, each gap between consecutive zeros contains exactly
one zero of $q'$, and $|q|$ attains its maximum over the gap only there.  The
points $c_i$ chosen in the proof of Theorem~\ref{thm:sharp-collinear-diameter}
are therefore the critical points of $q$, and the gap maxima are the moduli of
its ordered critical sequence $(q(c_1),\ldots,q(c_{n-1}))$, with
$c_1<\cdots<c_{n-1}$. In this distinct-root setting, the ordered sequence
determines the real polynomial up to an increasing real affine change
of variable \cite[Theorem~1]{eremenkoyuditskii2012}. The proof above uses
only the alternation count, not this classification theorem.  In this language, \emph{(9)} says that a
real polynomial of degree $n$ with $n$ distinct real zeros, leading
coefficient $a$ and zero diameter $D$ has a critical value of modulus at most
$C_n|a|(D/2)^n$.

The comparison polynomial $q_*$ is the monic polynomial of least deviation
from zero on $[-1/r_n,1/r_n]$, characterised by the equioscillation conditions
recalled in \cite[\S1]{eremenkoyuditskii2012}.  All $n-1$ of its critical
values have modulus $C_n$, so every critical point of $q_*$ lies on the level
curve $\{|q_*|=C_n\}$.  Extremals with this property also occur in the
level-curve length problem, where some extremal polynomial has all its
critical points on $\{|p|=1\}$ \cite[Lemma~5]{eremenkohayman1999}, and in the
sharp bound for $|f'|$ on a connected sublevel set $\{|f|\le1\}$ of a monic
polynomial, proved by Eremenko and Lempert
\cite[Theorem~1]{eremenkolempert1994}, whose
equality cases are $e^{-in\theta}T_n(2^{1/n-1}e^{i\theta}z+b)$ with $\theta$
real \cite[Theorem~A]{eremenko2007}.

\paragraph{Formal scope.}
\label{rem:sharp-collinear-formal-boundary}
Lean checks sign preservation, the alternating root-count mechanism, the
scaled Chebyshev endpoint and uniform bound, and
Theorem~\ref{prop:sharp-collinear-chebyshev-comparator}.  It does not check
the rigid normalization, the existence of the gap maxima, transport back to
the root line, or the equality-node locations used in
Theorem~\ref{thm:sharp-collinear-diameter}.  These are ordinary steps.  The
sharpness assertion concerns the maximum of $|f|$ on the selected adjacent-root
segment, and the displayed configurations attain it; no claim is made that
they exhaust the equality cases, and no claim is made that the selected
segment is a shortest path in $\{|f|<1\}$.

\subsection{Quintics with two missing coefficients}
\label{subsec:primitive-sparse-quintic}

The next inequality selects two indices from five real pairs. Its three
moment identities will come from the missing cubic and quadratic terms
of $z^5+az^4+bz+c$. For $0<r<2$ and
$0\le s_i\le1$, $x_i^2\le s_i$, put
\[
 E_i=s_i^4(s_i+r^2+2rx_i).
\]

\begin{theorem}[a consequence of three moment identities]
\label{prop:primitive-quintic-two-tail-energy-selector}
Suppose
\[
 \sum_{i=0}^4x_i=-r,\qquad
 \sum_{i=0}^4(2x_i^2-s_i)=r^2,\qquad
 \sum_{i=0}^4(4x_i^3-3s_ix_i)=-r^3.              \tag{10}
\]
Then at least one of the ten pairs $0\le i<j\le4$ satisfies
$E_i<1$ and $E_j<1$.
\end{theorem}
\leannote{Lean: \lproof{ErdosProblems/Erdos1041/PrimitiveQuinticInteriorTail.lean}{272}{primitiveInterior_exists_two_tailEnergy_lt_one}.}

This is the exact finite conclusion selected as
\href{\repobase/ExternalVerification1041SolvedFamilies/Solution.lean\#L24}{the formal two-index inequality}.
It selects two distinct indices.  It does not assert that the corresponding
complex root values are distinct.

The polynomial must have no $z^3$ or $z^2$ term in the coordinates used
here. For example, $z^5+bz+c$ is included, but a general quintic with all
coefficients nonzero is not. The root-location assumption alone does not
imply the three moment identities. Those identities are precisely what
lets the finite inequality find two radial segments.

\begin{theorem}[a quintic with two missing coefficients]
\label{thm:primitive-quintic-two-tail}
Let
\[
 p(z)=z^5+az^4+bz+c
\]
and suppose its five zero occurrences $w_0,\ldots,w_4$ lie in the closed unit
disc.  At least two distinct indices satisfy
\[
 |bw_i+c|\le1.                                    \tag{11}
\]
If $a\ne0$, two indices can be chosen with strict inequalities.  If $a=0$,
every index satisfies \emph{(11)}, and equality holds exactly when
$|w_i|=1$.

For open-disc zeros, two zero occurrences are joined inside $\{|p|<1\}$ by a
curve of length below $2$: use the two radial spokes through $0$ when their
values are distinct, and the constant path when the selected occurrences
have the same value.
\end{theorem}
\leannote{Lean: \lproof{ErdosProblems/Erdos1041/PaperCompleteR21/PrimitiveQuinticClosedDisc.lean}{260}{primitive_quintic_two_tail}, \lproof{ErdosProblems/Erdos1041/PaperCompleteR21/PrimitiveQuinticClosedDisc.lean}{313}{primitive_quintic_two_tail_of_polynomial}, \lproof{ErdosProblems/Erdos1041/PaperCompleteR21/PrimitiveQuinticClosedDisc.lean}{93}{two_tails_closedDisc_of_ne_zero}, \lproof{ErdosProblems/Erdos1041/PaperCompleteR21/PrimitiveQuinticClosedDisc.lean}{199}{tail_le_one_and_eq_iff_of_leading_zero}, and 1 further declaration in the \hyperref[sec:coverage]{coverage section}.}

\begin{proof}
If $a=0$, the root equation gives $|bw_i+c|=|w_i|^5$ for every
index, proving \emph{(11)} and its equality clause. For $a\ne0$, write
$a=re^{i\phi}$ with $r>0$ and set $z_i=e^{-i\phi}w_i$.
The missing $z^3$ and $z^2$ coefficients and Newton's identities give
\[
 \sum z_i=-r,\qquad \sum z_i^2=r^2,\qquad
 \sum z_i^3=-r^3.                                 \tag{12}
\]
At a zero,
\[
 |bw_i+c|=|w_i|^4|w_i+a|.
\]
Thus, for $x_i=\Re z_i$ and $s_i=|z_i|^2$, the squared modulus $|bw_i+c|^2$ is exactly
$E_i$ and \emph{(12)} becomes \emph{(10)}. The third moment gives
$r^3=|\sum_i z_i^3|\le5<8$, so $0<r<2$ as required by the finite
inequality.

We need to rule out four indices with squared modulus at least $1$.
The three known power sums determine the sum of any real harmonic
polynomial of degree at most three over the roots. On the unit circle,
$s^4(s+r^2+2rx)\ge1$ reduces to $x\ge-r/2$. We therefore use the harmonic
extension of $(-r/2-x)(1-x)^2$, which is nonpositive on that part of the
circle. In the variables $x=\Re z$ and $s=|z|^2$, this extension is
\[
 H_r(x,s)=(-r/2-x)(1-x)^2
 +(1-s)\left(1-\frac r4-\frac{3x}{4}\right).
\]
The useful feature of this choice is the expansion
\[
 H_r(\Re z,|z|^2)
 =1-\frac{3r}{4}+
 \Re\!\left[\left(r-\frac74\right)z+
             \left(1-\frac r4\right)z^2-\frac{z^3}{4}\right].
\]
It is harmonic, and summing this expression with the three moments in
\emph{(12)} gives
\[
 \sum_{i=0}^4H_r(x_i,s_i)=5-2r.                  \tag{13}
\]
On the unit circle its value is $(-r/2-x)(1-x)^2$, and
\[
 4-2r-(-r/2-x)(1-x)^2
 =(x+1)\left(\frac{(2-r)(3-x)}2+(1-x)^2\right)\ge0.
\]
The maximum principle therefore gives $H_r\le4-2r$ throughout the disc.
For an index with $s^4(s+r^2+2rx)\ge1$, one has $s>0$ and
$s^{-4}\ge1+4(1-s)$ by convexity on $(0,1]$. Consequently
\[
 2r(x+r/2)\ge s^{-4}-s\ge5(1-s).
\]
Put $d=x+r/2$, $u=1-x$, $A=1-r/4-3x/4$, and
$P=-u^2+(2r/5)A$.  The assumption $s^4(s+r^2+2rx)\ge1$ gives $0\le d<2$.  Since
$H_r=-du^2+(1-s)A$, it is nonpositive when $A\le0$; when $A>0$,
the preceding bound gives $H_r\le dP$.  The remaining estimate is the exact
sum-of-squares identity
\[
 \frac1{31}-P
 =\frac25\left(d-\frac{38}{31}
               +\frac14(u-\frac3{31})\right)^2
  +\frac{31}{40}\left(u-\frac3{31}\right)^2\ge0,
\]
which gives $P\le1/31$, and hence $H_r(x,s)\le2/31$ whether $P$ is
positive or nonpositive.  Thus $H_r\le2/31$ whenever the squared modulus is at least $1$.
If four indices had squared modulus at least $1$, \emph{(13)} would give
\[
 5-2r\le4-2r+\frac8{31}<5-2r,
\]
a contradiction. This proves the two-index conclusion when $a\ne0$.

For either value of $a$, at a zero $w$,
\[
 p(tw)=(1-t)c+(t-t^4)(bw+c)-t^4(1-t)w^5.
\]
For $0\le t\le1$ the three nonnegative weights sum to $1-t^5$.
The tail bound, $|w|\le1$, and $|c|\le1$ therefore keep each selected
radial segment in $\{|p|\le1\}$. Under the open-disc hypothesis,
$|c|<1$. Its coefficient $1-t$ is positive for $0\le t<1$, so
\[
 |p(tw)|<(1-t)+(t-t^4)+t^4(1-t)=1-t^5\le1.
\]
At $t=1$ the value is zero. Thus both segments lie in $\{|p|<1\}$,
and their total length is $|w_i|+|w_j|<2$.
\end{proof}

\paragraph{Formal scope.}
\label{rem:primitive-quintic-formal-boundary}
The moment identities and the inequality selecting two indices have
formal proofs.
The source snapshot additionally contains \href{https://github.com/wcook04/plectis-erdos/blob/3d6d938d696fed0fb71dd55115a18a73738ff223/lean/ErdosProblems/Erdos1041/PaperPrimitiveCompletionR10.lean#L235}{a complete implementation of the open-disc sparse-quintic path conclusion}, including the geometric assembly.
The formal-source index places the complete \emph{open-disc path conclusion}
in its checked build. That declaration does not by itself assert every
closed-disc equality clause of the longer theorem above.
That is supplied build evidence, not a new compilation performed for this
prose revision. The ordinary argument above is retained, and the source
link remains at its original commit.

\subsection{Polynomials obtained from a cubic by a power substitution}
\label{subsec:translated-cubic-quotient-fibres}

For $f(z)=P((z-h)^q)$, the roots associated with a nonzero root of $P$
form a regular $q$-gon centred at $h$. This rotational symmetry is a
substantial restriction: a generic polynomial of degree $3q$ does not
have it. It lets one contained segment for the cubic give two segments
in the original variable. The first theorem supplies that segment.

\begin{theorem}[a contained radial segment for a cubic]
\label{lem:cubic-safe-root-spoke}
If $r,s,v\in\mathbb C$ have modulus below one, at least one
$u\in\{r,s,v\}$ satisfies
\[
 \left|(tu-r)(tu-s)(tu-v)\right|\le1
 \qquad(0\le t\le1).
\]
\end{theorem}
\leannote{Lean: \lproof{ErdosProblems/Erdos1041/CubicQuotientFiberCase.lean}{161}{cubic_has_safe_root_spoke}.}

This is exactly
\href{\repobase/ExternalVerification1041SolvedFamilies/Solution.lean\#L14}{the formal cubic segment inequality}.

\begin{theorem}[a cubic composed with a power map]
\label{thm:translated-cubic-quotient-fibres}
Let $q\ge2$, $h\in\mathbb C$, $P$ be monic cubic, and
\[
 f(z)=P((z-h)^q).
\]
If every zero of $f$ lies in the open unit disc and $f$ has at least two
distinct zero values, then two zeros are joined through $h$ by a two-segment
path of length below $2$ inside $\{|f|<1\}$.  Equivalently this closes the
coefficient family
\[
 (z-h)^{3q}+A(z-h)^{2q}+B(z-h)^q+C
\]
in every degree $3q\ge6$.
\end{theorem}
\leannote{Lean: \lproof{ErdosProblems/Erdos1041/PaperCubicFibres.lean}{240}{complete_translated_cubic_quotient_fibres}.}

\begin{proof}
Fix a $q$th root $y$ of a quotient root and a primitive $q$th root of unity
$\zeta$.  The full fibre consists of $h+y\zeta^k$.  Since every fibre point
lies in the open unit disc,
\[
 \frac1q\sum_{k=0}^{q-1}|h+y\zeta^k|^2
   =|h|^2+|y|^2<1.                                \tag{14}
\]
Thus $|y|<1$, and the corresponding quotient root has modulus
$|y|^q<1$.  This verifies the open-disc hypothesis for all quotient roots
before selecting a radial segment for the cubic.
The identity retains more information than these separate bounds:
$|h|<1$ and every fibre radius is strictly below $\sqrt{1-|h|^2}$.
Consequently, any two contained fibre segments constructed below have total length
strictly less than $2\sqrt{1-|h|^2}$.  The position of the symmetry centre
therefore improves the length bound; replacing it by $2$ discards this information.

Write $P(w)=(w-r)(w-s)(w-v)$ and associate to $r$ the real number
$A_r=\Re(r\overline{s+v})$, cyclically.  The exact sum is
\[
 A_r+A_s+A_v=|r+s+v|^2-(|r|^2+|s|^2+|v|^2)>-3,
\]
so one of these numbers, say $A_r$, exceeds $-1$. For $0\le t\le1$,
\[
 \frac{|tr-s|^2+|tr-v|^2}{2}
 =t^2|r|^2+\frac{|s|^2+|v|^2}{2}-tA_r<1+t+t^2.
\]
AM--GM now gives $|tr-s|\,|tr-v|<1+t+t^2$.
Consequently, for $0\le t<1$,
\[
 |P(tr)|<(1-t)(1+t+t^2)=1-t^3\le1.              \tag{15}
\]
At $t=1$ the polynomial vanishes.  The entire spoke is therefore strictly
contained, including its origin endpoint.

For a nonzero quotient root satisfying the segment bound from \emph{(15)}, two distinct
fibre points satisfy
\[
 f(h+ty\zeta^k)=P(t^qr),
\]
and their two spokes through $h$ have total length
$2|y|<2\sqrt{1-|h|^2}\le2$.  If the selected
quotient root is zero, choose a nonzero quotient root $s$ and write
$P(w)=w(w-s)(w-v)$. For $0<t<1$,
\[
 |P(ts)|=t(1-t)|s|^2|ts-v|<2t(1-t)\le\frac12.
\]
The values at both endpoints are zero, so this nonzero root supplies
two distinct fibre points. If no nonzero quotient root exists, $f$ has
only one distinct zero value, contrary to the hypothesis.
\end{proof}

\paragraph{Formal scope.}
\label{rem:cubic-quotient-formal-boundary}
Lean checks the identity for these three real numbers, the choice of
a contained cubic segment, and the quartic counterexample,
and also \href{https://github.com/wcook04/plectis-erdos/blob/27c2fc5fe3c55fe547feaeb4c9bb68b3cf63a5bf/lean/ErdosProblems/Erdos1041/PaperCubicFibres.lean#L240}{the complete translated cubic fibre theorem}.  The latter supplies the finite set of roots in each fibre, root-of-unity average, nonzero-root selection, pullback and
explicit rectifiable path through the centre under the hypotheses displayed above.
Its exact source occurs in the successfully compiled dependency closure of
\href{https://github.com/wcook04/plectis-erdos/blob/27c2fc5fe3c55fe547feaeb4c9bb68b3cf63a5bf/verification/erdos1041-returned-r18-v5-full-audit-evidence.json}{the source-bound Lean build and axiom-audit record}.  The theorem gives length below $2$; the sharper
$2\sqrt{1-|h|^2}$ reading in the proof is recorded here as an ordinary
consequence of the mean-square identity, not as a separately audited endpoint.

\paragraph{Scope of the formal results.}
\label{bdry:solved-polynomial-families}
The finite declarations prove the Chebyshev comparison, the two-index
quintic inequality and the cubic segment estimate. The formal-source index also
records complete path theorems for the sparse quintic and the translated
cubic family in its checked build. These are stronger than the finite
inequalities alone. The sharper bound $2\sqrt{1-|h|^2}$ and the quartic
lifting arguments below retain the ordinary proof status stated beside them.
A formal declaration must still be compared with the exact hypotheses and
conclusion of the paper statement; a similar name is not enough.

\section{Further families and counterexamples to proposed proof steps}
\label{sec:frontier}

The following sections collect additional sufficient conditions and
counterexamples to particular path constructions. They are not a programme
for proving the unrestricted historical assertion. The underlying research
notes are retained at commit
\href{https://github.com/wcook04/plectis-erdos/tree/f214a6b45528dc5eefe20ffadc35f2e981627d4c/research_corpus/Erdos1041}{\texttt{f214a6b4}}
and summarised in \href{https://github.com/wcook04/plectis-erdos/blob/f214a6b45528dc5eefe20ffadc35f2e981627d4c/research_corpus/Erdos1041/FRONTIER.md}{the dated overview of these arguments}.
Where a claim is based on a computation rather than a proof, that
restriction is stated explicitly.

\subsection*{Roots close to a regular polygon}

Let $f(z)=\prod_{i=1}^n(z-a_i)$, $n\ge2$, with all $|a_i|\le1$.
Put $\rho_i=|a_i|$ and $D=|\operatorname{disc}(f)|/n^n$.  The
\href{https://github.com/wcook04/plectis-erdos/blob/f214a6b45528dc5eefe20ffadc35f2e981627d4c/research_corpus/Erdos1041/NearFeketeRadialAngularSplit.md}{Vandermonde determinant estimate}
says that, when $D\ge1-\eta$ and
$0\le\eta\le1/(80n^2)$,
\[
  1-\rho_i^2\le \frac{n\eta}{n-1},\qquad
  |a_i-a_j|\ge \frac{2-2\sqrt\eta-n\eta}{n-1},
\]
where the second inequality is for $i\ne j$. Moreover, the roots lie
within $7\sqrt\eta$ of a rotated regular $n$-gon, after a bijection.
These estimates come from near equality in Hadamard's determinant
inequality, as follows.

Take the Vandermonde matrix $V_{im}=a_i^m$, $0\le m<n$, and let
$G=VV^*$, $s_i=G_{ii}=\sum_{m=0}^{n-1}\rho_i^{2m}$ and
$H=\operatorname{diag}(s_i^{-1/2})G\operatorname{diag}(s_i^{-1/2})$.
The normalisation makes $H$ positive semidefinite with diagonal entries
one, and
\[
 D=\det H\prod_i\frac{s_i}{n}.
\]
Every factor is at most one by Hadamard's inequality, so $D\ge1-\eta$
forces each factor to be at least $1-\eta$. In particular,
\[
 (n-1)(1-\rho_i^2)\le n-s_i\le n\eta.
\]
Hadamard--Fischer applied to the $\{i,j\}$ principal block gives
$\det H\le1-|H_{ij}|^2$, hence $|G_{ij}|\le n\sqrt\eta$.
For $t=a_i\overline{a_j}$, the finite geometric sum then gives
\[
 n\sqrt\eta\ge\left|\sum_{m=0}^{n-1}t^m\right|
 \ge n-\frac{n(n-1)}2|1-t|.
\]
Together with
$1-t=1-\rho_j^2+\overline{a_j}(a_j-a_i)$, this proves the displayed
separation bound. These two bounds need only $0\le\eta<1$;
the smaller range is used for the matching with a regular polygon.

For that matching, put $\kappa=n^2\eta/(n-1)\le1/80$ and
$u_i=a_i/\rho_i$. The radial bound gives
$\rho_i^{2n}\ge1-\kappa$, while
$(1-t)G_{ij}=1-t^n$ gives
$|1-a_i^n\overline{a_j^n}|\le2n\sqrt\eta$ for $i\ne j$.
Using
$|1-be^{i\theta}|^2=(1-b)^2+b|1-e^{i\theta}|^2$ with
$b=\rho_i^n\rho_1^n\ge1-\kappa$, we obtain
\[
 |u_i^n-u_1^n|\le\frac{2n\sqrt\eta}{\sqrt{1-\kappa}}.
\]
Choose the nearest $n$th-root phase to $u_i$ among
$u_1e^{2\pi i j/n}$. The inequality
$|e^{i\theta}-1|\ge2|\theta|/\pi$ for $|\theta|\le\pi$
shows that its distance from $a_i$ is at most
\[
 \frac{n\eta}{n-1}+\frac{\pi\sqrt\eta}{\sqrt{1-\kappa}}
 \le7\sqrt\eta.
\]
The separation bound exceeds $14\sqrt\eta$ in the stated range,
so two roots cannot be assigned the same phase. There are $n$ roots and
$n$ phases, giving the required bijection. This is an ordinary argument
from standard determinant inequalities, not a new Lean declaration or
a claim of priority for the stability estimate.

The hypothesis $D\ge1-\eta$ with $\eta\le1/(80n^2)$ is a strong
quantitative assumption, not merely a requirement that roots lie near the
unit circle. The regular $n$-gon has $D=1$, whereas a configuration with
a repeated root has $D=0$. In particular, radial information alone cannot
supply this discriminant bound.

The comparison with unit-modulus roots is pointwise on a specified
segment, not an ordering of complex polynomial values. Write
$a_k=\rho_k u_k$, with $|u_k|=1$. Under
$D\ge1-\eta$ and $0\le\eta\le1/(10n^4)$, the same source proves
\[
 |f(su_i)|\le\prod_k|su_i-u_k|
 \qquad(0\le s\le\rho_i).
\]
The elementary step is
\[
 |su_i-\rho_ku_k|^2-|su_i-u_k|^2
 =(1-\rho_k)\bigl(2s\Re(u_i\overline{u_k})-1-\rho_k\bigr).
\]
Here is where the discriminant hypothesis enters. The two bounds just
proved imply
\[
 |a_i-a_k|>\frac1{n-1}\ (i\ne k),\qquad
 1-\rho_k\le\frac{n\eta}{n-1}<\frac1{(n-1)^2}.
\]
For $k=i$, the required inequality is $2s\le1+\rho_i$, which follows
from $s\le\rho_i\le1$. For $k\ne i$, it is immediate if
$\Re(u_i\overline{u_k})\le0$. Otherwise $s\le\rho_i$ gives
\[
\begin{split}
 \rho_k\bigl(2s\Re(u_i\overline{u_k})-1-\rho_k\bigr)
 &\le\rho_i^2-\rho_k-|a_i-a_k|^2\\
 &\le1-\rho_k-|a_i-a_k|^2<0.
\end{split}
\]
Multiplying the factorwise inequalities proves the comparison. This
argument uses the two elementary bounds, not the polygon-matching range;
it therefore also covers $n=2$, when $1/(10n^4)>1/(80n^2)$.
Moving the roots inward reduces this modulus bound, but does not select
the angular directions of two contained segments.

For exact regular-polygon directions, write
$a_i=\rho_i\omega\zeta^i$, with $|\omega|=1$ and
$\zeta=e^{2\pi i/n}$. In any degree $n\ge2$, the sufficient condition
\[
 0<\rho_k\le1,\qquad \rho_k\ge2\cos(2\pi/n)-1\quad\hbox{for every }k
\]
makes the same factorwise comparison hold. Indeed, for $k\ne i$ with
positive cosine, $s\le1$ gives
$2s\Re(\zeta^{i-k})\le2\cos(2\pi/n)\le1+\rho_k$;
nonpositive cosines and $k=i$ are handled as above. Since
$\prod_k(z-\omega\zeta^k)=z^n-\omega^n$, we obtain
\[
 |f(s\,\omega\zeta^i)|\le1-s^n
 \quad(0\le s\le\rho_i).
\]
Every root-to-origin segment is therefore contained in $\{|f|\le1\}$.
If all $\rho_i<1$, containment is strict: the displayed bound is below
$1$ for $s>0$, and $|f(0)|=\prod_i\rho_i<1$. Any two roots can then
be joined through the origin with length $\rho_i+\rho_j<2$.
For $2\le n\le6$ the lower bound on $\rho_k$ is nonpositive,
so every choice of positive radii at most one is allowed. For $n\ge7$
the argument instead requires each radius to be at least
$2\cos(2\pi/n)-1$. This is a sufficient condition for the factorwise
comparison, not a claimed necessary condition for a contained path.
The favourable index in the general unit-circle averaging identity may
move with $s$, however; an existence statement at one level does not
supply a fixed pair of roots.

\subsection*{The exact binomial chord calculation}
\hypertarget{binomial-chord-calculation}{}
For $f(z)=z^n-r^n$, $0<r<1$ and $n\ge2$, the construction in
Section~13 of the short note depends on the maximum over an entire
chord. The endpoint values alone do not give that maximum. Put
$\omega=e^{2\pi i/n}$ and $c=\cos(\pi/n)$. When $n\ge3$, define
$r_*=(1+c^n)^{-1/n}$ and $\varepsilon=(1-r^n)^{1/n}$; in the inner-chord
case $r\ge r_*$, set $t=\varepsilon/c\le r$.

Here is the chord estimate, including the smaller radii used in the
second construction. For $n\ge3$, let $0<s\le r$. A point on the chord
from $s$ to $s\omega$ can be written
\[
 z=\frac{sc}{\cos\theta}e^{i(\pi/n+\theta)},
 \qquad |\theta|\le\pi/n.
\]
Concavity of $\log\cos x$ on $[0,\pi/2)$ gives
$\cos(n\theta/2)\le\cos^{n/2}\theta$ for $0\le\theta<\pi/n$.
By symmetry and continuity,
\[
 1+\cos(n\theta)\le2\cos^n\theta
 \qquad (|\theta|\le\pi/n).
\]
Set $u=\cos^{-n}\theta\ge1$. Since
$z^n=-(sc)^nu e^{in\theta}$, this inequality gives
\[
 \begin{split}
 |z^n-r^n|^2-\bigl(r^n+(sc)^n\bigr)^2
 &\le (u-1)(sc)^n\bigl((sc)^n(u+1)-2r^n\bigr)\le0.
 \end{split}
\]
The last step uses $(sc)^nu\le s^n\le r^n$ and
$(sc)^n\le r^n$. Equality holds at the midpoint, where $\theta=0$.
Consequently the maximum on every such chord is exactly
$r^n+(sc)^n$. Taking $s=r$ proves the outer threshold; taking
$s=t=\varepsilon/c\le r$ proves the inner bound, and taking
$s=\lambda t<t$ makes it strict. Every radial leg between $s$ and
$r$ has modulus $r^n-|z|^n<1$. Its two legs and crossing chord have
length $2r-2s(1-\sin(\pi/n))<2r$, which proves the stated length
formula and shows that contraction preserves the strict length bound.

For $n=2$, the chord is a diameter and its maximum is $r^2<1$;
its length is $2r<2$. For $n\ge3$, the outer chord works exactly when
$r<r_*$, and its length is $2r\sin(\pi/n)$. At $r=r_*$ it reaches
level one at its midpoint. If $r\ge r_*$, the two radial legs and
inner crossing have length
\[
 2r-2\varepsilon\tan\left(\frac\pi4-\frac\pi{2n}\right)<2r<2.
\]
Contracting the crossing radius makes the level bound strict, as
proved above. The argument compares these constructions, not their
lengths with all admissible paths.

\subsection*{A bounded-radius concyclic class}

A different argument gives a chord when all roots lie on a circle of
sufficiently small radius.
Let $f$ be monic of degree $n\ge3$ with distinct zeros on a circle of radius
$\rho$.  If $2\rho^n\le1$, two adjacent zeros are joined by their straight
chord, whose length is at most $2\rho\sin(\pi/n)<2$, and the chord lies in
$\{|f|<1\}$.  (If a zero is repeated, the short-connection conclusion is
immediate; the distinct case is the substantive one.)

For each fixed degree $n$, the radius condition is $\rho\le2^{-1/n}$.
It allows arbitrary spacing on that circle, but does not cover every
radius $\rho<1$. No restriction on the circle's centre is imposed.

After translating and scaling, let $g$ be the monic polynomial with
roots $w_1,\ldots,w_n$ on the unit circle. Such a polynomial is
\emph{self-inversive}: $z^n\overline{g(1/\overline z)}$
is a constant multiple of $g(z)$, with the constant of modulus one.
Equivalently, $e^{-in\theta/2}g(e^{i\theta})$ is real up to a fixed phase.
Choose
\[
 c=(-1)^{n+1}\prod_j w_j,\qquad q(z)=z^n-c.
\]
Then $g$ and $q$ have the same leading and constant coefficients, and their
boundary values have a common real phase. If every root gap contained a
point with $|g|>|q|$, the real boundary function of $g-q$ would have
alternating signs at these points. Advancing the angle by $2\pi$ multiplies
that function by $(-1)^n$, so the alternation also holds across the final
gap. The intermediate value theorem would give $n$ distinct unit-circle
zeros of $g-q$. This is impossible: its degree is at most $n-1$, and
$g=q$ cannot satisfy the assumed strict inequalities. Thus some
adjacent-root arc satisfies $|g|\le|q|\le2$. The interior of that arc contains no zero of
$q$, since $g$ is nonzero there. As consecutive zeros of $q$ are separated
by angle $2\pi/n$, the selected arc has angular width at most $2\pi/n$.
This explains why the same selected pair also has the required short chord.
To pass from this arc to its chord $[a,b]$, let $\nu$ be the unit
normal pointing into the circular segment bounded by them. On the open
chord, logarithmic differentiation gives
\[
 \partial_\nu\log|g(z)|
 =\sum_{w_j\notin\{a,b\}}
   \frac{\operatorname{Re}\bigl(\nu\overline{z-w_j}\bigr)}{|z-w_j|^2}>0.
\]
The endpoint terms vanish because their directions are tangent to the
chord. Every remaining numerator is positive because the other roots
lie strictly on the far side of its line; $n\ge3$ ensures at least one
such term. The harmonic function $\log|g|$ has no interior maximum and
tends to $-\infty$ at $a,b$. A maximum cannot lie on the open chord
either, since moving into the circular segment increases the function.
It is therefore attained on the arc, and the maximum on the chord is
strictly smaller. Scaling back gives $|f|<2\rho^n\le1$ throughout the
chord, with $f=0$ at its endpoints. These are the arguments of
\href{https://github.com/wcook04/plectis-erdos/blob/f214a6b45528dc5eefe20ffadc35f2e981627d4c/research_corpus/Erdos1041/ConcyclicAlternation.md}{Theorem~C and Lemma~S, Sections~3--4 of the concyclic note}.

The argument is an ordinary proof outside Lean.  The
\href{https://github.com/wcook04/plectis-erdos/blob/f214a6b45528dc5eefe20ffadc35f2e981627d4c/research_corpus/Erdos1041/scripts/check_erdos1041_concyclic_exact_witness.py}{exact-rational checker}
checks finitely many identities used in the proof and configurations, while the
\href{https://github.com/wcook04/plectis-erdos/blob/f214a6b45528dc5eefe20ffadc35f2e981627d4c/research_corpus/Erdos1041/scripts/check_erdos1041_concyclic_alternation.py}{numerical checker}
is regression and stress-test evidence. The arc constant $2$ is attained
by the regular $n$-gon. Thus the bound obtained by passing through the
arc is $\rho\le2^{-1/n}$, although this cutoff tends to $1$ as
$n\to\infty$. This does not assert optimality of the chord bound or
exclude another contained path at a larger radius. Nor does arc
sharpness make the pair of radial segments shortest. For example, for
$g(z)=z^3-1/8$ the chord between $1/2$ and $e^{2\pi i/3}/2$ has length
$\sqrt3/2<1$, whereas the two radii have total length $1$. The binomial
calculation above gives a maximum of $9/64<1$ on that chord. The
length-$2$ equality at $z^n-1$, proved below, concerns the uncontracted
closed-level problem.

\subsection*{Power substitutions in trinomials}

The radial-segment identity and strict bound are formalized in
\lword{Erdos1041/CyclicTrinomialFiberCase.lean}{46}{trinomialRoot_spoke_factorization}{cancellation along a trinomial root segment}
and \lword{Erdos1041/CyclicTrinomialFiberCase.lean}{103}{trinomialRoot_spoke_norm_lt_one_of_norm_lt_one}{containment of a trinomial root segment}.

The trinomial identity also applies after a power substitution.  Fix
integers $1\le r<m$ and $q\ge1$, and consider the translated monic polynomial
\[
 f(z)=(z-h)^{qm}+a(z-h)^{qr}+c.
\]
The strict exponent inequality keeps the leading monomial separate.
The algebraic identity below also holds when $r=m$, but then the two
leading terms merge and the displayed polynomial need not be monic or
have degree $qm$.
Writing $w=(z-h)^q$ reduces the root equation to
$w^m+aw^r+c=0$.  At such a quotient root the middle coefficient can be
eliminated exactly: for $0\le u\le1$,
\[
 u^mw^m+au^rw^r+c
   =(1-u^r)c-(u^r-u^m)w^m.
\]
This is a nonnegative combination of $c$ and $-w^m$, with weights
$1-u^r$ and $u^r-u^m$ summing to $1-u^m$. The
\lword{Erdos1041/CyclicTrinomialFiberCase.lean}{103}{trinomialRoot_spoke_norm_lt_one_of_norm_lt_one}{formal radial-segment estimate}
gives strict containment when $|w|<1$ and $|c|<1$.

For $q\ge2$ and a nonzero quotient root satisfying these bounds, choose
two distinct solutions $y_1,y_2$ of $y^q=w$. Their common modulus is
$|w|^{1/q}<1$. To lift a quotient radial segment, set $u=t^q$:
\[
 |f(h+ty_i)|
 \le (1-t^{qm})\max\{|c|,|w|^m\}<1
 \qquad(0\le t\le1).
\]
Thus the two lifted segments join $h+y_1$ to $h+y_2$ through $h$
with length $2|w|^{1/q}<2$. The construction uses any quotient root
with the stated modulus bound; it does not prove that such a root
exists without an additional hypothesis. A zero quotient root instead
gives a root of multiplicity at least $q$ at $h$ and hence a constant
path between two occurrences. The formal source checks the
factorization, radial-segment estimate and finite length inequality.
It does not establish the existence of a suitable quotient root or
formalise the selection and lifting just described. Thus its scope is
the segment inequality and finite length calculation, not the complete
path construction or unrestricted Erd\H{o}s~\#1041.

\subsection*{A coefficient condition for four-term polynomials}

With one additional monomial, Abel summation gives a sufficient
condition involving the two lower coefficients.  Let
\[
  g(w)=w^m+aw^r+bw^s+c,\qquad m>r>s\ge1,
\]
assume all roots of $g$ lie in the open unit disk, and let $w_1,w_2$ be roots
of the two smallest moduli.  If
\[
  |c|+|b|\,|w_2|^s<1,
\]
then the complete radial spokes from $w_1$ and $w_2$ to the origin lie in
$\{|g|<1\}$, and their broken line has length strictly below $2$.  The simpler
coefficient condition
\[
  |b|+|c|\le1
\]
makes every root-to-origin segment contained; the coefficient $a$ has
no additional restriction beyond the assumed root locations. This coefficient
condition includes $b=0$ and $|c|\le1$, but is not automatic for all
four-term polynomials with roots in the disc.

The mechanism is visible in one identity.  At a root $w$ and for
$0\le u\le1$,
\[
\begin{split}
 g(uw)={}&(1-u^s)c
  +(u^s-u^r)(c+bw^s)
  +(u^r-u^m)(-w^m).
\end{split}
\]
Here the root equation has eliminated $aw^r$ before taking absolute
values. The three weights are nonnegative and sum to $1-u^m$.
The controlled terms are $c$, $c+bw^s$ and $-w^m$; their moduli are
strictly below one under the root-dependent condition. The same
identity explains why no bound on $a$ is needed.
Lean checks the \lword{Erdos1041/CyclicTetranomialCoefficientCase.lean}{27}{tetranomialRoot_spoke_factorization}{factorization}
and the resulting \lword{Erdos1041/CyclicTetranomialCoefficientCase.lean}{145}{tetranomialRoot_spoke_norm_lt_one_of_rootBudget}{strict spoke theorem} under
the exact weak exponent hypotheses $1\le s\le r\le m$; it also checks the
coefficient-only corollary above. Its weak bound $|b|+|c|\le1$ still
gives strict containment: the root-disc hypothesis implies $|c|<1$,
and, if $b\ne0$, $|w|^s<1$ gives
$|c|+|b||w|^s<|c|+|b|\le1$. If $b=0$, the required bound is
simply $|c|<1$.

An alternative criterion uses a complex power sum to select two
indices. Let $S$ index a finite family of roots, $N=|S|\ge2$, and
\[
  M=\sum_{i\in S}w_i^s.
\]
The \lword{Erdos1041/TetranomialL2Selector.lean}{33}{sum_normSq_const_add_mul}{squared-modulus identity} gives
\[
 \sum_{i\in S}|c+bw_i^s|^2
 =N|c|^2+|b|^2\sum_{i\in S}|w_i^s|^2
   +2\mathop{\rm Re}(\overline c\,bM).
\]
Consequently, if every $|w_i|<1$, $|c|<1$, and
\[
 N\bigl(|b|^2+|c|^2\bigr)
   +2\mathop{\rm Re}(\overline c\,bM)<N-1,
\]
then two distinct indices $i,j\in S$ satisfy
$|c+bw_i^s|<1$ and $|c+bw_j^s|<1$. Indeed, the displayed identity
bounds the sum of squared moduli strictly below $N-1$. If at most one
index had modulus below one, the other $N-1$ terms would each contribute
at least one. This explains both the threshold $N-1$ and why the
criterion selects two indices. Applying the
\lword{Erdos1041/TetranomialL2Selector.lean}{208}{exists_two_tetranomialRoot_safeSpokes_of_moment_coeff_budget}{resulting contained-segment theorem} to these inequalities proves that both root segments lie
strictly in $\{|g|<1\}$.  Retaining the signed cross term permits the inequality to hold for configurations outside the coefficient-only triangle
$|b|+|c|\le1$.  The formal hypotheses do not require $i\mapsto w_i$ to be
injective, so distinct indices need not denote distinct root values.
For the polynomial consequence, the family must be a submultiset of the
zeros, respecting their algebraic multiplicities. A repeated zero then
already gives a constant path between two occurrences. Repeating a simple
zero in an arbitrary indexed family cannot supply a second occurrence.
The two sufficient conditions are not ordered by strength. For example,
$g(w)=w^5+19/20$ has all roots in the open unit disc and satisfies the
coefficient condition with $b=0$. For its full five-root list, however,
the left side of the power-sum condition is
$5(19/20)^2=361/80>4=N-1$, so that condition fails. Its benefit is the
use of cancellation in the cross term, not inclusion of every case
covered by the coefficient condition. Conversely,
\[
 g(w)=(w^2-1/2)(w^3+4/5)
     =w^5-\tfrac12w^3+\tfrac45w^2-\tfrac25
\]
has all roots of modulus $1/\sqrt2$ or $(4/5)^{1/3}$, hence below
one, but $|b|+|c|=6/5>1$. Its full root list has
$M=\sum_iw_i^2=1$: the two quadratic roots contribute $1$, and the
three cubic roots contribute $0$. The power-sum left side is
$5(16/25+4/25)-16/25=84/25<4$. Thus this criterion does cover
root-disc polynomials excluded by the coefficient condition.

For $q\ge2$ and $f(z)=g((z-h)^q)$, assume also that all roots of $f$
lie in the open unit disc. If $y^q=w$ and $\omega=e^{2\pi i/q}$, the
roots $h+\omega^j y$ satisfy
\[
 \frac1q\sum_{j=0}^{q-1}|h+\omega^j y|^2=|h|^2+|y|^2<1.
\]
Thus every quotient root has $|w|=|y|^q<1$, which supplies the
root-disc hypothesis used above, and every lifted spoke has length
$|y|<1$. Two selected spokes join through $h$ with total length below
$2$; repeated root occurrences instead permit the constant path.
The Lean theorem assumes the finite root family, its complex power sum $M$,
and the displayed inequality. It neither derives that inequality from
arbitrary coefficients nor verifies the family's polynomial
multiplicities, the cyclic lifting, or the final path construction.
Completeness of the root list is not required for the selection argument:
a submultiset respecting multiplicity suffices.  Without one of the power-sum, root-dependent, or coefficient-only
inequalities, the tetranomial case is not covered by this argument; this family does not solve
unrestricted Erd\H{o}s~\#1041.

\subsection*{A quartic composed with a power map}\label{sec:translated-quartic-quotient-fibres}

Pendyala's degree-four theorem \cite[Theorem~1]{june2026}, with his four-point
radial lemma \cite[Lemma~1]{june2026}, can also be lifted through every
nontrivial cyclic power. Let $P$ be a monic quartic, let $q\ge2$ be an
integer, and set
\[
  f(z)=P((z-h)^q).
\]
Assume that every root of $f$ lies in the open unit disc.
A repeated root of $f$ gives the constant path, so suppose $f$ is
squarefree. Then the roots of $P$ are distinct and nonzero. For a root
$w=y^q$ of $P$, all points $h+\zeta^j y$, with $\zeta=e^{2\pi i/q}$,
are roots of $f$. Hence the open-disc hypothesis gives
\[
 \frac1q\sum_{j=0}^{q-1}|h+\zeta^jy|^2=|h|^2+|y|^2<1.
\]
In particular $|w|^{1/q}=|y|<1$, which is the bound needed for the
lifted length, and also $|w|<1$. Pendyala's chord-or-radial argument
supplies the quotient geometry in two cases.  If two quotient roots satisfy
$|w_i-w_j|<1$, every point $w$ of the chord $[w_i,w_j]$ lies in the unit disc;
the two endpoint factors of $|P(w)|$ have product at most $|w_i-w_j|^2/4$ and
the other two factors are below $2$, so $|P(w)|\le|w_i-w_j|^2<1$.  Otherwise
the quotient roots are pairwise at distance at least $1$.  Choosing $R$ with
$\max_k|w_k|<R<1$ and applying the four-point radial lemma to the points
$w_k/R$ gives distinct $i,j$ with $|P(tw_i)|\le R^4$ and $|P(tw_j)|\le R^4$ for
$0\le t\le1$, so the radial segments $[0,w_i]$ and $[0,w_j]$ lie in
$\{|P|<1\}$.  The extra issue is metric: a short chord in the
$w$-plane need not lift isometrically through $y\mapsto y^q$.

Put $\alpha=1/q$. On a chord avoiding $0$, choose a continuous
argument and hence a root lift $y$ with $y^q=w$. Differentiation gives
$|dy|=\alpha|w|^{\alpha-1}|dw|$, independently of the chosen root.
If the supporting line has distance $d$ from $0$, and $x$ is distance
along that line from the perpendicular foot, then
\[
  \bigl(\sqrt{d^2+x^2}\bigr)^{\alpha-1}
    \le x^{\alpha-1}\qquad(x>0).
\]
The negative exponent is why replacing $|w|$ by $x$ gives an upper
length bound. Lean checks \href{\repobase/research_corpus/Erdos1041/QuarticQuotientFiberCase.lean\#L37}{the integral of the power-map derivative}
\[
  \alpha\int_0^A x^{\alpha-1}\,dx=A^\alpha.
\]
If the foot lies on the chord, split there: the two
integrals are bounded by the $\alpha$-powers of the distances from the
foot to the endpoints, hence by $|a|^\alpha+|b|^\alpha$. If the foot
lies outside the chord, integrate on one side of it; extending the
interval to the foot bounds the result by the $\alpha$-power of the
farther endpoint's modulus. When the chord passes through $0$, choose
a continuous root on each half and join the lifts at $0$; the integral
is finite because $\alpha>0$. This proves the chord-lift estimate
\[
  \operatorname{length}(\widetilde{[a,b]})
    \le |a|^{1/q}+|b|^{1/q}.
\]
The \href{\repobase/research_corpus/Erdos1041/QuarticQuotientFiberCase.lean\#L59}{strict endpoint-length estimate} also proves that if
$0\le a,b<1$, $\alpha>0$, and a candidate length $L$ is at most
$a^\alpha+b^\alpha$, then $L<2$.  The close-pair chord therefore has a lift of length below $2$.
In the other case the quotient arms meet at $0$, so their lifts are radial
segments meeting at $h$, with total length
$|w_i|^{1/q}+|w_j|^{1/q}<2$. In both cases $f=P\circ((z-h)^q)$
transfers containment from the quotient curve. Distinct quotient endpoints
give distinct roots of $f$. This proves the path assertion in every degree
$4q\ge8$. The fibre average gives a little more: $|h|<1$ and
$|w_k|^{1/q}<\sqrt{1-|h|^2}$ for every quotient root. Consequently the same
constructed path has length
\[
 \operatorname{length}<2\sqrt{1-|h|^2}\le2.
\]
This refinement concerns the translated power-substitution family, not
arbitrary degree-$4q$ polynomials.

Pendyala proves the quartic geometric theorem and its four-point radial
lemma.  The local Lean module
checks the antitone density inequality, its integral, the strict powered
endpoint-length bound, and the final length estimate.  It does not formalize
Pendyala's geometric lemma, the continuous covering-space construction of the
root lift, or the ordinary chord/radial case assembly.  This is a family defined by a power substitution, not a proof of unrestricted Erd\H{o}s~\#1041.

\paragraph{The joining point in the quartic construction.}
Pendyala's proof of \cite[Theorem~1, pp.~2--3]{june2026} uses a
smallest enclosing disc and joins the selected roots through its centre.
For the power substitution we instead apply \cite[Lemma~1]{june2026}
to the centred disc $|w|\le R$. That lemma requires only the disc and
pairwise-distance hypotheses, not minimality of the disc. Keeping the
joining point at $0$ is what makes the radial lift above valid. The
endpoint-length bound for a single chord would not justify lifting an
arbitrary broken line through a nonzero quotient vertex.

\subsection*{A mean bound for critical values in every degree}

The next theorem gives the quadratic case of a Poisson argument for
critical-value moments. It allows a disc of any centre and radius,
repeated roots, and repeated critical points. After proving this case,
we derive the stronger exponent $4/(n-1)$ and explain how vanishing
complex power sums permit still larger exponents. None of these moment
bounds controls the lengths of inverse paths ending at the critical points.

Dubinin \cite[Theorem~2, (9), pp.~1172--1173]{dubinin2006critical}
proves the sharp unit-disc critical-value product inequality using the
resultant identity, the maximum-modulus principle and Schur's Vandermonde
inequality. The positive-moment estimate below implies his product bound
by AM--GM. When $n\ge3$, the product estimate alone cannot control a
positive moment: a list $T,T^{-1},1,\ldots,1$ has product one, but its
sum of $p$th powers is unbounded as $T\to\infty$ for every $p>0$.
This compares the information in two inequalities for nonnegative lists;
it does not claim that these lists occur as critical values of polynomials
in the theorem. In degree two there is just one critical value, and the
product and positive-moment bounds are equivalent.
For a zero prescribed at the origin, his Theorem~3 (p.~1174) gives
\[
 \left(\prod_j|f(c_j)|\right)^{1/(n-1)}
 \le(n-1)\left(\frac{|f'(0)|^2}{n^n}\right)^{1/(n-1)}.
\]
Dubinin explicitly identifies Tischler's earlier contribution
\cite[p.~444, as reported in Dubinin]{tischler1989}.
His Theorem~1 uses dissymmetrisation; that is not the proof mechanism of
Theorem~2.  Schur's original paper \cite{schur1918} is the historical
antecedent cited there, not a separately re-proved source in this record.
A sharp marked-zero positive-moment bound remains a separate question.
The addendum \cite[Section~3]{revision2026} reports a reduction to
boundary-root configurations, without solving the resulting angular
optimisation. That unavailable addendum is not a proof dependency of the
critical-value theorem below.

\begin{theorem}[a mean bound for critical values]
\label{res:critical-value-budget}
Let $f$ be monic of degree $n\ge2$, with roots in a closed disc of radius
$R\ge0$.  If $c_1,\ldots,c_{n-1}$ are its critical points counted with
multiplicity, then
\begin{equation}\label{eq:critical-value-quadratic-budget}
 \sum_{j=1}^{n-1}|f(c_j)|^{2/(n-1)}
 \le(n-1)R^{2n/(n-1)}.
\end{equation}
Consequently the lower exponents used elsewhere in the record satisfy
\begin{equation}\label{eq:critical-value-power-budget}
 \sum_{j=1}^{n-1}|f(c_j)|^{1/(n-1)}
 \le(n-1)R^{n/(n-1)}
\end{equation}
and
\[
 \sum_{j=1}^{n-1}|f(c_j)|^{1/n}\le(n-1)R.
\]
The constant is attained by $f(z)=(z-\tau)^n-\lambda$ with enclosing
disk centred at $\tau$ and radius $R=|\lambda|^{1/n}$.
\end{theorem}
\leannote{Lean: \lproof{ErdosProblems/Erdos1041/PaperCompleteR20/CriticalMeanWhole.lean}{14}{critical_value_three_budgets}, \lproof{ErdosProblems/Erdos1041/PaperCompleteR20/CriticalMeanWhole.lean}{36}{critical_value_three_budgets_sharp}.}

The finite inequality behind the theorem concerns arbitrary points of the
disk, without asking them to arise as critical points.  Its quadratic form is
\[
 \sum_{j=1}^m\left(\prod_{k=1}^m
       |1-\overline{c_j}c_k|\right)^{2/m}\le m,
 \qquad |c_j|\le1.
\]
The following pointwise bound explains why these finite inequalities
control critical values in every degree.

\begin{lemma}[reflected-derivative bound]\label{res:reflected-critical-value}
If $f$ is monic of degree $n\ge2$ with roots in the closed unit disk,
and $c_1,\ldots,c_{n-1}$ list its critical points with multiplicity, then
\begin{equation}\label{eq:critical-reflected-product}
 |f(c_j)|\le\prod_{k=1}^{n-1}|1-\overline{c_j}c_k|.
\end{equation}
\end{lemma}
\leannote{Lean: \lproof{ErdosProblems/Erdos1041/PaperReflectedCompletion.lean}{254}{reflected_critical_value}.}

\noindent\href{https://github.com/wcook04/plectis-erdos/blob/bb4e24651b37cd096d3749ef299c3317930b3a3b/ErdosProblems/Erdos1041/PaperReflectedCompletion.lean#L254}{The reflected-derivative inequality is checked in Lean, including closed-disc roots and critical-point multiplicities.}

\begin{proof}
First put all roots $a_i$ strictly inside the disk. We want a numerator
that equals $nf$ at the critical points and can be compared with $f'$
on the circle. Set
\[
 N(z)=nf(z)-zf'(z),\qquad
 G(z)=n\prod_k(1-\overline{c_k}z).
\]
The numerator is the polar derivative at $0$: the conventional polar
derivative at $\alpha$ is $nf(z)+(\alpha-z)f'(z)$; see
\href{https://arxiv.org/abs/1303.0068v1}{Rather--Gulzar, Section~1} for this
notation. Its leading term cancels, and $N(c_j)=nf(c_j)$.
Reflection gives $|G|=|f'|$ on the circle. Gauss--Lucas puts every
$c_k$ inside the disc, so $G$ has no zero on its closure. Thus $N/G$
is a holomorphic quotient to which the maximum-modulus principle applies.  On $|\zeta|=1$,
\[
 \operatorname{Re}\frac{\zeta f'(\zeta)}{f(\zeta)}
 =\sum_i\left(\frac12+
       \frac{1-|a_i|^2}{2|1-a_i\bar\zeta|^2}\right)\ge\frac n2.
\]
Writing the logarithmic derivative as $w$, the identity
$|n-w|^2-|w|^2=n^2-2n\operatorname{Re}w\le0$ gives
$|N|\le|G|$ on the circle.  The maximum-modulus principle applied to
$N/G$ gives the same comparison inside.  At $c_j$, $N(c_j)=nf(c_j)$;
dividing by $n$ and conjugating each factor proves
\eqref{eq:critical-reflected-product}.
For closed-disk roots apply this result to
$f_r(z)=r^nf(z/r)$, $0<r<1$, whose critical points are $rc_k$:
\[
 r^n|f(c_j)|\le\prod_k|1-r^2\overline{c_j}c_k|.
\]
Let $r\uparrow1$.  This also handles boundary critical points and all
multiplicities without selecting local branches.
\end{proof}

Set $m=n-1$.  The weighted Poisson proof below gives the quadratic finite
inequality for every $m\ge1$.  Gauss--Lucas and
Lemma~\ref{res:reflected-critical-value} then give
$\sum_j|f(c_j)|^{2/m}\le m$ on the unit disc.  Applying this to
$R^{-n}f(h+Rz)$ and scaling back proves
\eqref{eq:critical-value-quadratic-budget} for $R>0$; if $R=0$, all critical
values vanish.

The two displayed consequences are the lower power means of the same
nonnegative $m$-tuple.  Equivalently, after unit-disc normalization, apply
the monotonicity of normalized $L^p$ means from exponent $2/m$ first to
$1/m$ and then to $1/(m+1)=1/n$.  Scaling back contributes respectively
$R^{n/m}$ and $R$. These are consequences of the quadratic case;
the higher-power argument below strengthens that case without changing
these lower-exponent applications.

The quadratic mean and both lower-exponent consequences are recorded as checked in
\href{https://github.com/wcook04/plectis-erdos/blob/27c2fc5fe3c55fe547feaeb4c9bb68b3cf63a5bf/lean/ErdosProblems/Erdos1041/PaperCriticalValueMeanR10.lean#L101}{the complete critical-value mean}.  This includes arbitrary centre, zero radius,
and critical points counted with multiplicity.  The source version and
formal target are identified in \href{https://github.com/wcook04/plectis-erdos/blob/27c2fc5fe3c55fe547feaeb4c9bb68b3cf63a5bf/verification/erdos1041-returned-r18-v5-full-audit-evidence.json}{the source-bound Lean build and axiom-audit record}. The earlier candidate description in the source header predates that run.

The quadratic moment controls the number of large critical values.
For $R>0$, write $r_j=|f(c_j)|^{1/n}$. Then
\[
 \sum_j(r_j/R)^{2n/(n-1)}\le n-1,\qquad
 \#\{j:r_j\ge tR\}\le\frac{n-1}{t^{2n/(n-1)}}\quad(t>0).
\]
The counting bound follows by retaining just the indicated summands.
The fourth-power refinement proved below replaces the exponent
$2n/(n-1)$ here by $4n/(n-1)$, improving the count for $t>1$.
For $0<t\le1$ the trivial bound $n-1$ is at least as good.
The binomial equality family has every $r_j=R$, but these distributional
bounds still leave the geometry of the joining paths to be supplied.

\paragraph{The weighted Poisson inequality and its equality case.}
\hypertarget{weighted-poisson-proof}{}
For one point $c$ of weight $1$, the quantity to bound is
$(1-|c|^2)^2\le1$; equality holds only at $c=0$.
For two points $c,-c$ of equal weight, it is $1-|c|^4\le1$.
These examples suggest why a nonzero configuration should lose from equality.
Now let $w_j\ge0$, $\sum_jw_j=1$, and $|c_j|\le1$, with repetitions allowed.
Write $G(z)=\prod_k|1-\overline{c_k}z|^{w_k}$, taking a factor of
exponent zero to be one. We seek an analytic function with modulus $G$:
its logarithmic derivative will express the weighted Poisson kernel,
and its Taylor coefficients will measure the loss from equality.
For $|c|<1$ and $|\zeta|=1$, let
$P_c(\zeta)=(1-|c|^2)/|\zeta-c|^2$, and write $dm=dt/(2\pi)$ on
$\zeta=e^{it}$. When every centre is strictly inside the disc, choose
analytic logarithms zero at the origin and put
\[
 g(z)=\exp\sum_jw_j\log(1-\overline{c_j}z)
     =1+\sum_{\nu\ge1}a_\nu z^\nu,\qquad
 P=\sum_jw_jP_{c_j}.
\]
On the unit circle $P=1-2\Re(\zeta g'/g)$, by logarithmic
differentiation. The Poisson formula reproduces $g(c)$ from its boundary
values and gives $\int P_c\,dm=1$ for $|c|<1$. Expanding a square therefore
yields
\[
 \int |g|^2P_c\,dm-|g(c)|^2
   =\int |g-g(c)|^2P_c\,dm\ge0.
\]
This explains why the same kernel bounds the value at each centre.
Take $c=c_j$, multiply by $w_j$ and sum to obtain
\[
 \sum_jw_jG(c_j)^2\le\int |g|^2P\,dm
 =1-\sum_{\nu\ge1}(2\nu-1)|a_\nu|^2.
\]
To see the last equality, expand $|g|^2$ and
$\zeta g'\overline g$ and integrate on the circle. Orthogonality leaves
$1+\sum_{\nu\ge1}|a_\nu|^2$ and
$\sum_{\nu\ge1}\nu|a_\nu|^2$, respectively. The analytic function $g$ extends to a neighbourhood of the closed disc
in this interior case,
which justifies the termwise integrations.
For closed-disc centres replace $c_j$ by $rc_j$, $0<r<1$.
The actual coefficient of degree $\nu$ becomes $r^\nu a_\nu$, where the
fixed coefficients are defined from the analytic function near zero.
Pass to $r\uparrow1$ first with an arbitrary finite coefficient sum.
Nonnegative finite deficits then give summability and the infinite inequality;
no boundary holomorphic logarithm is required. Lean checks the
\href{https://github.com/wcook04/plectis-erdos/blob/1d70482583cc9c662cd3b86f66d75c54bb405aa6/lean/ErdosProblems/Erdos1041/WeightedCauchyTransportR10.lean\#L66}{exact radial transport of each Cauchy coefficient},
the
\href{https://github.com/wcook04/plectis-erdos/blob/1d70482583cc9c662cd3b86f66d75c54bb405aa6/lean/ErdosProblems/Erdos1041/WeightedBoundaryDeficitR10.lean\#L108}{finite boundary-deficit inequality},
and the
\href{https://github.com/wcook04/plectis-erdos/blob/1d70482583cc9c662cd3b86f66d75c54bb405aa6/lean/ErdosProblems/Erdos1041/WeightedBoundaryDeficitR10.lean\#L153}{summable full boundary deficit}.

If equality holds, all positive-degree coefficients vanish, so $g=1$ near
zero and $\sum_jw_j\overline{c_j}/(1-\overline{c_j}z)=0$ there.
Group equal centres before clearing denominators. A nonzero centre of positive
total weight gives a nonzero pole coefficient, which is impossible. Thus
equality is equivalent to $c_j=0$ on positive support. In particular,
strictly positive weights force all centres zero. A point of weight zero
is unconstrained: weight one at $0$ and weight zero at $1$ still give
equality.
Grouping equal centres is necessary: an argument that assumed they
were distinct would not prove the stated equality case.

Equal weights in $\sum_jw_jG(c_j)^2\le1$ give the quadratic finite
inequality above for every $m\ge1$, with equality only at the zero
configuration. The linear inequality follows separately. For
$M=\sum_jw_jG(c_j)$, the variance identity
\[
 M^2+\sum_jw_j(G(c_j)-M)^2=\sum_jw_jG(c_j)^2
\]
gives $M\le1$. Equality in this Cauchy--Schwarz step means that the
values $G(c_j)$ are constant at all indices of positive weight; equality
in the final bound also requires equality in the quadratic inequality. The
three-point and four-point specialisations follow with their equality
cases.  The formal-source index records the complete weighted theorem in its
checked build. The hypotheses permit repeated centres and zero weights;
the equality condition refers only to centres of positive weight.

This equality condition also identifies all equality cases of
Theorem~\ref{res:critical-value-budget} for a fixed containing disc
$\overline D(h,R)$ with $R>0$. Equality in the critical-value bound
forces equality in the quadratic finite inequality after normalisation,
so every critical point is $0$. Thus the normalised derivative is
$nz^{n-1}$ and the polynomial is $z^n-\lambda$, with $|\lambda|=1$
forced by equality. Restoring the centre and radius gives exactly
$f(z)=(z-h)^n-\lambda$ with $|\lambda|=R^n$.
Equality in either lower-power consequence also forces equality in the
quadratic bound, so it has the same family. For $R=0$ the only possible
polynomial is $(z-h)^n$. This classification follows from the ordinary
proof; it is not an additional assertion about the cited Lean endpoints.

The earlier three-point H\"older argument and four-point matching argument
are not used here. Their scalar identities do not justify the omitted
inequalities; the proof is the weighted Poisson calculation above.

\paragraph{Higher powers and vanishing complex power sums.}
The quadratic inequality is the formally recorded case, not the full
range of the analytic argument. To estimate $G^p$ for $p>0$, we need an
analytic function whose squared modulus is $G^p$. Using the same
logarithms as above, define
\[
 H_p=g^{p/2}=1+\sum_{\nu\ge1}b_\nu z^\nu.
\]
For interior centres, logarithmic differentiation and the Poisson
majorant give
\begin{align*}
 P&=1-\frac4p\Re\frac{\zeta H_p'}{H_p},\\
 \sum_jw_jG(c_j)^p
 &\le\int |H_p|^2P\,dm
 =1-\sum_{\nu\ge1}\left(\frac{4\nu}{p}-1\right)|b_\nu|^2.
\end{align*}
The identity follows from the same two orthogonality calculations used
for $g$; only the coefficient of the derivative term changes. For
$0<p\le4$ every displayed coefficient is nonnegative. Hence
\[
 \sum_jw_jG(c_j)^p\le1\qquad(0<p\le4).
\]
Boundary centres follow by replacing $c_j$ by $rc_j$: the degree-$\nu$
coefficient becomes $r^\nu b_\nu$, and the values on the left are
$G_r(rc_j)=\prod_k|1-r^2\overline{c_k}c_j|^{w_k}$.
First retain finitely many nonnegative terms, then let $r\uparrow1$.
Zero weights can be omitted throughout. This proves the closed-disc
inequality without differentiating a boundary logarithm. The argument
is the ordinary
\href{https://github.com/wcook04/plectis-erdos/blob/f36a98bf3d3e6f65f1074e3b800e3293d5b8a51a/ErdosProblems/Erdos1041/SharpPowerDiscProducts.md}{power-parameter Poisson identity},
not an additional conclusion of the cited complete Lean theorem.

In particular, equal weights and Lemma~\ref{res:reflected-critical-value}
give, with $m=n-1$,
\[
 \frac1m\sum_j|f(c_j)|^{4/m}\le\frac1m\sum_jG(c_j)^4\le1
\]
on the unit disc. For roots in $\overline D(h,R)$ this becomes
\[
 \sum_j|f(c_j)|^{4/(n-1)}\le(n-1)R^{4n/(n-1)}.
\]
As before, $R>0$ is handled by $R^{-n}f(h+Rz)$, and $R=0$ by the
vanishing of all critical values. The constant is attained by
$(z-h)^n-\lambda$ with $|\lambda|=R^n$. This proves sharpness of the
constant, not optimality of the exponent for each fixed degree.

At $p=4$ the coefficient of $|b_1|^2$ vanishes, so the previous
quadratic equality proof cannot simply be reused. If equality holds,
the finite-deficit bounds imply $H_4=1+b_1z$. If $b_1\ne0$, the
rational identity $H_4'/H_4=2g'/g$ shows, by comparing poles and their
residues, that all nonzero centres coincide at $a=-\overline{b_1}$
and have total weight $1/2$. The other half of the weight is at zero.
Their fourth-power average is
\[
 \tfrac12\bigl(1+(1-|a|^2)^2\bigr)
 =1-|a|^2+\tfrac12|a|^4<1\qquad(0<|a|\le1),
\]
a contradiction. Thus $H_4=1$, and the pole argument already used for
$g$ forces every centre of positive weight to be zero. Conversely that
configuration gives equality. It follows that the fourth-power
critical-value estimate has the same equality family as the quadratic
one for a fixed containing disc.

For $p>4$, the coefficient $4/p-1$ is negative, so the nonnegative-deficit
argument no longer applies. To extend its range we can force the first
Taylor coefficients to vanish. Let $s\ge1$ be an integer and suppose
\[
 \sum_jw_jc_j^k=0\qquad(1\le k<s).
\]
The expansion
\[
 \log H_p(z)=-\frac p2\sum_{k\ge1}
       \frac{\overline{\sum_jw_jc_j^k}}{k}z^k
\]
then has no terms of degree below $s$, so $b_1=\cdots=b_{s-1}=0$. All remaining
coefficient weights are nonnegative for $0<p\le4s$. The same proof
therefore gives $\sum_jw_jG(c_j)^p\le1$ throughout that range.
For a monic polynomial with roots in $\overline D(h,R)$, $R>0$,
the resulting endpoint is
\[
 \sum_j|f(c_j)|^{4s/(n-1)}\le(n-1)R^{4sn/(n-1)},
 \qquad \sum_j(c_j-h)^k=0\quad(1\le k<s).
\]
There is no cancellation hypothesis when $s=1$. For $s=2$, the condition
says that the critical-point centroid is the disc centre. Comparing the
coefficients of $f$ and $f'$ shows
\[
 \frac1{n-1}\sum_jc_j=\frac1n\sum_{i=1}^n z_i,
\]
so this is also the root centroid. Recentring can increase the required
radius: for $0<r<1$, the roots of $(z-r)^2(z+r)$ lie in
$\overline D(0,r)$, but their centroid is $r/3$, and a disc centred
there needs radius $4r/3$ to contain the root $-r$. Thus the
exponent-$8/(n-1)$ statement cannot be applied at the old centre with
the old radius. These are exact cancellations, not merely bounds on the
moments. All higher-power conclusions in this paragraph
are ordinary analytic results; the linked complete Lean moment endpoint
has exponent $2/(n-1)$.

\paragraph{Formal scope of the Poisson and integration steps.}
The square-expansion proof above applies whenever $g$ is holomorphic on
the disc and continuous on its closure, with $|c|<1$. The linked
\href{https://github.com/wcook04/plectis-erdos/blob/d4fed71423840f70f10edf27b9ad27c22fc4f49a/ErdosProblems/Erdos1041/PoissonNormSquare.lean#L114}{norm-square majorization}
uses the Poisson formula for a real part and the nonnegativity of a square.
The \href{https://github.com/wcook04/plectis-erdos/blob/d4fed71423840f70f10edf27b9ad27c22fc4f49a/ErdosProblems/Erdos1041/PoissonKernelBridge.lean#L27}{weighted kernel identity}
identifies the pointwise Poisson mixture.  Separately,
\href{https://github.com/wcook04/plectis-erdos/blob/d4fed71423840f70f10edf27b9ad27c22fc4f49a/ErdosProblems/Erdos1041/CircleSeriesTransport.lean#L49}{absolute series transport}
justifies termwise circle integration when the terms are circle integrable
and admit a summable uniform norm bound; its Taylor double-product form
assumes absolute summability of both coefficient sequences.
The complete weighted and critical-value targets are indexed in the
recorded checked dependency closure.  Their exact statements, dependencies and status
are recorded in the accompanying formal-scope map.  Historical totals of
named axiom prints are not an additional proof of a paper statement; no
fresh axiom audit was run for this revision.

\paragraph{A uniform central region in every degree.}
For every $m\ge1$, the
\href{https://github.com/wcook04/plectis-erdos/blob/3be82b1a7340284aea72e9a5c8493cb020843921/ErdosProblems/Erdos1041/FreePointCentralCompletion.lean#L39}{inequality for points in a smaller disc}
proves
\[
 |c_j|\le\sqrt{1-e^{-2}}\quad(1\le j\le m)
 \quad\Longrightarrow\quad
 \sum_{j=1}^m\left(\prod_{k=1}^m|1-\overline c_jc_k|\right)^{1/m}\le m.
\]
To see the mechanism, set
$h_j=m^{-1}\sum_k\log|1-\overline c_jc_k|$ and
$D_j=-\log(1-|c_j|^2)$. Write
$M_\nu=m^{-1}\sum_k c_k^\nu$ for the normalised complex power sums used in this
calculation. The absolutely convergent series for $\log(1-z)$ gives
\[
 h_j=-\Re\sum_{\nu\ge1}\frac{\overline c_j^{\,\nu}M_\nu}{\nu},
 \qquad
 -\frac1m\sum_jh_j=\sum_{\nu\ge1}\frac{|M_\nu|^2}{\nu}.
\]
In particular, $\sum_jh_j\le0$. Cauchy--Schwarz applied to the first
series, using
$\sum_{\nu\ge1}|c_j|^{2\nu}/\nu=D_j$, gives
$h_j^2\le(-m^{-1}\sum_i h_i)D_j$. The radius bound gives $D_j\le2$.
Put $r=\sqrt{-m^{-1}\sum_i h_i}$. If $r=0$, every $h_j$ vanishes.
For $r>0$, convexity bounds $e^h$ by the chord of its graph over
$[-\sqrt2r,\sqrt2r]$. Averaging and using $m^{-1}\sum_jh_j=-r^2$ gives
\[
 \frac1m\sum_j e^{h_j}
 \le\cosh(\sqrt2r)-\frac r{\sqrt2}\sinh(\sqrt2r)\le1.
\]
For the last inequality, the function $\cosh u-(u/2)\sinh u$ equals
$1$ at $u=0$ and has derivative $(\sinh u-u\cosh u)/2\le0$ for
$u\ge0$. Thus the radius bound supplies the estimates needed for the
geometric-mean inequality; no further series hypotheses are imposed.
This alternative proof includes the all-zero configuration but requires
the smaller radius $\sqrt{1-e^{-2}}$. The weighted theorem above covers
the whole closed unit disc. The next logarithmic argument imposes a
different sufficient radius condition; neither restriction belongs to
the weighted theorem.

For $m\ge1$, let $|c_i|^2\le a<1$, with $a\ge0$, and put
$L=-\log(1-a)$. The
\href{https://github.com/wcook04/plectis-erdos/blob/457e0c74b2666f7fb2d825c967e49e439a600b23/ErdosProblems/Erdos1041/FreePointUniformRadius.lean#L84}{uniform-radius theorem}
gives the geometric-mean inequality
\[
 e^L\le1+2L\quad\Longrightarrow\quad
 \sum_{i=1}^m\left(\prod_{j=1}^m|1-\overline c_i c_j|\right)^{1/m}\le m.
\]
For $a=1/2$ the scalar hypothesis holds, since
$2\le1+2\log2$; for $a=3/4$ it fails, since $4>1+2\log4$.
Thus it supplies a sufficient restriction on the squared radii, not a
necessary condition for the geometric-mean inequality.

The proof uses point-dependent logarithmic bounds. Write
$H_i=m^{-1}\sum_j\log|1-\overline c_i c_j|$ and
$D_i=-\log(1-|c_i|^2)$. The preceding logarithmic calculation gives
$\sum_iH_i\le0$ and $H_i^2\le-(D_i/m)\sum_jH_j$.
For nonnegative upper bounds $M_i\ge H_i$, Taylor's formula with
integral remainder gives
\[
 e^{H_i}\le1+H_i+\Phi(M_i)H_i^2,
 \qquad \Phi(t)=\int_0^1(1-s)e^{st}\,ds.
\]
Since $\Phi(M_i)\ge0$, summing and using the quadratic bounds yields
\[
 \sum_i e^{H_i}
 \le m+\left(1-\frac1m\sum_i\Phi(M_i)D_i\right)\sum_jH_j
 \le m
\]
whenever $\sum_i\Phi(M_i)D_i\le m$. This is the sufficient condition
in the
\href{https://github.com/wcook04/plectis-erdos/blob/0d34630e1cc9d2b1ac6edfa7cfdba83b31bdc8fe/ErdosProblems/Erdos1041/LogKernelCentralCertificate.lean#L99}{formal inequality with individual upper bounds}.
Taking $M_i=L$ gives $\Phi(L)L\le1$ as a sufficient condition,
since $H_i,D_i\le L$. For $L>0$, integration gives
$\Phi(L)=(e^L-1-L)/L^2$; hence this condition is exactly
$e^L\le1+2L$. At $L=0$ all centres vanish, so the conclusion is
immediate.

\paragraph{Four points without a matching or numerical split.}
The preceding weighted proof with $m=4$ and $w_j=1/4$ gives
\[
 \sum_{j=1}^4\left(\prod_{k=1}^4|1-\overline{c_k}c_j|\right)^{1/4}\le4,
 \qquad |c_j|\le1,
\]
with equality if and only if all four centres vanish. The checked
four-point theorem includes this equality case. In particular, the
inequality is strict as soon as one centre is nonzero, even with repeated
centres or centres on the unit circle.

This four-point consequence needs neither the former $21/25$ split
nor a perfect-matching H\"older step. The scalar stationary-point
inequalities from that approach are not a proof of its missing steps.
The two restricted arguments above remain independent alternatives.
The recorded compilation and axiom checks concern the weighted proof
and its three- and four-point consequences.

The mean inequality gives $\mu\le R^n$. In the unit-disc normalisation
$R=1$, this does not force the smaller quintic threshold $\mu<1/M_5$
below. For $f(z)=z^n-1$, every critical value has modulus one, so no
uniformly smaller bound holds on the closed-disc class. These are bounds
for moduli, not inverse-ray lengths; a path selection or a metric
comparison is still needed.

\subsection*{A sufficient critical-value threshold in degree five}

\paragraph{The distance bound for the two nearest roots.}
\hypertarget{two-nearest-root-distance}{}
Let $f$ have degree $n\ge2$ and all roots in the closed unit disc, and
let $c$ be a critical point with $f(c)\ne0$. List roots with multiplicity.
Its two nearest roots have
total distance at most $2$. The root-disc hypothesis, rather than a bound
on $|f(c)|$, is what gives this estimate.

Rotate so that $c=t\ge0$, and order the distances
as $0<d_1\le d_2\le\cdots\le d_n$. Logarithmic differentiation gives
$\sum_j(z_j-c)^{-1}=-f'(c)/f(c)=0$. Dividing
$|z_j-c|^2+2t\Re(z_j-c)\le1-t^2$ by $d_j^2$, then summing and using
the reciprocal balance, gives
\[
 n\le(1-t^2)\sum_jd_j^{-2}.
\]
Thus $t<1$ and $d_1\le\sqrt{1-t^2}\le1$. The same reciprocal balance,
with the nearest term isolated, gives $d_2\le(n-1)d_1$.

Suppose $d_1+d_2>2$. Since $d_2\le1+t$, we have $d_1>1-t$, and hence
\[
 1-t^2<d_1(2-d_1)<d_1d_2.
\]
The first strict inequality uses that $x(2-x)$ is increasing on $[0,1]$.
It follows that
\[
 n\le(1-t^2)\bigl(d_1^{-2}+(n-1)d_2^{-2}\bigr)
 <\frac{d_2}{d_1}+(n-1)\frac{d_1}{d_2}\le n.
\]
For the last step, put $x=d_2/d_1\in[1,n-1]$ and expand
$(x-1)(x-(n-1))\le0$. This contradiction proves the closed-disc bound.
For roots in the open unit disc, scaling by a containing radius
$0<R<1$ gives $d_1+d_2\le2R<2$.
This is the argument in
\href{https://github.com/wcook04/plectis-erdos/blob/a729f05c40398663fd586c9da487ad874d898639/research_corpus/Erdos1041/GlobalCriticalTwoNearestBudget.md}{the proof for the two nearest roots}; Lean now checks the \href{https://github.com/wcook04/plectis-erdos/blob/52d6c45ad203ba619cb5fe6ba485c0b5400519ea/lean/ErdosProblems/Erdos1041/PaperCompleteR20/TwoNearestPolynomial.lean\#L31}{strict distance bound from the critical polynomial}, including the root-disc reduction and the choice of the two nearest indices.

This argument selects the two nearest roots but does not put their
segments inside a polynomial sublevel set. The next calculation supplies
that separate containment estimate under a bound on $|f(c)|$.

\paragraph{Containment of the two selected segments.}
Let $f$ be monic of degree $n\ge3$ with roots in the closed unit disc.
The following estimate controls the segments from a critical point to
its two nearest roots. To specify the constant, put
\[
\begin{split}
 M_n^2=\max_{\substack{0\le t\le1\\4-n\le y\le n}}
 &(1-t)^2\bigl((1-t)^2+2ty+n(n-2)t^2\bigr)\\
 &\hspace{8mm}\cdot
 \left((1-t)^2+\frac{2t(n-y)}{n-2}\right)^{n-2}.
\end{split}
\]
Here $t$ parametrises the segment. The product estimates below use
$f'(c)=0$ and the ordering of distances, not the root-disc hypothesis;
the latter enters when bounding the combined segment length. To explain
the other parameter and the coefficient $n(n-2)$, suppose $f(c)\ne0$, order the roots
$a_1,\ldots,a_n$ by distance from $c$, and set
$v_j=(a_2-c)/(a_j-c)$. Then $v_2=1$, $|v_j|\le1$ for $j\ge2$,
and $\sum_jv_j=0$ by $f'(c)=0$. With $y=1-\Re v_1$, this gives
\[
 y=2+\sum_{j=3}^n\Re v_j\in[4-n,n],\qquad |v_1|\le n-1.
\]
Consequently
\[
\begin{split}
 |1-tv_1|^2&\le(1-t)^2+2ty+n(n-2)t^2,\\
 \frac1{n-2}\sum_{j=3}^n|1-tv_j|^2
 &\le(1-t)^2+\frac{2t(n-y)}{n-2}.
\end{split}
\]
The factor $v_2=1$ contributes $(1-t)^2$. Applying AM--GM to the
remaining $n-2$ squared factors in
$f(c+t(a_2-c))/f(c)=\prod_j(1-tv_j)$ gives the defining expression
for $M_n^2$. Thus the maximisation is an explicit bound for the second
segment, not an additional optimisation hypothesis about the roots.

For the nearer segment, put $x_j=(a_1-c)/(a_j-c)$. Now $x_1=1$,
$|x_j|\le1$ for every $j$, and $\sum_jx_j=0$. Hence
\[
 \sum_{j=2}^n|1-tx_j|^2
 \le (n-1)(1+t^2)+2t,
\]
so AM--GM gives
\[
 \left|\frac{f(c+t(a_1-c))}{f(c)}\right|^2
 \le(1-t)^2\left((1-t)^2+\frac{2tn}{n-1}\right)^{n-1}.
\]
This bound is at most $M_n^2$ for every $n\ge3$, not just for
numerically tested degrees. Indeed, set $y=n/(n-1)$ in the defining
maximum. This value is admissible because
$n/(n-1)-(4-n)=(n-2)^2/(n-1)\ge0$. The last factor then has base
$(1-t)^2+2tn/(n-1)$, and the preceding factor has this same base plus
$n(n-2)t^2\ge0$. Their product therefore dominates the displayed
nearer-segment bound. Thus $|f(z)|\le M_n|f(c)|$ on both segments, so
$0<|f(c)|\le1/M_n$ gives containment in $\{|f|\le1\}$; the
critical-point distance estimate gives total length at most $2$.
If $f(c)=0$, two root occurrences already coincide at $c$.
The quadratic case is the direct root segment treated earlier.
For degree five the maximum evaluates to
\[
 M_5=(1-t_*)(1+t_*)^3\sqrt{16t_*^2-4t_*+1},
 \qquad
 t_*=\frac{5}{16}+\frac{3\sqrt{105}}{80},
 \qquad
 \frac{1}{M_5}=0.2760461\ldots.
\]
The breakpoint comes from the constraint on $y$, not from a numerical
search. For fixed $0<t<1$, write the two bases in the definition of
$M_5$ as
\[
 A=(1-t)^2+2ty+15t^2,\qquad
 B=(1-t)^2+\frac{2t(5-y)}3.
\]
They are positive for $-1\le y\le5$, and
$\partial_y\log(AB^3)=2t(B-A)/(AB)$. Thus the unconstrained maximum
occurs at $A=B$, or $y=5/4-45t/8$. This reaches the lower endpoint
$-1$ at $t=2/5$; thereafter the maximum is at $y=-1$.
For $0\le t\le2/5$ the resulting expression is
$(1-t)(1+t/2+19t^2/4)^2$, whose derivative is
\[
 \frac{5t(14-19t)(19t^2+2t+4)}{16}\ge0.
\]
For $2/5\le t\le1$, the square of the resulting expression is
$(1-t)^2(1+t)^6(16t^2-4t+1)$, whose derivative is
\[
 -4t(1-t)(1+t)^5(40t^2-25t-2).
\]
The quadratic changes sign once on this interval, at $t_*$, so the
second branch increases from the common breakpoint to $t_*$ and then
decreases to zero. Since the first branch increases to that breakpoint,
$t_*$ gives the global maximum. This proves the displayed value without
sampling either optimisation variable.
As an existence criterion this is weaker than the $13/25$ theorem.
Its different content is the control of the two straight segments from
the specified critical point, rather than a curve selected by area
averaging. For $z^5-b$ it requires $|b|\le1/M_5$, so it excludes part
of a binomial family already covered by the trinomial theorem. No claim
is made that every quintic has a critical value below this level.

A near-regular pentagon illustrates why a prescribed joining point is
restrictive. For sufficiently small $\epsilon>0$, set
\[
 P_\epsilon(z)=z^5+\epsilon^5z^4+\epsilon^3z^3
              -\epsilon^3z^2-\epsilon^5z-1,
 \quad F_\epsilon(z)=r^5P_\epsilon(z/r),\quad r=1-\epsilon^9.
\]
All five roots of $F_\epsilon$ lie in the open unit disc, approach a
regular pentagon, and have escaping segments from the origin. Here is a
proof of that assertion, including the effect of the radial contraction.

The roots of $P_\epsilon$ are simple perturbations
$\zeta_\epsilon=\zeta_0+O(\epsilon^3)$ of the fifth roots of unity.
The identity
\[
 z^5\overline{P_\epsilon(1/\overline z)}=-P_\epsilon(z)
\]
makes the roots invariant under conjugate inversion. Uniqueness of the
root near each $\zeta_0$ therefore forces $|\zeta_\epsilon|=1$.
The roots $r\zeta_\epsilon$ of $F_\epsilon$ consequently have modulus
$r<1$. The displacement estimate follows from the implicit function
theorem at each simple root of $z^5-1$, since the first coefficient
perturbations have order $\epsilon^3$.

For $\zeta_0=1,e^{4\pi i/5},e^{-4\pi i/5}$, take $t=\epsilon/4$.
Substitution into the polynomial gives
\[
 |F_\epsilon(rt\zeta_\epsilon)|^2
 =1+\left(\frac{\Re\zeta_0^2}{8}-\frac1{512}\right)\epsilon^5
       +O(\epsilon^6).
\]
The coefficient is positive, since
$\Re\zeta_0^2\ge(\sqrt5-1)/4$.
For the remaining roots, with $\zeta_0=e^{2\pi i/5}$ or
$e^{-2\pi i/5}$, take $t=\epsilon^2/4$ instead. Then
\[
 |F_\epsilon(rt\zeta_\epsilon)|^2
 =1+\frac{3\sqrt5-5}{32}\epsilon^7+O(\epsilon^9)>1
\]
for sufficiently small $\epsilon$. The two scales are needed because
$\Re\zeta_0^2$ is negative for this second pair; its positive linear
contribution dominates at the smaller scale. Multiplication by
$r^{10}=1+O(\epsilon^9)$ cannot remove either positive leading term.
The remainders are uniform over the five root branches, so one sufficiently
small $\epsilon$ works for every segment. This is an ordinary asymptotic
proof; it gives no numerical upper bound for the permitted $\epsilon$.

The addendum \cite[Proposition~1]{revision2026} reports the stronger
assertion that no critical point of this family supplies two contained
straight arms of total length at most $2$. That proof is not supplied
with this record and was not available for checking in this revision.
The origin calculation above does not prove it: $P_\epsilon'(0)=-\epsilon^5$
shows that the origin is not a critical point. The stronger reported
assertion is not used in any theorem here. Neither the proved failure at
the origin nor the reported critical-point obstruction excludes curved
inverse images of rays or a different joining point.

A different local question allows a freely chosen, noncritical vertex:
for some $\delta>0$, does every monic quintic with open-disc roots within
$\delta$ of a rotated regular pentagon admit two contained straight arms
through some $h\in\mathbb C$ of total length below $2$?
Neither the origin nor a critical point is prescribed in this question.
The sufficient low-critical threshold preceding this paragraph is unchanged.

\paragraph{Area of a connected sublevel set.}
The following calculation applies in every degree $n\ge2$, independently
of the degree-five segment criterion. Let $f$ be monic of degree $n$ and let $t>0$ with
$K_t=\{|f|\le t\}$ connected. Translate its root centroid to $0$, so
$f(z)=z^n+c_{n-2}z^{n-2}+\cdots$, and let $\psi$ be the exterior map
normalised at infinity with positive leading coefficient. For a regular
connected lemniscate, the identity $f(\psi(\zeta))=t\zeta^n$ is
\cite[(2.3)]{eks2010} after scaling the value level. The following argument
also permits critical boundary points. The exterior Green functions
$(\log|f|-\log t)/n$ and $\log|\psi^{-1}|$ agree: both vanish at the
boundary and have a logarithmic pole of coefficient one at infinity.
Comparison there gives the leading coefficient $t^{1/n}$ for $\psi$.
The quotient $f(\psi(\zeta))/(t\zeta^n)$ has constant modulus one and
limit one at infinity, proving the polynomial identity. Its
$\zeta^{n-1}$ coefficient forces the constant Laurent coefficient of
$\psi$ to vanish. Thus
\[
 \psi(\zeta)=t^{1/n}\zeta+\sum_{k\ge1}a_k\zeta^{-k}.
\]

Gr\"onwall's area identity takes the form
\[
 \operatorname{Area}(K_t)=\pi\left(t^{2/n}-
 \sum_{k\ge1}k|a_k|^2\right).
\]
One can obtain it directly without assuming that level $t$ is regular.
For $r>1$, Green's area formula on the analytic Jordan curve
$\psi(re^{i\theta})$ gives the enclosed area
$\pi(t^{2/n}r^2-\sum_{k\ge1}k|a_k|^2r^{-2k})$.
These enclosed compact sets decrease to $K_t$ as $r\downarrow1$;
continuity of area from above and monotone convergence of the series
give the identity. Comparing the coefficient of $\zeta^{n-2}$ in the
exterior identity now gives
\[
 nt^{(n-1)/n}a_1+c_{n-2}t^{(n-2)/n}=0,
 \qquad a_1=-\frac{c_{n-2}}{n\,t^{1/n}}.
\]
The nonnegative coefficient sum quantifies the loss from P\'olya's
upper bound. It does not by itself control an internal path length.
No analogous coefficient formula for an individual component of a
disconnected sublevel set is established here. The
degree-five statements and their evidence classes are recorded in
\href{https://github.com/wcook04/plectis-erdos/blob/f214a6b45528dc5eefe20ffadc35f2e981627d4c/research_corpus/Erdos1041/Degree5AssemblyAndSharpenedCuts.md}{the degree-five estimates and their proofs}.

\subsection*{Numerical tests of selection by the smallest critical value}

Three examples distinguish critical-value bounds from path-length bounds.
For a critical point $c$ whose two inverse arms over $[0,f(c)]$ reach
roots, write $L_f(c)$ for their total length.
The first computation tests the rule of choosing the critical point
with uniquely smallest value modulus. The source gives an unnormalised
quintic $f_0$, a selected critical point $c_0$, and the numerical values
\[
 L_{f_0}(c_0)\approx2.057343275393654508,
 \qquad R\approx1.021393477405696164,
\]
where $R$ is the radius of a smallest disc containing the roots of $f_0$.
The first number is not the length after normalisation to the unit disc.
If $h$ is the centre of that disc, set
\[
 F(z)=R^{-5}f_0(h+Rz).
\]
This preserves monicity, sends the roots into the closed unit disc,
and divides every corresponding curve length by $R$. The reported
normalised values are therefore
\[
 L_F\!\left(\frac{c_0-h}{R}\right)\approx2.01425143287505,
 \qquad
 \left|F\!\left(\frac{c_0-h}{R}\right)\right|
 \approx0.899569245.
\]
A further contraction uses $F_s(z)=s^5F(z/s)$, with $0<s<1$: roots and
curve lengths are multiplied by $s$, and critical values by $s^5$.
Thus the reported excess would persist for $s$ sufficiently close to $1$,
with roots in the open unit disc.

These are numerical observations, not a certified counterexample. A small
residual at a computed fibre root does not enclose that root or establish
its membership in the tracked inverse branch. The inscribed-polyline
comparison needs both facts before it gives a rigorous lower bound for
the intended arms. Another critical point supplies a shorter pair in the
reported computation; this is evidence at that configuration, not a
universal selection theorem.

Second, the linked record reports numerical violations of
$\sum_cL_f(c)\le2(n-1)R$ in degrees four and five. Here $R$ is the radius of
a smallest disc containing the roots. The quartic computation gives
$\sum_cL_f(c)=6.006352157\ldots$ against the proposed bound $6$.
Several differential-equation solvers agree, and an inscribed-polyline
calculation using numerically computed fibre roots gives a similar excess.
These checks do not provide rigorous enclosures for those roots, their
branch assignment or the resulting length lower bound. They are evidence
against the aggregate inequality, not a certified refutation or a proof
that an open set of polynomials violates it. Its algebraic factor
\[
 \sum_{k=1}^{n-1}|f(c_k)|^{1/n}\le(n-1)R
\]
is proved for every $n\ge2$ in Theorem~\ref{res:critical-value-budget}.
The value estimate alone supplies no bound for the associated path-length
sum. The algebraic theorem therefore remains useful independently of how
the numerical aggregate examples are resolved.
The two computational records, to be read with the evidence limitations
just stated, are
\href{https://github.com/wcook04/plectis-erdos/blob/f214a6b45528dc5eefe20ffadc35f2e981627d4c/research_corpus/Erdos1041/MinimalHubArmBudgetRefutation.md}{the numerical test of selection by the smallest critical value}
and
\href{https://github.com/wcook04/plectis-erdos/blob/f214a6b45528dc5eefe20ffadc35f2e981627d4c/research_corpus/Erdos1041/SeparatrixAggregateReduction.md}{the comparison between sums of critical values and path lengths}.

Third, all five origin segments can escape even when the roots are
arbitrarily close to a regular pentagon. The two-scale calculation above
proves this for the family $F_\epsilon$ attributed to the earlier note
\cite[Proposition~A of the preserved earlier note]{revision2026}; an
isolated numerical example would not prove the assertion for arbitrarily
small defects. Failure for every critical joining point is a different,
stronger claim, reported in the addendum but not established by the
origin calculation. No conclusion about curved inverse-ray paths or an
arbitrary connector follows from either straight-segment assertion.

\subsection*{Limits of contained paths at fixed degree}
\hypertarget{closed-class-limits}{}
The compactness argument used in Section~17 of the short note is given
here in full. It concerns arbitrary contained curves, not a prescribed
pair of inverse rays.

Fix $n\ge2$, and let $\mathcal K_n$ be the compact coefficient class of
monic degree-$n$ polynomials with roots in the closed unit disc.  Define
\[
 \Lambda(f)=\inf_\gamma\operatorname{length}(\gamma),
\]
where $\gamma\subset\{|f|\le1\}$ joins two different listed root
occurrences. A repeated root permits a constant curve; an empty set of
competitors has infimum $+\infty$.

\paragraph{Passing a length bound to a limit.}
Let $f_j\to f$ in $\mathcal K_n$. If the lower limit of
$\Lambda(f_j)$ is finite, pass to a subsequence on which these infima
converge to it and choose rectifiable curves $\gamma_j$ with lengths at
most $\Lambda(f_j)+1/j$. Every closed unit sublevel lies in $|z|\le2$:
for $|z|>2$, the root factorization gives
$|f_j(z)|\ge(|z|-1)^n>1$. Constant-speed parametrisation on $[0,1]$
therefore gives uniformly bounded images and Lipschitz constants.
Arzel\`a--Ascoli supplies a uniformly convergent subsequence.

To retain the endpoints with their multiplicities, list the two chosen
occurrences first among the roots of $f_j$, and list the remaining roots
arbitrarily. Compactness of the product of $n$ closed unit discs gives a
further subsequence on which this entire list converges. The limiting
factorization is $f$, so its first two entries remain two occurrences,
even if their locations coincide. Uniform convergence of the curves and
of the polynomials on $|z|\le2$ gives $|f|\le1$ on the limiting curve
$\gamma$. For any partition $0=t_0<\cdots<t_M=1$,
\[
 \sum_{\ell=1}^M|\gamma(t_\ell)-\gamma(t_{\ell-1})|
 =\lim_j\sum_{\ell=1}^M|\gamma_j(t_\ell)-\gamma_j(t_{\ell-1})|
 \le\liminf_j\operatorname{length}(\gamma_j).
\]
Taking the supremum over partitions proves the required lower
semicontinuity of length and hence
$\Lambda(f)\le\liminf_j\Lambda(f_j)$. An infinite lower limit needs no
argument. With $f$ fixed, the same compactness proof shows that a finite
infimum is attained. Thus a bound $\Lambda\le2$ on a dense class gives
an actual path of length at most $2$ throughout $\mathcal K_n$.
For an open-disc polynomial, choose a containing disc of radius
$0<R<1$ and scale: the resulting path has length at most $2R<2$ and
lies in $\{|f|\le R^n\}\subset\{|f|<1\}$.

For $f(z)=z^n-1$, the inequality $|z^n-1|\le1$ implies
$2\Re(z^n)\ge|z|^{2n}>0$ whenever $z\ne0$. Thus the punctured sublevel
lies in $n$ disjoint angular sectors, one per root. A path between
different roots must pass through zero and has length at least $2$ by
the triangle inequality. The two radial segments attain that length.  The corresponding historical assertion on the closed class is
\[
 \boxed{\Lambda(f)\le\Lambda(z^n-1)=2\qquad(f\in\mathcal K_n).}
\]
The lower-semicontinuity argument is proved in
\href{https://github.com/wcook04/plectis-erdos/blob/f36a98bf3d3e6f65f1074e3b800e3293d5b8a51a/ErdosProblems/Erdos1041/GenericSufficiencyClosure.md}{the compactness proof}.

The displayed universal inequality is recorded to explain the earlier
reduction, not proposed here as a remaining open assertion. A counterexample
to the original problem would also refute this inequality and, by the same
compactness argument, any dense-class estimate implying it. The reduction
can still be used for separately specified polynomial classes.

\subsection*{Blaschke-product examples and limits in varying degree}
\hypertarget{blaschke-product-examples}{}
The following family separates the limitations of the sufficient
conditions from bounds for the actual shortest contained path.

The ordinary
\href{https://github.com/wcook04/plectis-erdos/blob/f36a98bf3d3e6f65f1074e3b800e3293d5b8a51a/ErdosProblems/Erdos1041/BlaschkePowerCriticalSpectra.md}{note on powers of Blaschke products}
uses the degree-$2N$ polynomials
\[
 F_N(z)=[z(z+b)]^N-(1+bz)^N,\qquad N\ge1,\quad 0<b<1,
\]
and their radial contractions. Its finite rational certificates and
its asymptotic proofs have different roles. The degree-eight certificate
excludes an unconditional critical-value exponent $8/7$. It does not
contradict the ordinary exponent $4/7$ proved by the Poisson argument above,
or the smaller exponent $2/7$ in the cited Lean theorem.
For $N\ge2$, the root centroid of $F_N$ is $-b/2$, as its
$z^{2N-1}$ coefficient is $Nb$. A radial contraction by $r>0$ changes
this to $-rb/2$, not zero. Thus the degree-eight example does not satisfy
the additional centroid condition that permits numerator $8$ above. The degree-twenty-four certificate gives critical values in
$|v+1|<1/12$, so $\mu>11/12$ while every radius-$4/3$ separation test
fails. These fixed-degree examples do not establish optimality uniformly
in degree. That conclusion uses the limiting measures in Sections~2--3
of the linked full proof.

For $b=\lambda/N$ with fixed $0<\lambda\le1$, Section~4 constructs
connectors of length at most
$[\log(2/\lambda)+\pi+o(1)]/N$. The two radial pieces reduce the
oscillatory term; a short circular arc then joins them in a sector where
the remaining factor has modulus below one. The unspecified starting
degree depends on $\lambda$. For $b=e^{-\alpha N}/N$ with fixed
$\alpha>0$, the same section instead proves
$\Lambda(F_N)\to2(1-e^{-\alpha/2})$ for the uncontracted polynomials
$F_N$ displayed above and the closed-level functional defined earlier.
The second assertion is a separate argument, not a substitution of a
varying parameter into the first. Both are ordinary asymptotic arguments,
not consequences of the finite certificates or formally checked analytic
theorems. They concern varying degree and do not contradict fixed-degree
compactness or prove the unrestricted assertion.

\subsection*{What the earlier inverse-ray approach would have required}

The earlier proposed sufficient condition concerned \emph{curved}
inverse images of rays, not straight segments.  On a generic ray-separated class let $L_f(c)$ be
the length of the two-root descent arc at an admissible critical point,
where $|f(c)|\le1$.  Then
\[
 \min_{c\ \mathrm{admissible}}L_f(c)\le2
\]
would imply the root-connector conclusion by the stated compactness
argument.  The unrestricted infimum over all contained curves can be
smaller; equality of these two optimisation problems is not assumed.
The obstruction to two straight segments through a critical point does
not by itself prove or disprove this curved-arc estimate.  Examples in which another critical point gives a shorter curve do not
prove the estimate universally. The critical value bound, containment and
length bound must hold for the same curve. The displayed universal
inequality is recorded as a historical sufficient condition, not proposed
as an open assertion: a counterexample to the original question would
refute it as well. The reduction remains useful on separately specified
polynomial classes.

\section{The Newton value equation}\label{sec:newton}

Sutherland uses the term Newton flow for $\dot z=-f(z)/f'(z)$ and
notes that its solution curves map under $f$ to radial lines
\cite[p.~42]{sutherland1992}. The value equation
\cite[Lemma~2.1]{kozen-stefansson1997} determines their orientation and
time parametrisation. The global Newtonian graph
\cite[\S2]{kozen-stefansson1997}, described with its endpoint conventions
below, is a separate input to the inverse-sheet discussion in
Section~\ref{sec:gap}.

A Newton trajectory is a differentiable curve satisfying
\[
 z'(t)=-\frac{f(z(t))}{f'(z(t))}
\]
where $f'(z(t))\ne0$. Let $I\subseteq\mathbb R$ be an interval on
which these assumptions hold;
put $w(t)=f(z(t))$.  Kozen and Stef\'ansson record the following identity as
a lemma of Shub, Tischler and Williams \cite[Lemma~2.1]{kozen-stefansson1997}.

\begin{theorem}[value equation]\label{res:value}
Let $f$ be a polynomial and $z:I\to\mathbb C$ a differentiable curve on
an interval $I$, with $f'(z(t))\ne0$ and
$z'(t)=-f(z(t))/f'(z(t))$ throughout $I$.  For $w=f\circ z$, one has
$w'(t)=-w(t)$ on $I$.
\end{theorem}
\leannote{Lean: \lproof{ErdosProblems/Erdos1041/PaperCompleteR20/NewtonRealTime.lean}{52}{newton_real_value_whole}.}

\noindent The computation is one line: $w'=f'(z)\,z'=f'(z)\cdot(-f(z)/f'(z))=-f(z)=-w$.
The kernel checks the local complex-parameter chain rule as
\lword{Erdos1041/NewtonFlowRaySeparation.lean}{50}{newtonFlow_value_hasDerivAt}{the derivative of the polynomial value along a Newton trajectory}, together with the
differential form of the first integral,
\lword{Erdos1041/NewtonFlowRaySeparation.lean}{64}{newtonFlow_scaledValue_hasDerivAt_zero}{the derivative of the exponentially rescaled value}:
\[
  \frac{d}{dt}\Bigl(e^{t}f(z(t))\Bigr)=0.
\]
Equivalently,
\[
  f(z(t))=e^{-(t-t_0)}f(z(t_0)),\qquad t_0\le t,\quad t_0,t\in I .
\]
Integration on the real interval gives the last identity as an ordinary
consequence of the differential equation.
For an existing trajectory with nonzero initial value, its value moves
inward on one positive ray; a zero value has no argument and remains zero.
Thus $|f|<1$ is preserved for later times in that trajectory's existing
interval. No global existence or description of a whole ray preimage follows
from this scalar equation. The real-time candidate endpoint is
\href{https://github.com/wcook04/plectis-erdos/blob/f36a98bf3d3e6f65f1074e3b800e3293d5b8a51a/ErdosProblems/Erdos1041/NewtonFlowTrajectory.lean}{the real-time value equation} in \href{https://github.com/wcook04/plectis-erdos/blob/f36a98bf3d3e6f65f1074e3b800e3293d5b8a51a/ErdosProblems/Erdos1041/NewtonFlowTrajectory.lean}{the real-time trajectory source}.

\begin{corollary}[ray separation]\label{res:ray}
Let $a<b$ and let the value trajectory $t\mapsto f(z(t))$ be continuous on
$[a,b]$. Assume the Newton equation and $f'(z(t))\ne0$ on $(a,b)$ only.
Then
\[
 f(z(b))=e^{a-b}f(z(a)).
\]
If these endpoint values are nonzero, they lie on one positive ray. Therefore
critical points with values on distinct positive rays cannot be endpoints
of such a finite connection. The trajectory in the $z$ plane need not be radial.
\end{corollary}
\leannote{Lean: \lproof{ErdosProblems/Erdos1041/PaperCompleteR20/NewtonRealTime.lean}{69}{newton_real_endpoint_whole}.}

\noindent The interior identity passes to the endpoints by continuity, not
by evaluating $-f/f'$ at a critical endpoint. The candidate declarations
\href{https://github.com/wcook04/plectis-erdos/blob/f36a98bf3d3e6f65f1074e3b800e3293d5b8a51a/ErdosProblems/Erdos1041/PaperNewtonEndpoints.lean}{value decay at continuous endpoints} and
\href{https://github.com/wcook04/plectis-erdos/blob/f36a98bf3d3e6f65f1074e3b800e3293d5b8a51a/ErdosProblems/Erdos1041/PaperNewtonEndpoints.lean}{no finite connection of distinct value rays}
in \href{https://github.com/wcook04/plectis-erdos/blob/f36a98bf3d3e6f65f1074e3b800e3293d5b8a51a/ErdosProblems/Erdos1041/PaperNewtonEndpoints.lean}{the trajectory endpoint source} make these
premises explicit. Their compilation and axiom checks are not reported here.
They construct no trajectories and supply no global monodromy theorem.
Kozen and Stef\'ansson draw the same ray consequence for the Newtonian graph:
under $f$, every edge maps onto a segment of a ray through the origin whose
endpoints are $0$ or critical values \cite[\S2]{kozen-stefansson1997}.

\section{Arguments, not moduli}\label{sec:arguments}

It is tempting to arrange a generic perturbation so that the critical values
are pairwise distinct, or that their moduli are pairwise distinct, and to
conclude that saddle connections are excluded.  Neither is enough.

Along a Newton trajectory the argument of $f$ stays constant while its
modulus decreases. By Corollary~\ref{res:ray}, arranging pairwise distinct
arguments of the nonzero critical values is therefore a sufficient way
to exclude finite connections between critical points. It is not asserted
to be necessary. This condition imposes no lower bound on the distances
between critical values and no upper bound on their moduli.

The weaker conditions really can coexist with a connection, even when
all roots lie in the open unit disc. Take
\[
 f(z)=z^3-\frac3{16}z+\frac3{64}.
\]
On $|z|=1$, the lower terms have total modulus at most $15/64<1$, so
Rouch\'e's theorem places all three roots inside. The critical points
are $-1/4$ and $1/4$, with distinct positive values $5/64$ and $1/64$.
On the intervening real interval $f$ is positive and strictly decreasing.
The equation
\[
 f(x(t))=\frac5{64}e^{-t},\qquad 0<t<\log5,
\]
therefore defines a differentiable path with
$x'(t)=-f(x(t))/f'(x(t))$. Its endpoint limits are $-1/4$ and $1/4$.
It is an actual critical-to-critical Newton trajectory, not merely a
pair of values on one ray. The time convention is that of the value
equation \cite[Lemma~2.1]{kozen-stefansson1997}; the singular field is
not evaluated at the endpoints.

A common translation of distinct values can separate their arguments.
The next calculation describes the translations to avoid.

\begin{theorem}[ray-collision locus]\label{res:locus}
Let $a\ne b$ be complex.  Every common translation $\beta$ for which $a+\beta$
and $b+\beta$ lie on the same positive ray has the form
\[
  \beta=\frac{ra-b}{1-r},
  \qquad r\in\mathbb{R}_{>0},\ r\ne1 .
\]
\end{theorem}
\leannote{Lean: \lproof{ErdosProblems/Erdos1041/NewtonFlowRaySeparation.lean}{107}{translated_samePositiveRay_parameterization}.}

Indeed, write $b+\beta=r(a+\beta)$ with $r>0$. Since $a\ne b$,
one has $r\ne1$, and solving for $\beta$ gives the displayed formula.
Equivalently, $\beta=-a+(a-b)/(1-r)$, so the forbidden translations
lie on the real affine line through $-a$ and $-b$.

\noindent The formula is checked as
\lword{Erdos1041/NewtonFlowRaySeparation.lean}{107}{translated_samePositiveRay_parameterization}{translations that put two values on one ray}. A finite union of these loci has empty interior. Adding a constant to
$f$ translates all critical values by that constant and leaves the critical
points unchanged. Section~\ref{sec:open} records the formal avoidance and
root-retention estimates. A constant cannot separate equal critical values;
coefficient perturbation and control of the remaining geometric margins
are separate requirements.

\section{Why the proposed spanning-tree estimate fails}\label{sec:gap}

Recent work on polynomial lemniscates separates component counts from metric
path questions. Ghosh and Ramachandran record that the open set
$\{|f|<1\}$ has $1+\#\{j:|f(c_j)|\ge1\}$ components, where
$c_1,\ldots,c_{n-1}$ are the critical points listed with multiplicity
\cite[Lemma~2.5]{ghosh2023}; the underlying component-wise Riemann--Hurwitz
count appears in the proof of \cite[Proposition~2.1]{eks2010}.  For the
binomial family $z^n-a$, the condition $|a|<1$ therefore puts the filled unit
lemniscate in the connected regime. Connectedness alone gives no path-length bound.
Bishop, Eremenko and Lazebnik describe the possible shapes: every rational
lemniscate is a lemniscate graph whose vertices are the critical points on the
level set, each of even degree at least four
\cite[Definition~1.2 and Proposition~2.4]{bishoperemenkolazebnik2025}, and a
lemniscate graph is realised by a polynomial lemniscate up to a homeomorphism
of the plane exactly when it is the boundary of its unbounded face
\cite[Corollary~1.6]{bishoperemenkolazebnik2025}.  This topological
description also gives no path-length bound.

The March manuscript's Proposition~12 claims the following
estimate.  For $u=-\log|f|$, a connected component $V$ of $\{u>c\}$ carrying
$m\ge2$ simple zeros, and the regularity and Morse hypotheses in that proposition, it
constructs, for every $\varepsilon>0$, an embedded spanning tree
$G_\varepsilon\subset V$ with
\[
 \operatorname{len}(G_\varepsilon)
 \le\frac1{2\pi}\int_{2\alpha}^{\infty}P_V(t)\,dt+\varepsilon.
\tag{5.1}\label{eq:prop12-bound}
\]
Here $P_V(t)=\mathcal H^1(V\cap\{u=t\})$ is the length of the level
set in $V$, and $\alpha>0$ is the truncation parameter in the proposed
estimate. The final theorem depends on this bound.

The estimate itself is false. Take $f_a(z)=z^2-a^2$ with $a=9/10$
and choose $2\alpha=-\log a$. For $0<c<\alpha/2$, the component
$V$ of $\{-\log|f_a|>c\}$ contains $\{|f_a|\le a\}$ and both roots.
The only critical point is a nondegenerate saddle at $0$; distinctness
of critical values is vacuous.

Here the level-length integral can be bounded directly. Put $r=e^{-t}$.
On a regular level the inverse parametrisations
$z=\pm\sqrt{a^2+re^{i\theta}}$ give
\[
 P_V(t)=r\int_0^{2\pi}|a^2+re^{i\theta}|^{-1/2}\,d\theta.
\]
The elementary identity
\[
 |a^2+re^{i\theta}|^2
 =(a^2+r)^2\cos^2(\theta/2)+(a^2-r)^2\sin^2(\theta/2)
\]
implies
\[
 |a^2+re^{i\theta}|^{-1/2}
 \le(a^2+r)^{-1/2}|\cos(\theta/2)|^{-1/2}.
\]
The angular factor has mean at most $2$. Indeed, symmetry reduces its
integral to $4\int_0^{\pi/2}(\cos s)^{-1/2}\,ds$, and
$\cos s\ge1-2s/\pi$ bounds the latter integral by $4\pi$.
Changing variables from $t$ to $r$ therefore gives
\[
 \frac1{2\pi}\int_{2\alpha}^{\infty}P_V(t)\,dt
 \le2\int_0^a\frac{dr}{\sqrt{a^2+r}}
 =4\bigl(\sqrt{a^2+a}-a\bigr)<\frac{41}{25}.
\]
The isolated critical level is handled by the corresponding improper
integral. Every connected set containing both roots has length at least
$2a=9/5$. Since $9/5-41/25=4/25$, any
$0<\varepsilon<4/25$ contradicts~\eqref{eq:prop12-bound}.  Thus neither a different local saddle neighbourhood nor a perfect
topological decomposition can recover the printed coefficient $1/(2\pi)$.
Lean checks the exact numerical inequality and the resulting contradiction
for a tree-length bound. The level-length majorant remains an explicit
analytic input; the integral evaluation is not checked by that formal proof.

At an interior index-one critical point $p$, the proof invokes a Morse chart
\[
 u=u(p)+x^2-y^2
\]
and replaces the saddle by a three-ended neighbourhood having one connected
lower cross-section and two connected upper cross-sections.  That local model
is false as written.  Because $p\in V$ and $V$ is open, a sufficiently small
closed disc around $p$ lies entirely in $V$.  In that disc the full Morse chart
has four sectors: two components of $u>u(p)$ and two components of $u<u(p)$.
The level set $\{u=u(p)\}$ is a lemniscate graph with a vertex of degree four
at $p$ \cite[Definition~1.2 and Proposition~2.4]{bishoperemenkolazebnik2025}.
A global component argument cannot delete one local sector from a disc already
contained in $V$.

This independently diagnoses a proof step, but the Cassini witness above also
refutes the proposition's printed metric statement.  A different route might
cut an adjoining regular annulus along a separatrix or regular flow arc before
forming the block, retain a four-pronged saddle neighbourhood and change the
assembly, or replace the local construction by the ray-cut decomposition
proposed below.  But no repair can retain~\eqref{eq:prop12-bound}; it must pay a
positive attachment cost, select only one short pair instead of spanning every
root, or use a different global metric inequality.
The shorter descriptions of the same three-ended block do not repair the
four-sector topology.

There is an ordinary topological theorem when the critical values have
pairwise distinct arguments and moduli, and the critical points are simple.  Its
ingredients are classical: the Newtonian graph of Shub, Tischler and Williams
\cite[\S2]{kozen-stefansson1997}, the component-wise Riemann--Hurwitz count in
the proof of \cite[Proposition~2.1]{eks2010}, and the division of the Riemann
surface of the inverse function into sheets by outward critical-value
slits.  Dubinin describes the sheet adjacency tree directly in
\cite[\S1, pp.~1169--1170]{dubinin2006critical} and uses a related
network for capacity in \cite[\S2, pp.~34--35]{dubinin2006capacity}.
The latter paper assumes bounded critical values and controls capacity;
it does not supply a Euclidean tree-length bound.  The theorem below specifies the sheets,
critical transpositions and descent arcs for the stated component.

The assumptions exclude a component containing a multiple critical point
or two critical values on the same ray. For $z^n-b$ with $n>2$ and
$0<|b|<1$, the component containing the origin fails the first condition.
They hold, for example, for
$f(z)=z^3-(3/25)z+1/500$: its critical points are $\pm1/5$ and its
critical values are $-7/500$ and $9/500$. On the unit circle the two
lower terms have total modulus at most $61/500<1$, so all three roots
lie in the open unit disc by Rouch\'e's theorem. In the statement,
``excellent'' means that the critical levels of the Morse function are
pairwise distinct. This regularity permits a decomposition into sheets;
it does not bound the total length of the resulting tree.

\hypertarget{inverse-sheet-component}{}
\begin{theorem}[inverse sheets with distinct critical-value arguments]
\label{res:attachment-aware-reeb}
Let $f$ be monic, and let $U$ be a component of $\{|f|<1\}$ containing
$k\ge2$ roots, counted with multiplicity. Suppose every critical point of
$f$ in $U$ is simple, its critical value is nonzero, and these critical
values have pairwise distinct arguments and pairwise distinct moduli.
All preimages and sheets below are taken inside $U$, and only critical
points in $U$ determine the cuts. Then:
\begin{enumerate}
\item $-\log|f|:U\mathbin{\backslash}f^{-1}(0)\to(0,\infty)$ is a proper excellent
Morse function, with exactly $k-1$ nondegenerate saddles;
\item cutting $\mathbb D\smallsetminus\{0\}$ along the critical-value rays
decomposes its preimage in $U$ into conformal strips;
\item cutting each ray only from its critical value to the unit circle gives
$k$ conformal sheets, one per root, whose critical transpositions form a tree;
\item for each critical point $c\in U$, the two inverse lifts of
$[0,f(c)]$ join two roots through $c$ inside
$U\cap\{|f|\le|f(c)|\}$, and the union of these arcs is an embedded
geometric realisation of that tree.
\end{enumerate}
Small neighbourhoods of the saddles can be chosen with diameter
$O(\sqrt\delta)$ at value radius $\delta$.
\end{theorem}
\leannote{Not formalised: the ray-disjointness, level-separation and saddle-scale steps of the proof are checked, and the Morse, monodromy and strip statements are not. See the \hyperref[sec:coverage]{coverage section}.}

\begin{proof}
The map $f:U\to\mathbb D$ is proper of degree $k$, and $U$ is simply
connected. To apply the count from the proof of
\cite[Proposition~2.1]{eks2010} only at a regular level, join all roots and
critical points in $U$ by finitely many compact paths in $U$. Choose a
regular $t<1$ above the maximum of $|f|$ on their union. One component of
$\{|f|<t\}$ then contains exactly the roots and critical points of $U$;
Riemann--Hurwitz there gives $k-1$ simple critical points. No regularity
of the level-one boundary is needed. The holomorphic Morse lemma makes
these points nondegenerate saddles of $-\log|f|$, and their distinct
moduli separate the critical levels. A compact range in $(0,\infty)$
stays away from both the level-one boundary and the deleted roots, proving
properness of the Morse function.

Remove $0$ and the rays determined by these $k-1$ values. Each remaining
sector is simply connected and contains no branch value of
$f:U\to\mathbb D$, even if it contains critical values coming from outside
$U$. Each component of its preimage in $U$ maps biholomorphically to the
sector. If the sector has angular range $(\theta_1,\theta_2)$, the
holomorphic coordinate $-\log f$ maps this component onto the
semi-infinite strip
\[
 \{x+iy:x>0,\ -\theta_2<y<-\theta_1\}.
\]
Thus its real coordinate is $-\log|f|$, while its imaginary coordinate
is $-\arg f$. The minus sign in both coordinates is needed for a
holomorphic map; keeping $+\arg f$ would reverse orientation.

For the sheet tree, remove only the outward slits
$\{tf(c):1\le t<1/|f(c)|\}$, for critical points $c\in U$.
The remaining value domain is star-shaped about $0$, hence simply
connected, and its preimage in $U$ splits into $k$ sheets labelled by the
roots. Local monodromy at each slit endpoint is a transposition of two
of these sheets. Deleting the finitely many fibres over the branch values
from $U$ leaves a connected domain, so the monodromy action is transitive.
The $k-1$ transpositions therefore give a connected graph with $k$ vertices,
which is a tree.

For $c\in U$, the inward segment from $f(c)$ to $0$ contains no other
branch value of $f:U\to\mathbb D$, by the distinct-argument hypothesis.
Distinct moduli were needed only to separate the Morse levels. Its two
inverse lifts converge to $c$ by the holomorphic Morse coordinate and
start at the two root labels exchanged by its transposition. These labels
are distinct, as are the roots, because a root in $U$ cannot also be a
critical point under the nonzero-critical-value hypothesis. Lifts belonging to
different critical points cannot cross away from roots, since an intersection
would give one nonzero value on two critical rays. Their union is therefore
an embedded tree. Each edge is rectifiable: at a simple critical endpoint,
parametrising the value segment by its distance $s$ from the critical value
gives speed $O(s^{-1/2})$, which is integrable. The root endpoint is regular,
and the derivative is bounded on compact subsegments between the endpoints.
This proves finiteness of each length, not a uniform bound for it.
The same local Morse coordinate identifies a value disc of
radius $\delta$ with a spatial neighbourhood of diameter
$O(\sqrt\delta)$. The constant and the permitted range of $\delta$
depend on $f$ and the chosen saddle. Indeed, in the coordinate
$f(z)-f(c)=u^2$, the inverse map has derivative of modulus
$\sqrt{2/|f''(c)|}$ at $u=0$. No bound uniform over degenerating
critical points is asserted. For a fixed polynomial there are finitely
many saddles, so their individual constants have a finite maximum.
\end{proof}

The complete ordinary argument is recorded in
\href{https://github.com/wcook04/plectis-erdos/blob/a729f05c40398663fd586c9da487ad874d898639/research_corpus/Erdos1041/AttachmentAwareReeb.md}{the proof of the decomposition into inverse sheets}.
Its Lean companion checks the ray-distance, averaging and finite-tree kernels;
it explicitly does not formalise Riemann--Hurwitz, monodromy, the strip
diffeomorphisms or the geometric assembly.  The theorem supplies topology,
not a useful length sum: the canonical inverse-ray tree can already have the
Cassini deficit discussed above.

The component restriction cannot be dropped. For example, take
\[
 f(z)=(z^2-4)^2-\tfrac14.
\]
Its critical points are $0,\pm2$, with values $63/4,-1/4,-1/4$.
The positive roots $\sqrt{7/2}$ and $\sqrt{9/2}$ lie in one component
$U$: on the intervening real interval, $-1/4\le f\le0$.
But $f(iy)=(y^2+4)^2-1/4>1$ for real $y$, so this component cannot
contain the negative roots or $-2$. Its only critical point is the simple
point $2$. Thus the theorem applies in $U$, although the other component
has the same critical value. This example concerns component topology;
no root-unit-disc hypothesis is imposed here.

Corollary~\ref{res:ray} also excludes finite saddle-to-saddle Newton
connections inside $U$. The theorem does not treat coincident arguments,
simultaneous critical levels or multiple saddles in that component. Nor
does it identify a slit domain with the quotient by all trajectories:
the quadratic example following Problem~\ref{prob:globalflow1041} satisfies
its hypotheses but has a non-Hausdorff orbit quotient. Boundary tangencies
and quantitative strip attachment remain separate questions. The embedded
tree supplies no length bound for the historical question.

\section{Numerical path searches}\label{sec:finite}

A search was run over random monic polynomials with roots in the unit disc.  For
each sample the region $\{|f|<1\}$ was rasterised and shortest grid paths were
computed between every pair of roots.  The reported grid distances were
\[
  \begin{array}{lrr}
    \text{degree } 5, & 500 \text{ trials}: & 1.1052648928\\
    \text{degree } 6, & 1500 \text{ trials}: & 0.8450414343\\
    \text{degree } 8, & 1500 \text{ trials}: & 0.6203916714\\
    \text{degree } 10, & 1500 \text{ trials}: & 0.4303640486 .
  \end{array}
\]
No counterexample candidate was found, and the measured values sit well below
the threshold $2$.  These are numerical candidate connectors.  A polygon whose
vertices satisfy $|f|<1$ need not lie in the open lemniscate, since an edge can
cross the boundary between two safe vertices, and root snapping adds further
edges of the same kind.  Turning a candidate into a proof needs a continuous
certificate on every edge, for instance strict positivity of the real
polynomial $1-|f(z(t))|^2$ on $[0,1]$ for each straight edge $z(t)=u+t(v-u)$,
established by exact coefficients or outward-rounded interval arithmetic with
subdivision, together with exact root enclosures.  No such certificate is
attached to the numbers above.  They are grid distances for the sampled
configurations, and the apparent decrease with degree is a property of the
sample, from which we draw no conjecture.  Even a finite collection of proved instances would not establish
the universal statement. These searches do not test \texttt{ani}'s degree-seven counterexample.
Useful stress tests for particular path constructions include near-degenerate
saddles, thin necks, boundary-critical configurations, and almost-connected
separatrices.  They should report those diagnostics.  Raster paths remain candidate
finders, never continuous certificates.

\section{Further questions about topology, length and perturbation}\label{sec:open}

The questions below separate topological, metric and perturbative steps
of earlier arguments. They are retained for their own content, not as a
claim that completing them must prove the unrestricted historical assertion.
A counterexample or a necessary additional hypothesis can be as informative
as a positive theorem.

\subsection*{1. Local saddle geometry and the refuted tree estimate}

\begin{problem}[corrected local saddle assembly]\label{prob:saddle1041}
The printed $1/(2\pi)$ spanning-tree estimate of Proposition~12 is false: the
Cassini polynomial $z^2-a^2$ at $a=9/10$ makes the proposed tree budget
strictly shorter than the distance between its roots.  Any later argument must
pay a positive attachment cost, select only one short pair instead of spanning
every root, or use a different global metric inequality.  A four-pronged or
cut-annulus local model may still be useful for another estimate, but it cannot
recover~\eqref{eq:prop12-bound}.
\end{problem}

The full-disc model $x^2-y^2$, an annulus in which lower branches rejoin, two
saddles joined by a separatrix and simultaneous saddle levels are mandatory
tests.  Repeating the one-lower/two-upper assertion does not answer the
problem.

\subsection*{2. Cutting a compact region into flow strips}

Let $p$ be a polynomial with simple roots and simple critical points, and
set $u=-\log|p|$ away from the roots. Simplicity of the roots makes every
critical value nonzero; critical points themselves may lie at $0$.
For regular values $c<T$, let $V$ be a component of $\{u>c\}$ and assume that
$\overline{V\cap\{c<u<T\}}$ is a compact genus-zero surface, its lower
boundary is one smooth Jordan curve, its upper boundary consists of $m$ smooth
level curves each enclosing one root, all interior critical points are nondegenerate index-one saddles,
the normalised gradient field, defined away from the saddles,
\[
 X=\frac{\nabla u}{|\nabla u|^2}=-\frac{p}{p'},
\]
is transverse to the level boundaries, and no maximal $X$-trajectory has two
saddle endpoints. The equality uses the identification of a planar vector
with a complex number: $\nabla u=-\overline{p'/p}$. In particular,
$du(X)=1$ wherever $X$ is defined.

\begin{problem}[finite strip decomposition after ray cuts]\label{prob:reeb1041}
For every $\eta>0$, construct disjoint Morse neighbourhoods $N_j$ with
$\sum_j\operatorname{diam}N_j<\eta$ and a finite set of complete separatrix or
regular-flow cuts so that the closure of each remaining component in
the cut-open surface is flow-diffeomorphic to a rectangle
\[
 [a_S,b_S]\times[0,1],\qquad u=t,
\]
with connected level sections and no uncut annular component.  Give an explicit
finite bound for the number of strips in terms of the number $s$ of
saddles and the number $m$ of upper boundary curves, or exhibit the minimal missing
hypothesis or a counterexample.
\end{problem}

No particular formula such as $2s+1$ is presumed. Boundary tangencies,
simultaneous levels and branch reunion through an annulus require explicit
cuts. The closure is taken in the cut-open surface, with its separate
boundary copies, not after regluing them in the plane. The quadratic
example following Problem~\ref{prob:globalflow1041} explains why replacing
this construction by the ordinary orbit quotient would fail.

\subsection*{3. Length estimates for joining flow strips}

For a strip $S$, write
\[
 \Gamma_t^S=S\cap\{u=t\},\qquad
 P_S(t)=\mathcal H^1(\Gamma_t^S),
\]
and write
\[
 \Phi_S(t)=\int_{\Gamma_t^S}|\nabla u|\,ds
\]
for its transverse flux. On a regular sub-band, $u$ is harmonic and the
side boundaries are flow lines. The divergence theorem therefore makes
$\Phi_S(t)$ independent of $t$. Normalise the transverse measure by
\[
 d\mu_{t_0}(x)=\frac{|\nabla u(x)|}{\Phi_S(t_0)}\,ds(x).
\]
The choice of measure is determined by the length calculation. Since
$u=t$ along the normalised flow, its Euclidean speed is
$1/|\nabla u|$. Transporting the flux measure to level $t$ therefore gives
\[
 \begin{aligned}
 \int_{\Gamma_{t_0}^S}\operatorname{len}(\gamma_x)\,d\mu_{t_0}(x)
 &=\int_{a_S}^{b_S}\int_{\Gamma_t^S}
    \frac1{|\nabla u|}\frac{|\nabla u|}{\Phi_S(t_0)}\,ds\,dt\\
 &=\frac1{\Phi_S(t_0)}\int_{a_S}^{b_S}P_S(t)\,dt .
 \end{aligned}
\]
Here $0<\Phi_S(t_0)<\infty$ for a nonempty regular strip. Tonelli's
theorem applies to the nonnegative integrand. At a critical endpoint
level, exhaustion through regular sub-bands gives the same identity,
allowing the common value to be infinite. A finite right side gives a
trajectory of length at most that mean; an identity with both sides
infinite gives no length estimate.
A strip flux need not be an integer multiple of $2\pi$. For $f(z)=z$
and $u=-\log|z|$, take an annular sector of angular
width $\theta$. On a circular level arc, $|\nabla u|=1/r$ and
$ds=r\,d\phi$, so its flux is $\Phi_S=\theta$.
For $f(z)=z^n$, the same calculation gives $\Phi_S=n\theta$. A closed
regular level curve enclosing $k$ roots has flux $2\pi k$ by the argument
principle. Restricting it to a strip replaces that full winding by the
strip's share of the flux.
The averaging measure therefore depends on the width of the strip.
Summing separate strip estimates cannot replace those widths by the
winding number of a closed level curve.  They do not reopen~\eqref{eq:prop12-bound}.  Cassini already
excludes the printed coefficient $1/(2\pi)$ for any embedded tree that contains
all $m$ roots: Problem~\ref{prob:saddle1041} records that no later gluing
argument can retain that estimate.

\begin{problem}[strip gluing after the printed coefficient]\label{prob:metric1041}
Assuming Problem~\ref{prob:reeb1041}, do one of the following, for every
$\eta>0$.
\begin{enumerate}
\item Produce an embedded tree $G$ containing all $m$ roots whose length obeys
  \begin{equation*}
 \operatorname{len}(G)
 \le C\int_{2\alpha}^{\infty}P_V(t)\,dt+\eta
\tag{6.1}\label{eq:metric-fanin1041}
 \end{equation*}
for an explicit constant $C>1/(2\pi)$, uniform over a stated class of
polynomials and an explicitly restricted range of truncations $\alpha$.
The restrictions must account for attachment cost and for a valid allocation
of the strip fluxes $\Phi_S$.
\item Connect only one pair of distinct roots, rather than spanning every root.
\item Give a different global metric inequality that is not~\eqref{eq:prop12-bound}.
\end{enumerate}
The printed coefficient $1/(2\pi)$ is not an open target.  Better additive
control of saddle, annular-cut and root-cap cost cannot recover~\eqref{eq:prop12-bound}.
\end{problem}

The restriction on $\alpha$ is essential in the first alternative.
For a fixed $f(z)=z^2-a^2$, with $0<a<1$, choose
$0<c<-2\log a$ and let $V$ be the component of $\{-\log|f|>c\}$
containing both roots. At sufficiently deep levels, the Cassini
parametrisation used above gives
\[
 P_V(t)\le\frac{2\pi e^{-t}}{\sqrt{a^2-e^{-t}}}.
\]
Consequently $\int_{2\alpha}^\infty P_V(t)\,dt\to0$ as
$\alpha\to\infty$, whereas every spanning tree has length at least
$2a$. No fixed $C$ can cover arbitrary truncations. This is a
limitation of the unrestricted formulation of (6.1), separate from the
earlier counterexample to the particular coefficient $1/(2\pi)$.

The printed collar slack
\[
 q=\frac1{2\pi}\int_{\alpha}^{2\alpha}P_V(t)\,dt>0
\]
used the same excluded coefficient.  It is not an unused error allowance in an otherwise
complete spanning-tree estimate.  Independently choosing a shortest trajectory
in each strip is not enough unless the attachment mismatch is controlled, and
controlling that mismatch does not restore the printed $1/(2\pi)$ tree bound.

\subsection*{4. Coefficient perturbation and stability}

The constant-translation stage is no longer open.  Once a finite critical-value
family is injective, Lean proves an arbitrarily small translation making every
value nonzero and pairwise positive-ray separated
(\lword{Erdos1041/NewtonFlowRaySeparation.lean}{197}{exists_small_translation_separating_arguments}{a small translation separating the arguments}).  It also proves the
explicit root-retention estimate: for a monic polynomial $f$ of degree
$n\ge1$, a constant shift of modulus below $\varepsilon>0$ keeps all roots
in the unit disc provided
\[
 ((n+1)\varepsilon)^{1/n}+\max_{f(b)=0}|b|<1
\]
(\lword{Erdos1041/NewtonFlowRaySeparation.lean}{287}{constant_perturbation_roots_in_unitDisk}{the root-retention estimate for the unit disc}).  Adding a constant leaves the difference of two critical values unchanged,
so it cannot separate values that were equal to begin with.

For a constant shift alone, factorisation gives the sharper elementary
bound
\[
 (f+\beta)(z)=0\quad\Longrightarrow\quad
 \prod_{f(b)=0}|z-b|=|\beta|
 \quad\Longrightarrow\quad
 \min_{f(b)=0}|z-b|\le|\beta|^{1/n},
\]
where the product counts roots with multiplicity. Hence
$\max_{f(b)=0}|b|+|\beta|^{1/n}<1$ already ensures retention in the
open unit disc. The factor $n+1$ in the recorded Lean estimate is not
needed for this ordinary constant-shift argument. The coefficient one
in the displacement bound is optimal, as $f(z)=z^n$ shows. The estimate
does not match labelled root lists or bound the cost of transporting
an existing curve.

Qualitatively, a small linear perturbation also supplies distinct
critical values as well as simple critical points. For $\deg f\ge2$,
a multiple critical point of $f+\lambda z$ must satisfy
$f''(c)=0$ and $\lambda=-f'(c)$; only finitely many parameters are
excluded. Near any other parameter, label the distinct critical points
by holomorphic functions $c_j(\lambda)$. Their values
\[
 v_j(\lambda)=f(c_j(\lambda))+\lambda c_j(\lambda)
 \quad\hbox{satisfy}\quad v_j'(\lambda)=c_j(\lambda).
\]
Thus $v_i-v_j$ cannot vanish identically for $i\ne j$, and its zeros
are isolated. Avoiding the finitely many pairs in a sufficiently small
parameter neighbourhood gives simple critical points with distinct
values. The neighbourhood can be chosen inside any prescribed disc
about $\lambda=0$, proving the required arbitrarily small choice.

If distinct moduli are also needed for the excellent Morse function,
the translation $\beta$ can avoid them at the same time as the ray
collisions. For each distinct pair $v_i,v_j$, the equality
$|v_i+\beta|=|v_j+\beta|$ is the real affine line
\[
 2\Re\bigl((v_i-v_j)\overline\beta\bigr)=|v_j|^2-|v_i|^2.
\]
For ray separation it is enough to avoid, for each pair, the real line
through $-v_i$ and $-v_j$: the ray-collision formula puts every forbidden
translation on that line. Avoid also the finitely many points $-v_j$
to keep the translated values nonzero. Thus the full excluded set is
contained in two real affine lines per pair of critical values,
together with finitely many points. Its complement meets every open
disc, so imposing distinct moduli costs no lower bound on the size
of the perturbation.
Root retention is also qualitative: choose a circle $|z|=R<1$
enclosing the original open-disc roots. If
$|\lambda|R+|\beta|<\min_{|z|=R}|f(z)|$, Rouch\'e's theorem
keeps all roots of $f+\lambda z+\beta$ inside that circle.
These ordinary arguments establish genericity and root retention;
they do not estimate the collar or the cost of transferring a selected
curve. The first two requirements below therefore have qualitative
proofs, while their simultaneous quantitative use is the remaining task.

For the next problem, fix a specified valid replacement for the refuted
metric inequality. Its slack is the difference between the permitted
length and the bound supplied by that replacement. The symbols $q$ and
$q_g$ refer to this same specified quantity for $f$ and $g$, respectively;
they are not the historical collar integral with coefficient $1/(2\pi)$.
Without a choice of replacement inequality, the requirement $q_g>q/2$
has no defined geometric content.

\begin{problem}[two-stage generic perturbation with slack]\label{prob:perturb1041}
For fixed root discs, compact collar $K$, regular levels
$\alpha/2,\alpha,2\alpha,5\alpha/2$, and a specified valid replacement metric
inequality with slack $q>0$ at $f$, prove that an arbitrarily small
\[
 g_{\lambda,\beta}(z)=f(z)+\lambda z+\beta
\]
can be chosen so that:

\begin{enumerate}
\item $f+\lambda z$ has simple relevant critical points and injective complex
critical values;
\item the checked constant translation $\beta$ makes those values nonzero and
pairwise ray-separated;
\item no critical point enters the protected collar or truncation boundary;
\item the relevant component remains in $K$, contains exactly the corresponding
perturbed roots and has slack $q_g>q/2$; and
\item root displacement and straight-line transfer back to the original roots
consume less than a prescribed fraction of $q/m$.
\end{enumerate}

If the one-coefficient perturbation $\lambda z$ cannot ensure all five
properties, give the smallest additional lower-coefficient direction that can,
or an explicit obstruction.
\end{problem}

The cited declarations supply finite planar avoidance and the constant
translation. The ordinary argument above additionally supplies a linear
perturbation with simple critical points and distinct critical values.
What remains for the stated problem is simultaneous control of the
protected collar, the component, the quantitative slack and the
root-to-root transfer cost. No formalisation of the new genericity
argument or solution of those stability requirements is claimed.

\subsection*{5. Newton flow on a compact region between regular levels}

Classical Newton-flow theory describes maximal trajectories and the graph
they form. Away from the roots and critical points, every orbit satisfies
\[
 p(z(t))=e^{-(t-t_0)}p(z(t_0)),\qquad
 u(z(t))=u(z(t_0))+t-t_0,\qquad u=-\log|p|,
\]
by the lemma of Shub, Tischler and Williams
\cite[Lemma~2.1]{kozen-stefansson1997}. For an initial point that is neither
a root nor a critical point, the maximal Newton trajectory has one of
four endpoint patterns: each finite-time end is a critical point, an
infinite forward end is a root, and an infinite backward end is the
point at infinity \cite[Lemma~2.2]{kozen-stefansson1997}. Stationary
trajectories at simple roots are not part of this classification. The
maximal nonconstant trajectories with endpoints among the roots and critical
points form a connected Newtonian graph with finitely many edges
\cite[Definition~2.1 and \S2]{kozen-stefansson1997}. Critical-to-critical
edges are included; pairwise distinct arguments exclude them by the value
equation, not by the definition of the graph.

The smooth field needed for Morse theory is not $-p/p'$, which is
undefined at a critical point with nonzero value. On the root-free band the
Euclidean gradient satisfies
\[
 \nabla u=-\overline{p'/p}
   =\frac{|p'|^2}{|p|^2}\left(-\frac{p}{p'}\right)
 \qquad(p'\ne0).
\]
It has the same oriented nonstationary trajectories as the Newton
field, with a different time parameter. At a simple critical point
$c$, its derivative is the real Hessian of $u$, whose eigenvalues are
$\pm|p''(c)/p(c)|$. Thus it extends as a hyperbolic saddle there.
A finite critical endpoint in Newton time is not a finite arrival at
an equilibrium of the smooth field. These two parametrisations must be
distinguished in a statement about a Morse--Smale flow.

\begin{problem}[relative global Newton-flow theorem]\label{prob:globalflow1041}
Let $p$ be a nonconstant polynomial, let $a<b$ be finite regular values
of $u=-\log|p|$, and suppose every critical point in the compact band
\[
 \{z:a\le u(z)\le b\}
\]
is simple, with the critical values in the band on pairwise distinct
positive rays. Describe the maximal trajectories of the smooth gradient
field $\nabla u$ relative to the lower and upper level boundaries,
each of which may have several components. Construct a cut-open model
by regular flow strips, specifying the separate copies of cut boundaries
and the equivalence relation used for its trajectory parameter space.
Prove that this parameter space is Hausdorff, or identify the obstruction
to the proposed identifications. The ordinary orbit quotient of the
uncut band is not required to be Hausdorff. State the boundary convention
and the comparison with Newton trajectories through the positive change
of time above, without evaluating $-p/p'$ at its singular points.
\end{problem}

The endpoint classification on this band follows from the cited Newton
classification and the time change. Each nonconstant trajectory runs
from the lower boundary to the upper boundary, from the lower boundary
to a saddle, or from a saddle to the upper boundary. Indeed, Newton time
equals the change in $u$, so the part of a Newton trajectory in the
band has bounded Newton-time interval. At either end it must meet a
regular boundary or tend to a critical point. Ray separation excludes
two different critical endpoints, and strict increase of $u$ excludes
two ends at the same critical point. In smooth-gradient time, regular
boundary crossings occur in finite time, whereas convergence to a
saddle takes infinite time by uniqueness for the smooth differential
equation. The remaining trajectories are the stationary saddles.
This is an ordinary consequence of the classical classification, not
an additional checked Lean theorem.

The planar transversality condition also follows from the hypotheses.
Along a nonstationary gradient trajectory, the value of $p$ has constant
argument, so a trajectory between two different saddles would contradict
ray separation. A homoclinic trajectory is excluded
by the strict increase of $u$. The stable-manifold theorem gives local
stable and unstable curves tangent to the Hessian eigenspaces. Their
global branches for different saddles therefore do not meet; at one
saddle they meet transversely,
since their tangent lines are the two distinct Hessian eigenspaces.
The regular lower boundary is an entrance and the regular upper boundary
an exit. With this boundary convention, the gradient satisfies the
Morse--Smale transversality condition. This does not make the quotient
by entire trajectories Hausdorff.

For an explicit obstruction, take $p(z)=z^2-1/4$ and the compact band
\[
 \log2\le-\log|p(z)|\le\log8.
\]
Its only critical point is the simple saddle at $0$, whose value is
$-1/4$; both boundary levels are regular and ray separation is vacuous.
The positive real separatrix, including its upper-boundary endpoint, is
\[
 \bigl(0,1/(2\sqrt2)\bigr],\qquad
 \frac{dx}{d\tau}=\frac{2x}{1/4-x^2}.
\]
Integration gives
$\tau-\tau_0=\tfrac18\log(x/x_0)-(x^2-x_0^2)/4$, so this
trajectory tends to $0$ as smooth-gradient time $\tau\to-\infty$.
Thus this orbit is not closed in the band: its closure also contains the
stationary orbit $\{0\}$. In the quotient topology its image is a
nonclosed singleton, so the orbit quotient is not even $T_1$, hence not
Hausdorff. Reparametrising the nonstationary trajectories leaves this
quotient unchanged.

On an actual cut-open flow rectangle, by contrast, collapsing the
trajectories gives the transverse interval. Identifying boundary copies
is a further quotient and requires a separate separation argument.
The \href{https://arxiv.org/abs/1902.08841v5}{Reeb graph} uses yet another
equivalence relation: it identifies connected components of level sets,
not complete trajectories. Neither that graph nor the embedded Newtonian
graph can be substituted for the orbit quotient without specifying a
new equivalence relation.

The recorded checked algebra gives the pointwise value equation and
excludes distinct endpoint rays once the exponential endpoint relation
is assumed. The ordinary value-equation proof derives that relation for
an existing real-time trajectory, using continuity at critical endpoints;
the separate endpoint candidate has no reported compilation here.
Global solution theory and the finite graph are the classical results
cited above. The relative endpoint classification and transversality are proved
above. They do not construct the required finite cut-open space, justify
its boundary identifications or control the length needed to join the
pieces. The quadratic example settles
the uncut-orbit-space question negatively; it does not obstruct a
construction that retains the required separate boundary copies.

This record does not establish the unrestricted conclusion of
Erd\H{o}s~\#1041.  Its sufficient conditions have different hypotheses and
containment levels: low critical value, separated simple critical value,
and the individually proved polynomial families are not a cover of the
remaining class.  In particular, the scaled $(5/2)\mu^{1/n}$ bound is at
level $(25/13)\mu$, which is smaller than $2\mu$; it therefore improves
both the constant and the level of the $71/10$ construction.  The latter
is retained for its independent proof and component-sensitive estimates.
The critical-value moment does not choose a contained pair.
The ordinary generic slit-sheet topology has classical antecedents; the
earlier proposed sufficient estimate concerned the lengths of particular
inverse-ray curves. Their minimum need not equal the infimum over all
contained curves. The degree-seven counterexample above rules out the corresponding universal
estimate for these particular curves too, without invalidating conditional
estimates on smaller classes.

\paragraph{Adjacent inverse-map methods.}
Crane \cite[Lemma~2.1, Theorem~2.2 and \S\S3--4]{crane2007smale}
uses inverse-branch hyperbolic geometry, Dubinin's radial-slit input and
capacity for a derivative-normalised Smale ratio.  Dubinin's four-point
result \cite[Theorem~1 and Corollary~4]{dubinin2013fourpoint} assumes
bounded critical values and estimates distortion.  These are relevant
methods, not direct sources of the present positive moment or connector
constant, and their normalisations are not interchangeable with a root-disc
hypothesis.

\section*{Statements and declarations}

Lean checks exact formal statements, not citation choices or mathematical
priority. The formal-source index records the complete quadratic weighted
inequality, the exponent-$2/(n-1)$ critical-value mean, the translated
cubic family and the open-disc sparse-quintic path conclusion in its
checked build. The fourth-power and higher-moment refinements proved here
remain ordinary analytic results. The complete degree-three path also has a release-source
proof under monicity, degree three and the open-disc root hypothesis,
distinct from the earlier implementation outside that build; its release status is explained above. Finite inequalities
do not replace a comparison of
the full geometric statements. The analytic path arguments and the research
addendum retain the ordinary proof status stated locally. No fresh Lean
build or independent mathematical review was performed for this revision.

\section*{Acknowledgements}
I thank Wouter van Doorn for advice on mathematical exposition, in particular
on explaining restrictive hypotheses, avoiding private terminology and
introducing notation only when it helps the reader. His remarks concerned
another note; this acknowledgement does not imply that he reviewed or
endorsed the mathematics of the present paper.

\appendix
\section{Guide to the formal sources}\label{app:sources}

The formal-source index describes sources at snapshot \texttt{6b78209ab63a}.
The links below retain their earlier commits and line numbers. The index
distinguishes complete statements in its recorded checked build from
implementations present outside that build.
\begin{itemize}
\item Complete statements recorded in the supplied checked build: \href{https://github.com/wcook04/plectis-erdos/blob/3d6d938d696fed0fb71dd55115a18a73738ff223/lean/ErdosProblems/Erdos1041/PaperCriticalValueMeanR10.lean#L101}{critical-value mean with exponent $2/(n-1)$}, \href{https://github.com/wcook04/plectis-erdos/blob/3d6d938d696fed0fb71dd55115a18a73738ff223/lean/ErdosProblems/Erdos1041/PaperWeightedRefinementsR10.lean#L17}{the quadratic weighted inequality for points in a disc}, and \href{https://github.com/wcook04/plectis-erdos/blob/3d6d938d696fed0fb71dd55115a18a73738ff223/lean/ErdosProblems/Erdos1041/PaperCubicFibres.lean#L240}{translated cubic-fibre path}.
\item Complete implementations recorded outside that build: \href{https://github.com/wcook04/plectis-erdos/blob/3d6d938d696fed0fb71dd55115a18a73738ff223/lean/ErdosProblems/Erdos1041/PaperCubicCompletion.lean#L297}{degree-three path}.
The supplied release snapshot additionally contains \href{https://github.com/wcook04/plectis-erdos-lean/blob/52f29ad173b04e3bac941b3663f2b9aebe5de0bb/ErdosProblems/Erdos1041/PaperCubicMonic.lean#L32}{the cubic theorem under monicity, degree three and the open-disc hypothesis},
with an actual proof and an expanded solution statement. Its recorded
status is \texttt{release\_only}; it is not promoted to the checked-build
category here.
\item Also recorded in the supplied checked build: \href{https://github.com/wcook04/plectis-erdos/blob/3d6d938d696fed0fb71dd55115a18a73738ff223/lean/ErdosProblems/Erdos1041/PaperPrimitiveCompletionR10.lean#L235}{the sparse-quintic path}.
\item The index also distinguishes finite inequalities, declarations supplied
only in the release sources, and ordinary analytic arguments. None of these
categories is interchangeable with a complete formal path theorem.
\end{itemize}
No source-link reachability check or fresh kernel build is implied by these
snapshot links.  The original paper-wide commit macro is retained for legacy
links; newer target links name their own snapshot explicitly.

\subsection*{Formalisation coverage and remaining dependencies}
\label{sec:coverage}
\papersectiontarget{coverage}

Every theorem, lemma, proposition and corollary of this record other than
those listed below has a Lean statement of the same assertion, with the same
hypotheses, checked by the Lean kernel using only the axioms
\texttt{propext}, \texttt{Classical.choice} and \texttt{Quot.sound}, in the
development at revision \texttt{181078b6b009}.  This includes
Lemma~\ref{res:circle-slice-packing}.

Theorem~\ref{res:low-critical-thirteen-twentyfifths} is checked in degree two,
with no hypothesis on the positions of the roots, and the exact rational
arithmetic of its stopping-time comparison is checked.  The analytic argument
in higher degrees is not yet formalised.
Corollary~\ref{res:low-critical-scale-free} and
Corollary~\ref{res:scaled-low-critical-path} are verified as reductions to
that theorem, which is itself in this list; the degree-two case of the second
corollary is checked with no input at all.

Theorem~\ref{res:critical-value-separation} and
Corollary~\ref{res:critical-value-thresholds} are verified as reductions to a
single named statement, which is what the analytic construction in the proof
of Theorem~\ref{res:critical-value-separation} produces: for every normalised
separation datum, a square-root connector of a given length together with a
component area satisfying the two displayed bounds.  That construction is not
formalised.  It uses the conformal bijection supplied by the Riemann mapping
theorem, the Bergman segment inequality, the exterior Blaschke capacity gap
and P\'olya's area--capacity inequality~\cite[Theorem~6]{crane}.  The first
clause of Corollary~\ref{res:critical-value-thresholds}, inequality
\eqref{eq:disk-family-coefficient} for every $n\ge3$, every $w_0\in[0,1]$ and
every $4/3\le S\le2$, is unconditional, and so is every numerical inequality
entering the second clause.

Theorem~\ref{res:constant-factor-path},
Corollary~\ref{res:constant-factor-arity} and
Corollary~\ref{res:constant-factor-capacity} are verified as reductions to
named statements, each of which is what the level and direction averaging in
their own proofs produces.  That construction is not formalised; its
ingredients are P\'olya's area inequality~\cite[Theorem~1]{crane}, the Koebe
distortion theorem and the coarea formula.  The degenerate case $\mu=0$, the
degree-two case and every numerical constant are checked.
Theorem~\ref{res:conjecture-p-consumer} is verified as a reduction to the
two-component split at the first critical level, which its own proof
constructs and which is not formalised; its ingredients are the component
count in the proof of \cite[Proposition~2.1]{eks2010}, the boundary extension
of the inverse branch and the Jordan curve theorem.  The half-perimeter
selection lemma is checked.  Theorem~\ref{res:dual-arity-floor} is verified as
a reduction to \eqref{eq:lc-radius-and-budget}, which its own proof derives
from the standing failure hypothesis by the Blaschke factorisation of a
Riemann map and P\'olya's area inequality; that derivation is not formalised.
The hyperbolic law of cosines it uses is proved in the development.  For
Theorem~\ref{res:attachment-aware-reeb} the development checks the
ray-disjointness, level-separation and saddle-scale steps of its proof, and
the Morse, monodromy and strip statements themselves are not yet formalised.

Formal proofs of the named inputs are not included in the verified
development, so the endpoints that carry one remain conditional on it.
Theorem~\ref{res:low-critical-thirteen-twentyfifths} and
Theorem~\ref{res:attachment-aware-reeb} carry no such input; for each of them
part of the argument is formalised and the rest is not.  A theorem whose own
statement is conditional is formalised exactly as stated; the results listed
here are asserted unconditionally.

% The exact former declaration list is retained in the editorial transfer archive.
% Part: extended record
% BEGIN GENERATED CONCORDANCE (claude_lane/build_concordance.py; do not edit by hand)
\paragraph{Concordance of statements and Lean declarations.}
Each result of this record that has a kernel-checked Lean statement of the same assertion is listed
below with the declarations that jointly state it.  Each name links to its declaration at revision
\texttt{181078b6b009}.  Where the Lean statement is stronger than the printed one and implies it by
an immediate specialisation, the entry says so.  An entry marked \emph{compared} was also checked
independently: a restatement of the same declarations against Mathlib alone, together with its
proof, was verified by Comparator (\texttt{leanprover/comparator}) in a clean continuous
integration environment, in the run named by its number.  Comparator trusts the restated
statement, so the correspondence between the printed result and that statement is the one this
concordance records.

\begingroup\small\raggedright
\par\noindent\hangindent=1.5em The statement at~\ref{res:ani-degree-seven-counterexample-long}: \href{https://github.com/wcook04/plectis-erdos/blob/181078b6b009d905809cf2e007ad309386a3d1e0/lean/ErdosProblems/Erdos1041/Counterexample/CatalogueAdapter.lean\#L25}{\nolinkurl{erdos1041_ani_degree_seven}}.  \emph{Compared}, run \href{https://github.com/wcook04/plectis-erdos-lean/actions/runs/35643815458}{\texttt{35643815458}}.
\par\noindent\hangindent=1.5em Theorem~\ref{res:trinomial-all-degree} (the Lean statement is stronger): \href{https://github.com/wcook04/plectis-erdos/blob/181078b6b009d905809cf2e007ad309386a3d1e0/lean/ErdosProblems/Erdos1041/PaperTrinomial.lean\#L26}{\nolinkurl{all_spokes}}, \href{https://github.com/wcook04/plectis-erdos/blob/181078b6b009d905809cf2e007ad309386a3d1e0/lean/ErdosProblems/Erdos1041/PaperTrinomial.lean\#L38}{\nolinkurl{complete_trinomial}}.  \emph{Compared}, run \href{https://github.com/wcook04/plectis-erdos-lean/actions/runs/35643815458}{\texttt{35643815458}}.
\par\noindent\hangindent=1.5em Lemma~\ref{res:circle-slice-packing} (the Lean statement is stronger): \href{https://github.com/wcook04/plectis-erdos/blob/181078b6b009d905809cf2e007ad309386a3d1e0/lean/ErdosProblems/Erdos1041/PaperCompleteR21/HyperbolicLawOfCosines.lean\#L563}{\nolinkurl{circle_slice_packing}}, \href{https://github.com/wcook04/plectis-erdos/blob/181078b6b009d905809cf2e007ad309386a3d1e0/lean/ErdosProblems/Erdos1041/PaperCompleteR21/HyperbolicLawOfCosines.lean\#L141}{\nolinkurl{cosh_dist_polar}}, \href{https://github.com/wcook04/plectis-erdos/blob/181078b6b009d905809cf2e007ad309386a3d1e0/lean/ErdosProblems/Erdos1041/PaperCompleteR21/HyperbolicLawOfCosines.lean\#L166}{\nolinkurl{dist_polar_I}}, \href{https://github.com/wcook04/plectis-erdos/blob/181078b6b009d905809cf2e007ad309386a3d1e0/lean/ErdosProblems/Erdos1041/PaperCompleteR21/HyperbolicLawOfCosines.lean\#L177}{\nolinkurl{exists_polar}}, \href{https://github.com/wcook04/plectis-erdos/blob/181078b6b009d905809cf2e007ad309386a3d1e0/lean/ErdosProblems/Erdos1041/PaperCompleteR21/HyperbolicLawOfCosines.lean\#L162}{\nolinkurl{polar_zero_zero}}, \href{https://github.com/wcook04/plectis-erdos/blob/181078b6b009d905809cf2e007ad309386a3d1e0/lean/ErdosProblems/Erdos1041/PaperCompleteR21/HyperbolicLawOfCosines.lean\#L315}{\nolinkurl{circle_slice_packing_abstract}}.
\par\noindent\hangindent=1.5em Theorem~\ref{res:degree-three} (the Lean statement is stronger): \href{https://github.com/wcook04/plectis-erdos/blob/181078b6b009d905809cf2e007ad309386a3d1e0/lean/ErdosProblems/Erdos1041/PaperCubicCompletion.lean\#L297}{\nolinkurl{cubic_paper_complete}}.  \emph{Compared}, run \href{https://github.com/wcook04/plectis-erdos-lean/actions/runs/35544127144}{\texttt{35544127144}}.
\par\noindent\hangindent=1.5em Theorem~\ref{prop:sharp-collinear-chebyshev-comparator}: \href{https://github.com/wcook04/plectis-erdos/blob/181078b6b009d905809cf2e007ad309386a3d1e0/lean/ErdosProblems/Erdos1041/SharpCollinearChebyshev.lean\#L133}{\nolinkurl{exists_peak_le_comparisonBound}}.  \emph{Compared}, run \href{https://github.com/wcook04/plectis-erdos-lean/actions/runs/35643815458}{\texttt{35643815458}}.
\par\noindent\hangindent=1.5em Theorem~\ref{thm:sharp-collinear-diameter}: \href{https://github.com/wcook04/plectis-erdos/blob/181078b6b009d905809cf2e007ad309386a3d1e0/lean/ErdosProblems/Erdos1041/PaperCompleteR21/CollinearDiameterWhole.lean\#L367}{\nolinkurl{sharp_collinear_root_diameter}}, \href{https://github.com/wcook04/plectis-erdos/blob/181078b6b009d905809cf2e007ad309386a3d1e0/lean/ErdosProblems/Erdos1041/PaperCompleteR21/CollinearDiameterWhole.lean\#L967}{\nolinkurl{sharp_collinear_root_diameter_monic}}, \href{https://github.com/wcook04/plectis-erdos/blob/181078b6b009d905809cf2e007ad309386a3d1e0/lean/ErdosProblems/Erdos1041/PaperCompleteR21/CollinearDiameterWhole.lean\#L928}{\nolinkurl{exists_collinear_factorisation}}, \href{https://github.com/wcook04/plectis-erdos/blob/181078b6b009d905809cf2e007ad309386a3d1e0/lean/ErdosProblems/Erdos1041/PaperCompleteR21/CollinearDiameterWhole.lean\#L212}{\nolinkurl{exists_gap_le_comparisonBound}}, \href{https://github.com/wcook04/plectis-erdos/blob/181078b6b009d905809cf2e007ad309386a3d1e0/lean/ErdosProblems/Erdos1041/PaperCompleteR21/CollinearDiameterWhole.lean\#L836}{\nolinkurl{sharp_collinear_equality_attained}}, \href{https://github.com/wcook04/plectis-erdos/blob/181078b6b009d905809cf2e007ad309386a3d1e0/lean/ErdosProblems/Erdos1041/PaperCompleteR21/CollinearDiameterWhole.lean\#L769}{\nolinkurl{chebyshev_configuration_attains}}, \href{https://github.com/wcook04/plectis-erdos/blob/181078b6b009d905809cf2e007ad309386a3d1e0/lean/ErdosProblems/Erdos1041/PaperCompleteR21/CollinearDiameterWhole.lean\#L684}{\nolinkurl{monicScaledChebyshev_eq_prod}}, \href{https://github.com/wcook04/plectis-erdos/blob/181078b6b009d905809cf2e007ad309386a3d1e0/lean/ErdosProblems/Erdos1041/PaperCompleteR21/CollinearDiameterWhole.lean\#L871}{\nolinkurl{collinearDiameterBound_sharpConstant}}, \href{https://github.com/wcook04/plectis-erdos/blob/181078b6b009d905809cf2e007ad309386a3d1e0/lean/ErdosProblems/Erdos1041/PaperCompleteR21/CollinearDiameterWhole.lean\#L878}{\nolinkurl{sharpConstant_le_of_collinearDiameterBound}}.  \emph{Compared}, run \href{https://github.com/wcook04/plectis-erdos-lean/actions/runs/35643815458}{\texttt{35643815458}}.
\par\noindent\hangindent=1.5em Corollary~\ref{cor:collinear-erdos-1041}: \href{https://github.com/wcook04/plectis-erdos/blob/181078b6b009d905809cf2e007ad309386a3d1e0/lean/ErdosProblems/Erdos1041/PaperCompleteR21/CollinearDiameterWhole.lean\#L515}{\nolinkurl{collinear_erdos_1041}}, \href{https://github.com/wcook04/plectis-erdos/blob/181078b6b009d905809cf2e007ad309386a3d1e0/lean/ErdosProblems/Erdos1041/PaperCompleteR21/CollinearDiameterWhole.lean\#L984}{\nolinkurl{collinear_erdos_1041_monic}}, \href{https://github.com/wcook04/plectis-erdos/blob/181078b6b009d905809cf2e007ad309386a3d1e0/lean/ErdosProblems/Erdos1041/PaperCompleteR21/CollinearDiameterWhole.lean\#L928}{\nolinkurl{exists_collinear_factorisation}}, \href{https://github.com/wcook04/plectis-erdos/blob/181078b6b009d905809cf2e007ad309386a3d1e0/lean/ErdosProblems/Erdos1041/PaperCompleteR21/CollinearDiameterWhole.lean\#L367}{\nolinkurl{sharp_collinear_root_diameter}}.  \emph{Compared}, run \href{https://github.com/wcook04/plectis-erdos-lean/actions/runs/35643815458}{\texttt{35643815458}}.
\par\noindent\hangindent=1.5em Theorem~\ref{prop:primitive-quintic-two-tail-energy-selector}: \href{https://github.com/wcook04/plectis-erdos/blob/181078b6b009d905809cf2e007ad309386a3d1e0/lean/ErdosProblems/Erdos1041/PrimitiveQuinticInteriorTail.lean\#L272}{\nolinkurl{primitiveInterior_exists_two_tailEnergy_lt_one}}.  \emph{Compared}, run \href{https://github.com/wcook04/plectis-erdos-lean/actions/runs/35643815458}{\texttt{35643815458}}.
\par\noindent\hangindent=1.5em Theorem~\ref{thm:primitive-quintic-two-tail}: \href{https://github.com/wcook04/plectis-erdos/blob/181078b6b009d905809cf2e007ad309386a3d1e0/lean/ErdosProblems/Erdos1041/PaperCompleteR21/PrimitiveQuinticClosedDisc.lean\#L260}{\nolinkurl{primitive_quintic_two_tail}}, \href{https://github.com/wcook04/plectis-erdos/blob/181078b6b009d905809cf2e007ad309386a3d1e0/lean/ErdosProblems/Erdos1041/PaperCompleteR21/PrimitiveQuinticClosedDisc.lean\#L313}{\nolinkurl{primitive_quintic_two_tail_of_polynomial}}, \href{https://github.com/wcook04/plectis-erdos/blob/181078b6b009d905809cf2e007ad309386a3d1e0/lean/ErdosProblems/Erdos1041/PaperCompleteR21/PrimitiveQuinticClosedDisc.lean\#L93}{\nolinkurl{two_tails_closedDisc_of_ne_zero}}, \href{https://github.com/wcook04/plectis-erdos/blob/181078b6b009d905809cf2e007ad309386a3d1e0/lean/ErdosProblems/Erdos1041/PaperCompleteR21/PrimitiveQuinticClosedDisc.lean\#L199}{\nolinkurl{tail_le_one_and_eq_iff_of_leading_zero}}, \href{https://github.com/wcook04/plectis-erdos/blob/181078b6b009d905809cf2e007ad309386a3d1e0/lean/ErdosProblems/Erdos1041/PaperCompleteR21/PrimitiveQuinticClosedDisc.lean\#L82}{\nolinkurl{tail_norm_of_leading_zero}}.  \emph{Compared}, run \href{https://github.com/wcook04/plectis-erdos-lean/actions/runs/35643815458}{\texttt{35643815458}}.
\par\noindent\hangindent=1.5em Theorem~\ref{lem:cubic-safe-root-spoke}: \href{https://github.com/wcook04/plectis-erdos/blob/181078b6b009d905809cf2e007ad309386a3d1e0/lean/ErdosProblems/Erdos1041/CubicQuotientFiberCase.lean\#L161}{\nolinkurl{cubic_has_safe_root_spoke}}.  \emph{Compared}, run \href{https://github.com/wcook04/plectis-erdos-lean/actions/runs/35643815458}{\texttt{35643815458}}.
\par\noindent\hangindent=1.5em Theorem~\ref{thm:translated-cubic-quotient-fibres}: \href{https://github.com/wcook04/plectis-erdos/blob/181078b6b009d905809cf2e007ad309386a3d1e0/lean/ErdosProblems/Erdos1041/PaperCubicFibres.lean\#L240}{\nolinkurl{complete_translated_cubic_quotient_fibres}}.  \emph{Compared}, run \href{https://github.com/wcook04/plectis-erdos-lean/actions/runs/35544127144}{\texttt{35544127144}}.
\par\noindent\hangindent=1.5em Theorem~\ref{res:critical-value-budget} (the Lean statement is stronger): \href{https://github.com/wcook04/plectis-erdos/blob/181078b6b009d905809cf2e007ad309386a3d1e0/lean/ErdosProblems/Erdos1041/PaperCompleteR20/CriticalMeanWhole.lean\#L14}{\nolinkurl{critical_value_three_budgets}}, \href{https://github.com/wcook04/plectis-erdos/blob/181078b6b009d905809cf2e007ad309386a3d1e0/lean/ErdosProblems/Erdos1041/PaperCompleteR20/CriticalMeanWhole.lean\#L36}{\nolinkurl{critical_value_three_budgets_sharp}}.  \emph{Compared}, run \href{https://github.com/wcook04/plectis-erdos-lean/actions/runs/35643815458}{\texttt{35643815458}}.
\par\noindent\hangindent=1.5em Lemma~\ref{res:reflected-critical-value}: \href{https://github.com/wcook04/plectis-erdos/blob/181078b6b009d905809cf2e007ad309386a3d1e0/lean/ErdosProblems/Erdos1041/PaperReflectedCompletion.lean\#L254}{\nolinkurl{reflected_critical_value}}.  \emph{Compared}, run \href{https://github.com/wcook04/plectis-erdos-lean/actions/runs/35643815458}{\texttt{35643815458}}.
\par\noindent\hangindent=1.5em Theorem~\ref{res:value} (the Lean statement is stronger): \href{https://github.com/wcook04/plectis-erdos/blob/181078b6b009d905809cf2e007ad309386a3d1e0/lean/ErdosProblems/Erdos1041/PaperCompleteR20/NewtonRealTime.lean\#L52}{\nolinkurl{newton_real_value_whole}}.  \emph{Compared}, run \href{https://github.com/wcook04/plectis-erdos-lean/actions/runs/35643815458}{\texttt{35643815458}}.
\par\noindent\hangindent=1.5em Corollary~\ref{res:ray} (the Lean statement is stronger): \href{https://github.com/wcook04/plectis-erdos/blob/181078b6b009d905809cf2e007ad309386a3d1e0/lean/ErdosProblems/Erdos1041/PaperCompleteR20/NewtonRealTime.lean\#L69}{\nolinkurl{newton_real_endpoint_whole}}.  \emph{Compared}, run \href{https://github.com/wcook04/plectis-erdos-lean/actions/runs/35643815458}{\texttt{35643815458}}.
\par\noindent\hangindent=1.5em Theorem~\ref{res:locus}: \href{https://github.com/wcook04/plectis-erdos/blob/181078b6b009d905809cf2e007ad309386a3d1e0/lean/ErdosProblems/Erdos1041/NewtonFlowRaySeparation.lean\#L107}{\nolinkurl{translated_samePositiveRay_parameterization}}.  \emph{Compared}, run \href{https://github.com/wcook04/plectis-erdos-lean/actions/runs/35643815458}{\texttt{35643815458}}.
\par\endgroup
% END GENERATED CONCORDANCE

\section{Numerical estimates for the packing inequality}
\label{sec:record-circle-slice}

This appendix records numerical estimates for the packing argument in
Theorem~\ref{res:low-critical-thirteen-twentyfifths}. The short paper does not depend on the floating-point comparisons below.
Those comparisons are not certified upper bounds; the separate table of
certified lower root counts is identified as such. The computations indicate
where particular relaxations lose strength in the displayed range
$13/25<\mu<1$. This is not an uncovered interval: the rigorous
comparison above also covers $\mu\le529/1000$. The floating-point
comparisons neither enlarge that certified range nor assert the
historical path bound throughout the rest of the interval.

At $a=1$ the constants of the failure inequalities have the rounded
values $\delta\approx0.4586751$, $D\approx5.3770730$, and
$d_{\mathrm{low}}\approx2.1700770$; the identity $\cosh(D/2)=e^{2}$ is exact. The quantity being estimated is
$\sum_j\lambda(d_j)$ for $k$ roots at this fixed area, not a Euclidean
path length. The four columns compare the earlier bound
\[
 \max\!\left\{\frac\delta2+(k-1)\lambda(D-d_{\mathrm{low}}),
                  k\operatorname{artanh}(e^{-2})\right\},
\]
where $\lambda(d)=-\log\tanh(d/2)$ as in
Section~\ref{sec:low-critical-closure}; the hyperbolic packing bound used
for the earlier threshold $2/5$; the linear-programming value
$M_{\mathrm{circ}}$ of the relaxation of
Lemma~\ref{res:circle-slice-packing}; and the largest reported objective
value at a numerically constructed configuration. The last column is a
search result, not a certified feasible value or a proved optimum.

\begin{center}
\begin{tabular}{@{}rrrrr@{}}
\toprule
$k$ & earlier bound & packing & $M_{\mathrm{circ}}$ & search \\
\midrule
2 & 0.3103 & 0.5899 & 0.3102 & 0.3103 \\
3 & 0.4085 & 0.6878 & 0.3893 & 0.3784 \\
4 & 0.5447 & 0.7475 & 0.4450 & 0.4306 \\
5 & 0.6809 & 0.7906 & 0.4846 & 0.4694 \\
6 & 0.8170 & 0.8244 & 0.5155 & 0.4986 \\
7 & 0.9532 & 0.8521 & 0.5408 & 0.5208 \\
8 & 1.0894 & 0.8757 & 0.5622 & 0.5413 \\
10 & 1.3617 & 0.9142 & 0.5972 & 0.5759 \\
14 & 1.9064 & 0.9710 & 0.6487 & 0.6228 \\
20 & 2.7234 & 1.0298 & 0.7022 & 0.6717 \\
30 & 4.0851 & 1.0955 & 0.7624 & 0.7283 \\
100 & 13.617 & 1.2871 & 0.9428 & n/a \\
\bottomrule
\end{tabular}
\end{center}

Both right-hand columns come from floating-point searches. In addition,
$M_{\mathrm{circ}}$ uses a $2600$-point $d$-grid and a $1400$-point
$r$-grid. Restricting the allowed $d$ values and checking only finitely
many circle constraints have opposite effects on the feasible set.
Without a separate continuous feasibility or dual certificate, the
computed value is not a proved lower or upper bound for the continuous
relaxation. At $k=2$, the grid value $0.31018$ differs from the analytic
value, approximately $0.310338$. Their proximity supplies no rigorous
bracket.

The certified lower bounds for $k$, rounded up to integers, are as
follows at $a=1$.

\begin{center}
\begin{tabular}{@{}lrrrrrrrr@{}}
\toprule
$x$ & 0.40 & 0.45 & 0.50 & 0.55 & 0.60 & 0.63 & 0.65 & 0.70 \\
\midrule
earlier bound & 3 & 4 & 4 & 5 & 5 & 5 & 5 & 6 \\
packing bound & 2 & 2 & 2 & 2 & 2 & 3 & 3 & 4 \\
circle-slice bound & 3 & 4 & 6 & 8 & 10 & 12 & 14 & 18 \\
\bottomrule
\end{tabular}
\end{center}

In one explicit-Euler floating replica with step $10^{-3}$ and a geometric grid
of $40$ initial areas from $10^{-6}$ to $1$, the largest computed time to reach the area cap
falls from $0.89703$ with the earlier packing lower bound for $k$ to
$0.65503$ with the certified circle-slice bounds, and to $0.63003$ if
the relaxation's linear-programming value replaces those bounds.  Shrinking the relaxation by the worst
measured slack factor $1.045$ moves it only to $0.60703$.  These computations
suggest that this particular relaxation may stop near
$e^{-0.607}=0.545$; they do not prove that it does, and they exclude neither a stronger lower bound on the root count nor a larger certified
range.

The search uses only ten or fourteen radii in its two settings and a
coarse $a$-grid with upward rounding. The reported loss from the first
restriction at $a=1$, $k=8$ is about one per cent; the reported rounding
loss is comparable. These are empirical estimates, not certified error
bounds relative to the unrestricted continuous optimum.

The continuous optimisation remains unevaluated: maximise
$\sum_j\lambda(d_j)$ subject to $d_j\ge d_{\mathrm{low}}$ and
$\sum_jw(d_j,r)\le\pi$ for every $r>0$, allowing the relaxed radial
distributions used in the computation. The numerical optimiser spreads
mass over several radii and suggests logarithmic growth in $k$.
Neither that growth law nor its coefficient is proved. Evaluating the
relaxation would clarify this packing argument; its effect on the final
path threshold would still need to be established.

% Part: family_catalogue
\section{Guide to the results}\label{sec:erdos-1041-complete-family-map}
The table distinguishes the main conclusions and their limitations.
Each theorem retains its own hypotheses; the scope of formal verification
is recorded separately.
\begin{center}\small
\begin{tabular}{p{.24\linewidth}p{.65\linewidth}}
\toprule
Family or mechanism & Conclusion and limitations \\
\midrule
Trinomials & The root equation controls every root-to-origin segment. \\
Low critical values & Area growth and a finite certificate give the stated
cutoff $13/25$; the recorded certificate also covers $\mu\le529/1000$.
The scaled consequence uses level $(25/13)\mu$. \\
Separated simple value & A square root removes local branching;
Bergman and capacity estimates bound length. The chosen critical point
must be simple and its value isolated. \\
Collinear and sparse families & Segment constructions for roots on a line
or specified coefficients, not for arbitrary sparse polynomials. \\
Critical-value means & An unconditional exponent $4/(n-1)$ with sharp
constant; higher exponents require specified complex power-sum cancellations.
No choice of a short contained path follows. \\
Counterexamples and conditional reductions & The Cassini example refutes
the stated spanning-tree coefficient. The degree-eight example refutes
the proposed Gamma-function perimeter constant, not every uniform
perimeter bound. An assumed uniform subcritical perimeter bound
$\beta\sigma^{1/n}$ gives length at most $\beta\rho$ in $K_\mu$.
Prescribed straight-segment rules also fail; the inverse-ray bound
remains a historical sufficient condition. \\
\bottomrule
\end{tabular}\end{center}
A critical-value mean, a contained segment and a formally checked finite
inequality do not supply the same conclusion.

% Part: back
\begin{thebibliography}{99}
\bibitem{polya1928} G.~P\'olya,
  \emph{Beitrag zur Verallgemeinerung des Verzerrungssatzes auf mehrfach
  zusammenh\"angende Gebiete}, Sitzungsberichte der Preussischen Akademie der
  Wissenschaften, Physikalisch-Mathematische Klasse (1928), printed
  pp.~228--232 and 280--282.
  \url{https://archive.org/details/sitzungsbericht1928preu}.
\bibitem{crane} E.~Crane,
  \emph{The areas of polynomial images and pre-images}, Bull. London Math.
  Soc. \textbf{36} (2004), no.~6, 786--792,
  doi:\href{https://doi.org/10.1112/S0024609304003509}{10.1112/S0024609304003509};
  preprint arXiv:\href{https://arxiv.org/abs/math/0302189v1}{math/0302189v1},
  whose statement numbers are cited.
\bibitem{bloom} T.~F. Bloom, \emph{Erd\H{o}s Problems}, problem~1041.
  \url{https://www.erdosproblems.com/1041}
\bibitem{ehp1958} P.~Erd\H{o}s, F.~Herzog, and G.~Piranian,
  \emph{Metric properties of polynomials}, J. Analyse Math. \textbf{6}
  (1958), 125--148,
  doi:\href{https://doi.org/10.1007/BF02790232}{10.1007/BF02790232}.
\bibitem{ghosh2023} S.~Ghosh and K.~Ramachandran,
  \emph{Number of Components of Polynomial Lemniscates: A Problem of Erd\H{o}s,
  Herzog, and Piranian}, J. Math. Anal. Appl. \textbf{540} (2024), no.~1,
  128571,
  doi:\href{https://doi.org/10.1016/j.jmaa.2024.128571}{10.1016/j.jmaa.2024.128571};
  preprint arXiv:\href{https://arxiv.org/abs/2312.13673v1}{2312.13673v1},
  whose statement numbers are cited.
\bibitem{sutherland1992} S.~Sutherland,
  \emph{Bad Polynomials for Newton's Method}, in B.~Bielefeld and M.~Lyubich
  (eds.), \emph{Conformal Dynamics Problem List}, Stony Brook IMS preprint
  (1992), 42--44.
  \url{https://www.math.stonybrook.edu/preprints/ims92-7.pdf}
\bibitem{kozen-stefansson1997} D.~Kozen and K.~Stef\'ansson,
  \emph{Computing the Newtonian graph}, J. Symbolic Comput. \textbf{24}
  (1997), no.~2, 125--136,
  doi:\href{https://doi.org/10.1006/jsco.1997.0118}{10.1006/jsco.1997.0118};
  authors' copy \url{https://www.cs.cornell.edu/kozen/Papers/newton.pdf}.
\bibitem{eks2010} P.~Ebenfelt, D.~Khavinson, and H.~S. Shapiro,
  \emph{Two-dimensional shapes and lemniscates}, in \emph{Complex Analysis and
  Dynamical Systems IV, Part~1}, Contemp. Math. \textbf{553}, Amer. Math. Soc.,
  Providence, RI, 2011, 45--59,
  doi:\href{https://doi.org/10.1090/conm/553/10931}{10.1090/conm/553/10931};
  preprint arXiv:\href{https://arxiv.org/abs/1003.4567v1}{1003.4567v1},
  whose statement numbers are cited.
\bibitem{dubinin2006capacity} V.~N. Dubinin,
  \emph{Lemniscates and inequalities for the logarithmic capacities of
  continua}, Mat. Zametki \textbf{80} (2006), no.~1, 33--37,
  doi:\href{https://doi.org/10.4213/mzm2777}{10.4213/mzm2777}; English
  translation, Math. Notes \textbf{80} (2006), no.~1, 31--35,
  doi:\href{https://doi.org/10.1007/s11006-006-0105-8}{10.1007/s11006-006-0105-8}.
\bibitem{march2026} \texttt{shtuka},
  \emph{A Short Path Joining Two Zeros Inside a Polynomial Lemniscate},
  manuscript posted 24 March 2026, 48~pp.
  \url{https://shtuka123.github.io/1041/main.pdf}.
  The file at this URL has since been replaced by a shorter partial version
  that no longer contains Proposition~12; the durable public record of the
  March version, its defect, and the author's 26 March 2026 concession is the
  discussion thread at
  \url{https://www.erdosproblems.com/forum/thread/1041}.
\bibitem{june2026} V.~S. Pendyala,
  \emph{A Degree-Four Lemniscate Path Theorem},
  arXiv:\href{https://arxiv.org/abs/2606.24875v1}{2606.24875v1} (2026),
  doi:\href{https://doi.org/10.48550/arXiv.2606.24875}
  {10.48550/arXiv.2606.24875}.
\bibitem{dubinin} V.~N. Dubinin, \emph{Some inequalities for polynomials and
  rational functions associated with lemniscates}, Zap. Nauchn. Sem. POMI
  \textbf{404} (2012), 83--99; English translation, J. Math. Sci. \textbf{193}
  (2013), no.~1, 45--54,
  doi:\href{https://doi.org/10.1007/s10958-013-1432-4}{10.1007/s10958-013-1432-4}.
\bibitem{borwein1995} P.~Borwein,
  \emph{The arc length of the lemniscate $\{|p(z)|=1\}$},
  Proc. Amer. Math. Soc. \textbf{123} (1995), no.~3, 797--799,
  doi:\href{https://doi.org/10.1090/S0002-9939-1995-1223265-3}
  {10.1090/S0002-9939-1995-1223265-3}.
\bibitem{eremenkohayman1999} A.~Eremenko and W.~Hayman,
  \emph{On the length of lemniscates}, Michigan Math. J. \textbf{46}
  (1999), no.~2, 409--415,
  doi:\href{https://doi.org/10.1307/mmj/1030132418}{10.1307/mmj/1030132418};
  preprint arXiv:\href{https://arxiv.org/abs/0805.2295}{0805.2295}.
\bibitem{fryntovnazarov2008} A.~Fryntov and F.~Nazarov,
  \emph{New estimates for the length of the Erd\H{o}s--Herzog--Piranian lemniscate},
  in \emph{Linear and Complex Analysis}, Amer. Math. Soc. Transl. Ser.~2
  \textbf{226}, Amer. Math. Soc., Providence, RI, 2009, 49--60,
  doi:\href{https://doi.org/10.1090/trans2/226/05}{10.1090/trans2/226/05};
  preprint arXiv:\href{https://arxiv.org/abs/0808.0717v1}{0808.0717v1}.
\bibitem{tao2025} T.~Tao,
  \emph{The maximal length of the Erd\H{o}s--Herzog--Piranian lemniscate in
  high degree}, arXiv:\href{https://arxiv.org/abs/2512.12455v1}{2512.12455v1}
  (2025).
\bibitem{eremenkoyuditskii2012} A.~Eremenko and P.~Yuditskii,
  \emph{Comb functions}, Contemp. Math. \textbf{578} (2012), 99--118,
  doi:\href{https://doi.org/10.1090/conm/578/11472}{10.1090/conm/578/11472};
  preprint arXiv:\href{https://arxiv.org/abs/1109.1464v1}{1109.1464v1}.
\bibitem{eremenkolempert1994} A.~Eremenko and L.~Lempert,
  \emph{An extremal problem for polynomials}, Proc. Amer. Math. Soc.
  \textbf{122} (1994), no.~1, 191--193,
  doi:\href{https://doi.org/10.1090/S0002-9939-1994-1207536-1}
  {10.1090/S0002-9939-1994-1207536-1}.
\bibitem{eremenko2007} A.~Eremenko,
  \emph{A Markov-type inequality for arbitrary plane continua},
  Proc. Amer. Math. Soc. \textbf{135} (2007), no.~5, 1505--1510,
  doi:\href{https://doi.org/10.1090/S0002-9939-06-08640-0}
  {10.1090/S0002-9939-06-08640-0};
  preprint arXiv:\href{https://arxiv.org/abs/math/0606745v1}{math/0606745v1}.
\bibitem{bishoperemenkolazebnik2025} C.~J. Bishop, A.~Eremenko, and
  K.~Lazebnik, \emph{On the shapes of rational lemniscates}, Geom. Funct.
  Anal. \textbf{35} (2025), no.~2, 359--407,
  doi:\href{https://doi.org/10.1007/s00039-025-00704-2}
  {10.1007/s00039-025-00704-2};
  preprint arXiv:\href{https://arxiv.org/abs/2407.14610v1}{2407.14610v1}.
\bibitem{dubinin2006critical} V.~N. Dubinin,
  \emph{Inequalities for critical values of polynomials}, Sb. Math.
  \textbf{197} (2006), no.~8, 1167--1176,
  doi:\href{https://doi.org/10.1070/SM2006v197n08ABEH003793}{10.1070/SM2006v197n08ABEH003793}.
\bibitem{crane2007smale} E.~Crane,
  \emph{A bound for Smale's mean value conjecture for complex polynomials},
  Bull. London Math. Soc. \textbf{39} (2007), no.~5, 781--791,
  doi:\href{https://doi.org/10.1112/blms/bdm063}{10.1112/blms/bdm063};
  author preprint \url{https://people.maths.bris.ac.uk/~maetc/SMVCbound.pdf}.
\bibitem{dubinin2013fourpoint} V.~N. Dubinin,
  \emph{Four-point distortion theorem for complex polynomials},
  arXiv:\href{https://arxiv.org/abs/1301.3985v1}{1301.3985v1} (2013).
\bibitem{kuznetsova2003} O.~S. Kuznetsova and V.~G. Tkachev,
  \emph{Length functions of lemniscates}, Manuscripta Math.
  \textbf{112} (2003), 519--538,
  doi:\href{https://doi.org/10.1007/s00229-003-0411-3}{10.1007/s00229-003-0411-3};
  preprint arXiv:\href{https://arxiv.org/abs/math/0306327}{math/0306327}.
\bibitem{tischler1989} D.~Tischler,
  \emph{Critical points and values of complex polynomials},
  J. Complexity \textbf{5} (1989), no.~4, 438--456.
\bibitem{revision2026} W.~Cook / Plectis,
  \emph{Three refinements for the lemniscate-path programme}, 16 September 2026,
  accompanying unreviewed research note, Sections~1--3, with the earlier
  \emph{Structural obstructions} note preserved in the same revision bundle.
\bibitem{schur1918} I.~Schur,
  \emph{\"Uber die Verteilung der Wurzeln bei gewissen algebraischen
  Gleichungen mit ganzzahligen Koeffizienten}, Math. Z. \textbf{1}
  (1918), no.~4, 377--402.
\end{thebibliography}
